
\chapter{Domain Analysis}
\section{Utterance-Level Shoutcast Coding}
\label{sec:appendix-coding}
% \begin{landscape}
\setlength{\tabcolsep}{3pt}     % tighter horizontal padding
\renewcommand{\arraystretch}{1.0} % tighter vertical spacing

{\scriptsize

  \setcounter{AppRow}{0}
  % Real row number column with vertical rule: [#] | Type Topic Info Utter.Type Spk Utterance
  \begin{longtable}{@{}>{\stepcounter{AppRow}\theAppRow}r|L{2.3cm}L{2.2cm}L{1.4cm}L{2.2cm}L{0.5cm}U{5.6cm}@{}}
    \caption[Co-cast coding Villese vs DauT G2]{Coding of T90 and Nili co-casting Villese vs DauT -- King of the Desert 3 quarterfinals Game 2}\label{tab:t90-nili-coding}\\
    \multicolumn{7}{@{}l}{\footnotesize Sources:}\\[\SourceLinkGap]
    \multicolumn{7}{@{}l}{\footnotesize \textbullet T90's PoV: \url{https://www.youtube.com/watch?v=IPMITjSW0WA&t=1656s}}\\[\SourceLinkGap]
    \multicolumn{7}{@{}l}{\footnotesize \textbullet Nili's PoV: \url{https://www.youtube.com/watch?v=d3V1kHQjUsQ&t=1656s}}\\[\AfterSourceLinksVSpace]
    \multicolumn{1}{r}{\textbf{\#}} &
    \makecell[l]{\textbf{Comm.}\\\textbf{Category}} &
    \makecell[l]{\textbf{Subtype}} &
    \makecell[l]{\textbf{Info}\\\textbf{Used}} &
    \makecell[l]{\textbf{Utter.}\\\textbf{Type}} &
    \makecell[l]{\textbf{Spk}} &
    \multicolumn{1}{l}{\makecell[l]{\textbf{Utterance}}} \\
    \midrule
    \endfirsthead
    \multicolumn{7}{@{}l}{\textit{Table \thetable\ (continued)}}\\[0.3em]
    \toprule
    \multicolumn{1}{r}{\textbf{\#}} &
    \makecell[l]{\textbf{Comm.}\\\textbf{Category}} &
    \makecell[l]{\textbf{Subtype}} &
    \makecell[l]{\textbf{Info}\\\textbf{Used}} &
    \makecell[l]{\textbf{Utter.}\\\textbf{Type}} &
    \makecell[l]{\textbf{Spk}} &
    \multicolumn{1}{l}{\makecell[l]{\textbf{Utterance}}} \\
    \midrule
    \endhead

    \midrule
    \multicolumn{7}{r}{\small Continued on next page} \\
    \endfoot

    \bottomrule
    \endlastfoot
    & Match Introduction    &                      & Descriptive   & what              & T90     & Game two between Villese and DauT.                                                                                                                                                                                       \\
    & Technical Issues      &                      &             &                   & T90     & I think I need to speed up a little bit, yep.                                                                                                                                                                            \\
    & Civilization Analysis &                      & Descriptive   & what              & T90     & Uh, Villese in the blue is playing as Malians and DauT is playing as the Khmer.                                                                                                                                          \\
    & Historical Account    &                      & Background  &                   & T90     & I mean, it is the quarterfinals, Nili, so it makes sense that we're going to be seeing some repeats, but every single time I watch Villese, it's like I'm seeing a repeat of a civ matchup that Max played.              \\
    & Historical Account    &                      & Experiential  &                   & Nili    & Well, it makes sense, right?                                                                                                                                                                                             \\
    & Historical Account    &                      & Experiential &                   & Nili    & Preparation is there, and if two people are so connected and sharing every single thought they have, then you will see the same, right?                                                                                  \\
    & Historical Account    &                      & Experiential &                   & Nili    & We saw it in Hidden Cup, and that's why you also always thought, "Okay, is this Villese/TheMax or is this DauT/TaToH?" because those two are always doing the same.                                                      \\
    & Historical Account    &                      & Experiential &                   & T90     & Sure, yeah, and it's great because we can compare these strategies to previous games so viewers can think, "Okay, this worked and this didn't." What will this player do with this situation?                            \\
    & Event Description     & Laming               & Domain      & what              & T90     & We actually see a bit of laming here from Villese.                                                                                                                                                                       \\
    & Off-Topic             &                      &             &                   & T90     & Uh, normally I go out with this in game one, Nili, so I'll just say it real quick again.                                                                                                                                 \\
    & Off-Topic             &                      &             &                   & T90     & Thanks to Memb and Chrazini for all the work they put in with King of the Desert 3.                                                                                                                                      \\
    & Off-Topic             &                      &             &                   & T90     & This is one of four quarterfinals today and, for everyone on YouTube, of course, they'll come on different days, I'm sure.                                                                                               \\
    & Off-Topic             &                      &             &                   & T90     & And then, uh, thanks to Microsoft and Pinztec. This has been a great event so far.                                                                                                                                       \\
    & Off-Topic             &                      &             &                   & Nili    & Semifinals and finals next weekend as well, so you don't want to miss those either.                                                                                                                                      \\
    & Off-Topic             &                      &             &                   & Nili    & Lovely stuff obviously from the whole crew and, well, very enjoyable for so many viewers.                                                                                                                                \\
    & Off-Topic             &                      &             &                   & Nili    & Feedback has been really good for this tournament.                                                                                                                                                                       \\
    & Off-Topic             &                      &             &                   & T90     & Yeah, I mean the structure's been really clean.                                                                                                                                                                          \\
    & Off-Topic             &                      &             &                   & T90     & I think that's something that all tournament organizers have realized over the past couple of years: just making the structure clean for the players makes it simple for viewers, and there are the results.             \\
    & Off-Topic             &                      &             &                   & T90     & There are already thousands of people watching at the moment.                                                                                                                                                            \\
    & Civilization Analysis &                      &             &                   & T90     & Malians versus Khmer, so what's the big difference in your eyes in the civ matchup?                                                                                                                                      \\
    & Off-Topic             &                      & Experiential &                   & T90     & You actually played this against Dogao in a rated game... uh, not too long ago. I casted it.                                                                                                                             \\
    & Off-Topic             &                      & Experiential &                   & Nili    & I played against someone, I can't remember.                                                                                                                                                                              \\
    & Off-Topic             &                      & Experiential &                   & T90     & You were Khmer in that game.                                                                                                                                                                                             \\
    & Off-Topic             &                      & Experiential &                   & Nili    & Uhhh.                                                                                                                                                                                                                    \\
    & Off-Topic             &                      & Experiential &                   & T90     & Nili is like: "I remember the wins."                                                                                                                                                                                     \\
    & Off-Topic             &                      & Experiential &                   & Nili    & Then I assume I lost                                                                                                                                                                                                     \\
    & Civilization Analysis &                      & Domain      & how to            & Nili    & Yeah, it should be like "Khmer can go crazy", right?                                                                                                                                                                     \\
    & Civilization Analysis &                      & Background  & how to            & Nili    & Yesterday we saw Lierrey going for the drush.                                                                                                                                                                            \\
    & Civilization Analysis &                      & Domain      & how to            & Nili    & You can open with archers.                                                                                                                                                                                               \\
    & Civilization Analysis &                      & Domain      & how to            & Nili    & Scouts is kind of the standard play.                                                                                                                                                                                     \\
    & Civilization Analysis &                      & Domain      & how to            & Nili    & Men-at-arms is the only play that is something we'll very rarely see.                                                                                                                                                    \\
    & Event Description     & Battle               & Descriptive   & what              & Nili    & While we see some scout shenanigans here.                                                                                                                                                                                \\
    & Event Description     & Economy Delay        & Descriptive   & what              & T90     & Yeah, DauT saw the barracks and tried to snipe the weak villagers there, or the unloomed villagers.                                                                                                                      \\
    & Civilization Analysis &                      & Domain      & how to            & T90     & Yeah, normally what I talk about is the overall flexibility of Malians.                                                                                                                                                  \\
    & Civilization Analysis &                      & Descriptive   & how to            & T90     & All their wood buildings are cheaper.                                                                                                                                                                                    \\
    & Civilization Analysis &                      & Descriptive   & how to            & T90     & They have camels, they have knights, they have good infantry, they have monks, whereas Khmer, it's more so about their farms and their knights, I feel like in 1v1s.                                                     \\
    & Civilization Analysis \footnote{Kept as Civilization Analysis since we established the "Max/Villese share strategies" connection earlier} &                      & Background  &                   & T90     & Um, but I am thinking back now to how TheMax played when he played together Khmer with DauT, and I think it was the first time they met in the group stage. Max really loved drush FC with Malians.                       \\
    &                       &                      &             &                   & Nili    & Hmm hmm                                                                                                                                                                                                                  \\
    & Strategy Prediction   &                      & Analytical  & what              & T90     & This looks more like a men-at-arms build, though,                                                                                                                                                                        \\
    & Event Description     & Tech. Progress       & Domain      & what              & T90     & now that I see Loom coming in, but for a second there, because the barracks was pretty early,\footnote{This utterance describes a current event, but hidden due to an overlay bug.}                                                                                                                            \\
    & Strategy Prediction   &                      & Domain      & what could happen & T90     & I thought maybe we'd be seeing some type of a drush FC from Villese.                                                                                                                                                     \\
    & Strategy Prediction   &                      &             &                   & Nili    & It is interesting.                                                                                                                                                                                                       \\
    & Historical Account    &                      & Background  &                   & Nili    & I just saw DauT did just practice against Malians, against TaToH in our morning session.                                                                                                                                 \\
    & Historical Account    &                      &             &                   & T90     & Really? Okay.                                                                                                                                                                                                            \\
    & Historical Account    &                      &             &                   & Nili    & So, and I think I can spoil it that he actually lost to TaToH's Malians there, so it will be interesting if he can properly adjust.                                                                                      \\
    & Historical Account    &                      &             &                   & Nili    & TaToH did play quite nicely there in the end.                                                                                                                                                                            \\
    & Historical Account    &                      &             &                   & Nili    & Champs just dominating whatever DauT was putting on the map, and I'm curious how DauT will adjust.                                                                                                                       \\
    & Strategy Prediction   &                      & Analytical  & what              & Nili    & Look at that, that would be an archer opening.                                                                                                                                                                           \\
    & Strategy Prediction   &                      & Domain      & what              & Nili    & Three on gold by him, so mind games.                                                                                                                                                                                     \\
    & Strategy Prediction   &                      & Domain      & what              & T90     & An archer opening.                                                                                                                                                                                                       \\
    & Strategy Prediction   &                      & Analytical  & what              & T90     & He saw the barracks, so he expects infantry.                                                                                                                                                                             \\
    & Event Description     & Unit Production      & Domain      & what              & T90     & Villese is creating infantry                                                                                                                                                                                             \\
    & Civilization Analysis &                      &             & what              & T90     & But it's also going to be Malian men-at-arms, as they do get +1 pierce armor starting in Feudal, and then it extends to other ages.                                                                                      \\
    & Game Discussion       &                      &             & how to            & Nili    & Yeah, well, it's just in my opinion, that is one of the most overrated bonuses.                                                                                                                                          \\
    & Game Discussion       &                      &             &                   & T90     & Okay, fair, fair,                                                                                                                                                                                                        \\
    & Game Discussion       &                      &             & how to            & Nili    & Because the moment you get an archer out on this level, a man-at-arms will never get to an archer, yeah, and will never get the kill.                                                                                    \\
    & Game Discussion       &                      &             & how to            & Nili    & It just takes longer for the archer to kill.                                                                                                                                                                             \\
    & Game Discussion       &                      &             & how to            & T90     & No, I agree, and the only thing then is you have units that won't ever kill anything but can be annoying. But if players are able to use one archer, regardless of how annoying that is, yeah, it might not go anywhere. \\
    & Game Discussion       &                      & Background  & how to            & T90     & I know it's not common, but if this was a player like Vivi and Vivi scouted that, uh, villagers were on gold, I think Vivi would come forward and tower there.                                                           \\
    & Strategy Prediction   &                      & Background  & what could happen & T90     & I think he would tower right outside the gold because it's going to take forever for DauT to make the archers, and DauT would have no real defense to a forward.                                                         \\
    & Event Description     & Movement             & Descriptive   & what              & T90     & And here comes Villese. Villagers forward from Villese.                                                                                                                                                                  \\
    & Player Analysis       &                      & Background  & what              & T90     & This is not common from him.                                                                                                                                                                                             \\
    & Strategy Prediction   &                      & Analytical  & what could happen & Nili    & Oh, what's this? Pressure on the gold there. That is indeed maybe taking him by surprise.                                                                                                                                \\
    & Event Description     & Scouting Info        & Descriptive   & what              & Nili    & The scout is there. Does he see it?                                                                                                                                                                                      \\
    & Event Description     & Scouting Info        & Descriptive   & what              & Nili    & He is not reacting. Sees the wolf, and now he knows.                                                                                                                                                                     \\
    & Game Discussion       &                      & Domain      & how to            & T90     & Okay, so just to bring up why men-at-arms and archers is a bit more common as a whole in our meta nowadays, it's because you have some military while you transition.                                                    \\
    & State Update          & Military State       & Descriptive   & what              & T90     & DauT does not have military at the moment.                                                                                                                                                                               \\
    & Strategy Discussion       &                      & Domain      & what could happen & T90     & He's got to wait for the range and he's got to wait for the archers and then maybe wait for upgrades.                                                                                                                    \\
    & Event Description     & Building Production  & Descriptive   & what              & T90     & But Villese is going for a forward archery range.                                                                                                                                                                         \\
    & Event Description     & Economy Delay        & Descriptive   & what              & Nili    & Oh, and the palisade wall I saw coming was so low! Oh god, DauT maybe... some more hits would have been nice there.                                                                                                      \\
    & Event Description     & Economy Delay        & Descriptive   & what              & Nili    & Villagers fighting a bit back.                                                                                                                                                                                           \\
    & Event Description     & Economy Delay        & Descriptive   & what              & Nili    & And DauT will be stuck with that gold that he has in the bank, but 215 is actually not too bad.                                                                                                                          \\
    & Strategy Discussion       &                  & Domain      & what could happen & T90     & Yeah, that's true because you'll probably make what, like two or three archers max, and then you'll make skirmishers if you're expecting archers from Villese.                                                           \\
    & Strategy Prediction   &                      & Analytical  & what could happen & T90     & But yeah, Villese might not even build the tower on the gold now.                                                                                                                                                        \\
    & Strategy Prediction   &                      & Analytical  & what could happen & T90     & He might just go right to that woodline.                                                                                                                                                                                 \\
    & Event Description     & Scouting Info        & Analytical  & judgement         & Nili    & DauT, you can see he's already scared. Beautiful use of his scout.                                                                                                                                                       \\
    & Strategy Prediction   &                      & Analytical  & what could happen & T90     & He could snipe a skirmisher, but a lot more won't happen here.                                                                                                                                                           \\
    & Event Description     & Tech State     & Descriptive   & what              & T90     & Fletching in for DauT.                                                                                                                                                                                                   \\
    & Map Analysis          &                      & Descriptive             & what              & T90     & Berries are on the other side. Woodlines are pretty good for now, but he's on a lot of straggler trees and, yeah,                                                                                                        \\
    & Event Description     & Building Production  &             & what              & T90     & He'll make another lumber camp.\footnote{This describes a current event despite the future tense. Possibly due to the building being placed and paid for, but not "hammered into existence."}                                                                                                                                                                                         \\
    & Historical Account    &                      & Background\footnote{This is a rethorical question, thus classified as Historical}  &                   & T90     & Nili, when's the last time you saw Villese go forward in a tournament?                                                                                                                                                   \\
    & Historical Account    &                      & Background  &                   & Nili    & I don't know, I can't remember seeing a game.                                                                                                                                                                            \\
    & Historical Account    &                      & Background  &                   & Nili    & He's such a meta, standard, sitting-at-home player.                                                                                                                                                                      \\
    & Historical Account    &                      & Background  &                   & T90     & Yep                                                                                                                                                                                                                      \\
    & Historical Account    &                      & Background  & judgement         & Nili    & I love that he's mixing it up.                                                                                                                                                                                           \\
    & Historical Account    &                      & Background  & judgement         & T90     & Yeah, I like it as well, and I think it wasn't part of the plan. I think it was him adapting to the situation, understanding what DauT was going for.                                                                    \\
    & Recap                 &                      & Descriptive   & what              & T90     & And no kills yet in this game, Nili                                                                                                                                                                                      \\
    & Recap                 &                      & Descriptive   & what              & T90     & But a lot has happened.                                                                                                                                                                                                  \\
    & Recap                 &                      & Descriptive   & what              & Nili    & Yeah, actually true. A lot of mind games, tech switches, map awareness, trying to set yourself up for the long game.                                                                                                     \\
    & Event Description     & Building Production  & Descriptive   & what              & Nili    & DauT now with the defensive tower.                                                                                                                                                                                       \\
    & Event Description     & Battle               & Analytical  & what could happen & Nili    & Can he... (get it up)? Well, he did get two volleys off, so repair will be a bit more tricky there for Villese.                                                                                                          \\
    & Strategy Discussion   &                      & Domain      & judgement         & Nili    & and DauT is not doing the greatest job there with his lumber camps.                                                                                                                                                      \\
    & Strategy Discussion   &                      & Analytical  & judgement         & Nili    & I absolutely hate how he designed his woodline there without a gate, especially with Khmer. You can garrison in houses there.                                                                                            \\
    & Strategy Discussion   &                      & Analytical  & judgement         & T90     & You know what else I think is ugly, and we'll see if it even matters in this game,                                                    \\
    & Event Description   &                      & Descriptive  & what              & T90         & as both players have Fletching now and are fighting.                                                    \\
    & Strategy Discussion   &                      & Analytical  & judgement         & T90     & I hate forward archery ranges.                                                    \\
    & Strategy Discussion   &                      & Analytical  & judgement         & T90     & I understand it. I understand, like here, he wanted military, he wanted ranged units out fast, and I get that. But if it doesn't work out, Nili, he's going to have no defense at home if he's not walled.               \\
    & Strategy Discussion   &                      & Domain      & what could happen & T90     & But it might not be an issue.                                                                                                                                                                                            \\
    & Event Description     & Battle               & Descriptive   & what              & T90     & He snipes an archer,                                                                                                                                                                                                     \\
    & Event Description     & Military State       & Descriptive   & what              & T90     & and he's still got some ranged units behind,                                                                                                                                                                             \\
    & Event Description     & Civilization Context & Descriptive   & what              & T90     & and these are Malian men-at-arms paying off just a little bit more than we expected.                                                                                                                                     \\
    & Event Description     & Battle               & Descriptive   & what              & Nili    & Oh yeah, and look at that, they are finding quite some kills.                                                                                                                                                            \\
    & Event Description     & Battle               & Descriptive   & what              & Nili    & DauT now lost his whole army.                                                                                                                                                                                            \\
    & Event Description     & Military State       & Descriptive   & what              & Nili    & Only three military for him. That's only skirmishers for now.                                                                                                                                                            \\
    & Event Description     & Battle               & Descriptive   & what              & Nili    & At least he won the tower war                                                                                                                                                                                            \\
    & Strategy Prediction   &                      & Analytical  & what could happen & Nili    & but I wouldn't be surprised if Villese is trying to build another one a bit further to the right-hand side.                                                                                                              \\
    & State Update          & Eco State            & Descriptive   & what              & T90     & There's a massive eco difference in one category, and that's food.                                                                                                                                                       \\
    & State Update          & Eco State            & Descriptive   & what              & T90     & DauT has got 12 on food and his berries are about to run out,                                                                                                                                                            \\
    & Civilization Analysis &                      & Analytical  & what could happen & T90     & But you know, Khmer farms in messy games do so well for you, and that's going to be a big concern if players are leading up towards Castle Age, which is the norm.                                                       \\
    & Event Description     & Eco State            & Descriptive   & what              & T90     & Also, Villese with an idle on the house on the left. Villese with an idle in the house to the right.                                                                                                                     \\
    & Skill Discussion      &                      & Analytical  & why did           & T90     & So he's trying to do a lot right now, and even if he doesn't want the full wall at this time, he has two idles simply due to how much he's trying to juggle at the moment.                                               \\
    & Skill Discussion      &                      & Domain      &                   & Nili    & It's not easy juggling all that, obviously.                                                                                                                                                                              \\
    & Event Description     & Battle               & Descriptive   & judgement         & Nili    & Now the micro at the front... Miss-juggling there as well, some losses. DauT is on point with the skirmisher micro.                                                                                                      \\
    & Event Description     & Opinion              &             & judgement         & T90     & Yeah, very good,                                                                                                                                                                                                         \\
    & Event Description     & Building Production  & Descriptive   & what              & T90     & and also adding the stable, not really hiding the fact he's going for the stable either.                                                                                                                                 \\
    & Strategy Prediction   &                      & Descriptive   & what could happen & T90     & That's going to be something Villese will spot, as he added one at home himself.                                                                                                                                         \\
    & Strategy Discussion   &                      & Domain      & why didn't        & T90     & Villese needs to take care of that scout, but DauT is just not allowing it. That scout has done so much.                                                                                                                 \\
    & Strategy Discussion   &                      &             &                   & Nili    & Absolutely right that you mentioned it as well.                                                                                                                                                                          \\
    & Joke                  &                      &             &                   & Nili    & I think that is actually the hero unit of the game, though he sadly dies now that we start singing tales about him.                                                                                                      \\
    & Recap                 &                      & Analytical  & what              & Nili    & But he did get so much scouting! He saw the forward, saw the tower, saw exactly what was happening, and had such a good idea where the army was at all times.                                                            \\
    & Event Description     & Eco Raid             & Descriptive   & what              & T90     & All right, uh, is that finally a vill kill for Villese?                                                                                                                                                                  \\
    & Event Description     & Eco Raid             & Descriptive   & what              & T90     & DauT does not micro.                                                                                                                                                                                                     \\
    & Skill Discussion      &                      & Domain      & judgement         & T90     & That's sloppy from DauT because he probably should have hopped inside a house.                                                                                                                                           \\
    & Technical Issues      &                      & Domain      & why did           & T90     & And maybe it was sloppy because he had an issue.                                                                                                                                                                         \\
    & Technical Issues      &                      &             &                   & Nili    & That is weird.                                                                                                                                                                                                           \\
    & Technical Issues      &                      & Descriptive   &                   & T90     & It's paused.                                                                                                                                                                                                             \\
    & Event Description     & Eco Raid             & Descriptive   & why didn't        & Nili    & The house was so close there as well.                                                                                                                                                                                    \\
    & Event Description     & Eco Raid             & Descriptive   & why didn't        & Nili    & Didn't really go for the men-at-arms.                                                                                                                                                                                    \\
    & Technical Issues      &                      & Domain      & why did           & T90     & Maybe he tried to garrison and his garrison hotkey didn't work and he's like, "What? What?" I don't know.                                                                                                                \\
    & Event Description     & Eco Raid             & Analytical  & judgement         & T90     & But yeah, uh, I guess a bit sloppier DauT has been so far this game, but I still prefer his position because of that food count. This is insane.                                                                         \\
    & Event Description     & Eco State            & Descriptive   & what              & T90     & He has 300 food and climbing.                                                                                                                                                                                            \\
    & Strategy Discussion   &                      & Domain      & what could happen & T90     & He can make scouts with that. He could do so much.                                                                                                                                                                       \\
    & Strategy Discussion   &                      & Analytical  & what              & T90     & Villese finds himself with a forward archery range that is useless now and is trying to desperately wall up at home.                                                                                                      \\
    & Event Description     & Economy              & Descriptive   & what              & Nili    & Oh, DauT is now expanding towards the top there, going onto gold.                                                                                                                                                        \\
    & Event Description     & Building Production  & Descriptive   & what              & Nili    & Builds an outpost in kind of a no-man's land,                                                                                                                                                                            \\
    & Strategy Discussion   &                      & Domain      & what could happen & Nili    & just to see maybe some aggression or a counter-attack here against Villese.                                                                                                                                              \\
    & Event Description     & Economy                & Descriptive   & what              & Nili    & Over-chopping                                                                                                                                                                                                            \\
    & Event Description     & Scouting Info        & Descriptive   & what              & T90     & Yeah, however, how many archers are in here? Only two.                                                                                                                                                                   \\
    & Strategy Discussion   & Scouting Info        & Analytical  & what could happen & T90     & I guess with ranged units and scouts, you could force a reaction from Villese, but what DauT does not see... yet... let's see if he can have another hero scout!                                                         \\
    & Event Description     & Scouting Info        & Descriptive   & what              & T90     & Look, he sees the stable.                                                                                                                                                                                                \\
    & Strategy Discussion   &                      & Analytical  & what could happen & T90     & Yeah, so he'll know: "I've got to be a little bit careful here with only two archers in my group. If Villese has three or four scouts, he'll push me back."                                                              \\
    & Event Description     & Military State       & Descriptive   & what              & Nili    & DauT has two or three scouts now.                                                                                                                                                                                        \\
    & Event Description     & Battle               & Descriptive   & what              & Nili    & Goes for the house.                                                                                                                                                                                                      \\
    & Strategy Discussion   & Opinion              & Analytical  & judgement         & Nili    & That feels like a weird thing to focus on.                                                                                                                                                                               \\
    & Player Analysis       &                      & Background  &                   & T90     & You know what's so apparent to me when I watch DauT play? It's how he controls games.                                                                                                                                    \\
    & Strategy Discussion   &                      & Background  & what could happen & T90     & Like, he is not... obviously breaking in would be great, but he's just trying to force a reaction from Villese.                                                                                                          \\
    & Event Description     & Scouting Info        & Domain      & what              & T90     & He's delaying the woodline, he sees what Villese is up to.                                                                                                                                                               \\
    & Strategy Discussion   &                      & Domain      & what              & T90     & He's just waiting, and now he feels like he's in full control.                                                                                                                                                           \\
    & Strategy Discussion   &                      & Background  &                   & T90     & It's just so interesting how obvious that is when I watch DauT, whereas with other players it might be about other aspects of their game.                                                                                \\
    & Strategy Discussion   &                      & Analytical  &                   & Nili    & What surprises me is the balls that he is showing over and over again.                                                                                                                                                   \\
    & State Update          & Base Design                     & Domain      & what              & Nili    & He is so open at home.                                                                                                                                                                                                   \\
    &                       &                      &             &                   & T90     & Yeah,                                                                                                                                                                                                                    \\
    & State Update          & Military State                     & Descriptive   & what              & Nili    & He doesn't have the craziest army count.                                                                                                                                                                                 \\
    & State Update          & Military State       & Descriptive   & what              & Nili    & He does have only one defensive tower.                                                                                                                                                                                   \\
    & State Update          & Base Design         & Domain      & what              & Nili    & Stone is open, wood is open, farms are open, and he's thinking about clicking up.                                                                                                                                        \\
    & Strategy Discussion   &                      & Analytical  & why did           & T90     & Yeah, I feel like this might come back to how much control he feels he has.                                                                                                                                              \\
    & Event Description     & Age Up           & Descriptive   & what              & T90     & And as he's on his way to Castle Age.                                                                                                                                                                                    \\
    & Recap                 &                      & Descriptive   & what              & T90     & We quickly went from, "Well, Villese is in a nice position, he's denying resources, he's being aggressive," to "What does Villese do now, Nili?". Crazy.                                                                 \\
    & Strategy Discussion   &                      & Domain      & how to            & T90     & He doesn't need walls because of his understanding of the game.                                                                                                                                                          \\
    & Strategy Discussion   &                      & Descriptive   & what              & Nili    & The crazy thing is, DauT is getting the tower by his gold.                                                                                                                                                               \\
    & Strategy Prediction   &                      & Analytical  & what could happen & Nili    & He just needs one other defensive tower at one of his woodlines, and then those scouts shouldn't be able to do too much, but Villese is going for a massive investment.                                                  \\
    & Event Description     & Tech Progress          & Descriptive   & what              & Nili    & Look at that: Bloodlines and Scale Barding Armor.                                                                                                                                                                        \\
    & Event Description     & Tech Progress           & Descriptive   & what              & T90     & Okay, so Villese is investing more.                                                                                                                                                                                      \\
    & Event Description     & Eco State            & Descriptive   & what              & T90     & He also had less food income for a while.                                                                                                                                                                                \\
    & Strategy Prediction   &                      & Analytical  & what could happen & T90     & So Villese's opportunity would be all-in Feudal if he sees that DauT is up.                                                                                                                                              \\
    & Strategy Prediction   &                      & Analytical  & what could happen & T90     & After clearing up this army, which should happen soon... surprised he's waiting so much then.                                                                                                                            \\
    & Strategy Prediction   &                      & Domain      & what could happen & T90     & He'll probably need to add spearmen and scouts and archers and everything he can in Feudal to stop DauT, who will probably make scorpions and knights.                                                                   \\
    & Event Description     & Battle               & Domain      & what              & T90     & And Villese, you've got to engage! There it is.                                                                                                                                                                          \\
    & Event Description     &                      & Descriptive   & what\footnote{"what" because it describes a house actually being destroyed}                & Nili    & It feels like a house is going down.                                                                                                                                                                                   \\
    & Event Description     &                      & Domain      & what could happen & Nili    & Jumps out there, and that shouldn't be a fair fight. I think Villese should clean this one up.                                                                                                                           \\
    & Event Description     &                      & Domain      & what could happen & T90     & Yeah, he should easily clear that, but I think DauT will be happy with how long it took Villese to decide to take that fight, and it wasn't quite as convincing as I thought it would be either.                         \\
    & Event Description     & Tech State           & Descriptive   & what              & Nili    & DauT had Forging, was very tightly positioned, and the scouts couldn't find the craziest engagement.                                                                                                                     \\
    & Event Description     & Military State       & Descriptive   & what              & Nili    & Yeah, that is actually quite an even trade-off.                                                                                                                                                                          \\
    & Event Description     &                      & Analytical  & judgement         & Nili    & Not what Villese needed.                                                                                                                                                                                                 \\
    & Event Description     &                      & Descriptive   & judgement         & T90     & No. Oh my goodness, like more investment into scouts and upgrades, and Villese only survives with two scouts, one of which is on 13 HP.                                                                                  \\
    & Strategy Discussion   &                      & Analytical  & judgement         & T90     & Yeah, I guess it was wise that he waited then because it could have been worse without Scale Barding, but holy cow, he's not going to be happy to see that DauT is in Castle Age.                                        \\
    & Strategy Discussion   &                      & Analytical  & what could happen & Nili    & Yeah, the difference in Castle Age will be like three minutes or something.                                                                                                                                              \\
    & Event Description     & Movement             & Descriptive   & what              & Nili    & DauT now sends another scout forward and already has some solid upgrades on those.                                                                                                                                       \\
    & Event Description     &                      & Descriptive   & what              & Nili    & He will now take the fight and, yeah, will punish the low HP scout.                                                                                                                                                      \\
    & Strategy Prediction   &                      & Analytical  & how to            & T90     & So, how worried should DauT be about the towers? Because if the game goes on and Villese is just going to wait for the resources to go to Castle, but if the game goes on, like these towers are denying gold currently. \\
    & Strategy Prediction   &                      &             &                   & T90     & You think he's real concerned there, or do you think that is not an issue for him at all?                                                                                                                                \\
    & Strategy Prediction   &                      & Analytical  & what could happen & Nili    & If they were completely walled, he should be concerned, but yeah,                                                                                                                                                        \\
    & Event Description     & Battle               & Descriptive   & what              & Nili    & for some reason, Villese has an opening to his tower there, and the knight is now happily chopping away.                                                                                                                 \\
    & Strategy Discussion   &                      & Analytical  & why did           & T90     & I think he had to save the vills because they were building the tower on the right,                                                                                                                                      \\
    & Game Discussion       &                      & Domain      & why didn't \footnote{\label{foot:hypothetical}Hypothetical, so why didn't instead of what could happen}   & Nili    & He can open a diagonal wall. The villager can jump in, but the knight has such a tougher time to attack.                                                                                         \\                       \\
    & Game Discussion       &                      & Domain      &                   & T90     & Yeah, I was going to say you couldn't garrison if it was fully walled, but that's a smart play. A smart play from you, if you were playing anyways.                                                                      \\
    & Game Discussion       &                      & Domain      & what              & T90     & Yeah, Villese is immediately on the back foot.                                                                                                                                                                           \\
    & Game Discussion       &                      & Background  & what              & T90     & Another game where it was pretty competitive, but DauT is looking really clean                                                                                                                                           \\
    & Event Description     & Building Production  & Descriptive   & what              & T90     & here and he's adding a second TC.                                                                                                                                                                                        \\
    & Event Description     & Civilization Context & Descriptive   & what              & T90     & So not only with the Khmer farms but also with the extra villagers soon.                                                                                                                                                 \\
    & Event Description     & Military State       & Descriptive   & what              & T90     & And he has the knights.                                                                                                                                                                                                  \\
    & Event Description     &                      & Descriptive   & judgement         & T90     & Huge lead for DauT.                                                                                                                                                                                                      \\
    & Event Description     & Battle               & Descriptive   & what              & Nili    & Yeah, and oh! What, what, what? How did that CA get in? Where...? What is that?                                                                                                                                          \\
    & Event Description     &                      & Descriptive   & what              & Nili    & It's fully walled.                                                                                                                                                                                                       \\
    & Event Description     &                      & Descriptive   & what              & T90     & It's in the middle of the house foundation, Nili. He just wanted a home.                                                                                                                                                 \\
    & Event Description     &                      & Descriptive   & what              & Nili    & and he can't move the CA. I think he clicked it twice or three times already.                                                                                                                                            \\
    & Joke                  &                      & Domain      & how to            & T90     & No wonder you never see cavalry archers from Huns nowadays. They have no houses to hop into!                                                                                                                             \\
    & Event Description     &                      & Analytical  & what              & T90     & What is this? That actually helped DauT to break through possibly, but now it hurts him if he wants to CA, but I think he'll be happy with this.                                                                         \\
    & Technical Issues      &                      &             &                   & T90     & Oh my god, that's ridiculous. Thanks DE, and oh,                                                                                                                                                                         \\
    & Event Description     &                      & Descriptive   & what              & Nili    & the knights are in at the other side as well, Tristan (T90)                                                                                                                                                              \\
    & Event Description     &                      & Domain      & what could happen & T90     & That's gotta be game over, Nili.                                                                                                                                                                                         \\
    & Event Description     &                      & Domain      & judgement         & T90     & I mean, I was feeling like it was game over before DauT got in, but now this is insane.                                                                                                                                  \\
    & Event Description     &                      & Descriptive   & what              & Nili    & Military count 14 to zero.                                                                                                                                                                                               \\
    & Event Description     &                      & Descriptive   & what              & Nili    & 39 villagers, 20 of them idle.                                                                                                                                                                                           \\
    & Event Description     &                      & Descriptive   & what              & Nili    & DauT is just going for the kill of every single one of them under the TC.                                                                                                                                                \\
    & Strategy Prediction   &                      & Descriptive   & what              & T90     & Yeah, and of course, three TC eco soon, I would imagine, but I see only two.                                                                                                                                             \\
    & Game Discussion       &                      &             & judgement         & T90     & Villese, I feel bad for him because of that bug, but I'm not sure how much it would have mattered here, Nili.                                                                                                            \\
    & Strategy Prediction   &                      & Analytical  & how to\footnote{"how-to" since its too predictive (but also trying to keep the hype up in a game that is over)}            & T90     & He's going to have to, if he wants to continue to play here which he is doing, get quite a few camels out, and he'll need to clear up these knights and then have a huge counter-attack.                                 \\
    & Strategy Prediction   &                      & Domain      & what could happen & Nili    & But how can you fuel that counter-attack?                                                                                                                                                                                \\
    & State Update          & Eco State                     & Descriptive   & what              & Nili    & You don't have a single villager on gold.                                                                                                                                                                                \\
    & State Update          & Eco Comparison                     & Descriptive   & what              & Nili    & You now have some farmers as well, but he will be 20 villagers behind with an even army.                                                                                                                                 \\
    & Strategy Discussion   &                      & Domain      & judgement         & Nili    & That's not a thing. That's not a recipe for a comeback.                                                                                                                                                                  \\
    & Strategy Discussion   &                      & Domain      & what could happen & T90     & Yeah, I, um, I don't have any answers for you. I'm just saying he needs a counter-attack, but I don't know how he's going to do it.                                                                                      \\
    & Strategy Prediction   &                      & Descriptive   & what              & T90     & This is what he needs. Um, okay, he re-walls it but pays the price and loses a vill.                                                                                                                                     \\
    & Strategy Prediction   &                      & Domain      & judgement         & T90     & And DauT is like, "Chase me, as I'll send in more units and I'll continue to boom." This is very well played.                                                                                                            \\
    & Game Discussion       &                      & Domain      & what              & Nili    & Villese can't even kill that CA now, right? Unless he deletes that house, because he can never get to that CA.                                                                                                           \\
    & Game Discussion       &                      &             & what could happen & T90     & Delete it, DauT, come on, delete it. Delete the cavalry archer.                                                                                                                                                          \\
    & Event Description     & Battle               & Descriptive   & what              & Nili    & He gets a kill there. Oh god.                                                                                                                                                                                            \\
    & Game Discussion       &                      & Analytical  & why didn't        & T90     & I feel like, I mean, it doesn't really matter again, but I feel like if that would ever happen, I would delete my bugged unit.                                                                                           \\
    & Game Discussion       &                      &             & why didn't        & T90     & But then again, it's a massive tournament, so, uh, take advantage of what you can get, I suppose.                                                                                                                        \\
    & State Update          &  Eco Comparison                    & Descriptive   & what              & T90     & Fortunately, it didn't really matter in the grand scheme of things. As there is 32 vills for Villese and 56 for DauT.                                                                                                    \\
    & Player Analysis       &                      & Domain      & judgement         & Nili    & Villese just doesn't want to give it up. Villese is doing a good job of holding, but it literally is only holding.                                                                                                       \\
    & Player Analysis       &                      &             &                   & T90     & Yeah.                                                                                                                                                                                                                    \\
    & Strategy Prediction   &                      & Domain      & what could happen & Nili    & There's no way of him ever pushing on the other side of the map.                                                                                                                                                           \\
    & Event Description     & Building Production  & Descriptive   & what              & Nili    & Now goes on to a monastery, still without any villagers on gold, really, and DauT's happily booming behind this one. He's in full control.                                                                               \\
    & Civilization Analysis &                      & Domain      & what              & T90     & Khmer farms are just disgusting.                                                                                                                                                                                         \\
    & State Update          & Tech State          & Descriptive   & what              & T90     & I'm watching the food count as I'm casting here. No Wheelbarrow.                                                                                                                                                         \\
    & State Update          & Eco State            & Descriptive   & what              & T90     & And the food count is just climbing and climbing and climbing and climbing.                                                                                                                                              \\
    & Strategy Prediction   &                      & Analytical  & what could happen & T90     & DauT will maybe need to heal up a bit with his monks, get a few conversions, but he has to maintain that eco lead which is nearing double and probably will get to double what Villese has.                              \\
    & Strategy Prediction   &                      &             &                   & Nili    & That's so ugly. Let's take a look.                                                                                                                                                                                       \\
    & Event Description     & Battle               & Descriptive   & what              & Nili    & Camels trying to find something. DauT, how on point is he? Oh, he doesn't even care.                                                                                                                                     \\
    & Event Description     & Battle               & Descriptive   & what              & Nili    & You see it, it just will be built. Knights are being pulled away, and for Villese, well, this is kind of the moment where you said he needed to make something happen.                                                   \\
    & Strategy Discussion   &                      & Domain      & what could happen & Nili    & The counter-attack should have come in, but can't find any damage.                                                                                                                                                       \\
    & Strategy Discussion   &                      & Domain      & how to            & T90     & Yeah, just not possible. I mean, we're getting to the point where DauT will have enough knights to take out the camels there, so I think we're at that point officially with all the numbers.                            \\
    & Event Description     & Battle               & Descriptive   & what              & T90     & So, it's unfortunate for Villese as he... oh my god, look at the camel in the house!                                                                                                                                     \\
    & Event Description     & Battle               & Descriptive   & what              & T90     & He had to delete the house to take out the cavalry archer.                                                                                                                                                               \\
    & Event Description     &                      &             &                   & Nili    & Ah, okay,                                                                                                                                                                                                                \\
    & Event Description     & Military State       & Descriptive   & what              & T90     & and the camel's now on 4 HP.                                                                                                                                                                                             \\
    & Strategy Prediction   &                      & Analytical  & what could happen & Nili    & Oh goddi. Uh, what is the next step here for Villese? Like, is he trying to play this one TC?                                                                                                                            \\
    & Strategy Prediction   &                      & Analytical  & what could happen & T90     & Yeah, one TC, give it your best shot, all-in camel.                                                                                                                                                                      \\
    & Strategy Prediction   &                      & Analytical  & what could happen & T90     & I would normally say camel-pike or something, but he can't even afford to get the Pikeman upgrade and make camels at the same time.                                                                                      \\
    & Strategy Prediction   &                      & Descriptive   & why did           & T90     & So, can we talk about how Villese has stables and monasteries where there should be farms? That I don't understand.                                                                                                      \\
    & Joke                  &                      &             &                   & Nili    & Do you really want to discuss farm placement?                                                                                                                                                                            \\
    & Joke                  &                      &             & why did           & T90     & Okay, fair enough. I opened myself up to that joke, but it's a bit suspect how he's placed two buildings where farms should be.                                                                                          \\
    & Joke                  &                      & Background  & why did           & T90     & Not as bad as Nicov building a mill directly next to his Town Center yesterday, that was funny, but I think there's a lot going on, safe to say, Nili. I can't focus on the specifics.                                   \\
    & Strategy Discussion   &                      & Domain      & why did           & Nili    & Well, he kind of had to rush those out.                                                                                                                                                                                  \\
    & Strategy Discussion   &                      & Analytical  & why did           & Nili    & Wanted to have them not too far in the back but also not scouted at the front.                                                                                                                                           \\
    & Event Description     & Building Production  & Analytical  & what              & Nili    & The monastery obviously has a really weird timing there.                                                                                                                                                                 \\
    & Event Description     & Eco State            & Descriptive   & what              & Nili    & and DauT is now at 4 TCs 35 villager lead.                                                                                                                                                                               \\
    & Strategy Prediction   &                      & Domain      & why did           & Nili    & He doesn't need to take any fights, right?                                                                                                                                                                               \\
    & Strategy Discussion   &                      & Analytical  & why did           & Nili    & He knows, "Okay, I just need to sit back and there's no way of you winning."                                                                                                                                             \\
    & Event Description     & Building Production  & Analytical  & what              & T90     & Yep, exactly. And he's even going for the safe approach of adding barracks                                                                                                                                               \\
    & Strategy Discussion   &                      & Domain      & why did           & T90     & so he can go for some pikes, just in case the camels were to be too much.                                                                                                                                                \\
    & Strategy Discussion   &                      & Domain      & how to            & T90     & But camels are a pretty poor raiding unit, so at least it's half camel, half knight now.                                                                                                                                 \\
    & Game Discussion       &                      & Analytical  &                   & T90     & I'm just... Nili, we should talk about the next game at this point because this game feels very over.                                                                                                                    \\
    & Game Discussion       &                      & Analytical  & what could happen & T90     & This would be the throw of all throws if DauT were to lose this.                                                                                                                                                         \\
    & Game Discussion       &                      & Experiential &                   & Nili    & Yeah, that is sometimes the funny situation as a caster, right? We kind of don't want to lie to our viewers and we kind of want to say, "Okay, what are the next steps to come back?"                                    \\
    & Game Discussion       &                      & Experiential &                   & Nili    & Obviously for us, we get the intel, we see everything, we know how bad it looks.                                                                                                                                         \\
    & Game Discussion       &                      & Experiential & why didn't        & Nili    & For Villese, he thinks, "Okay, maybe it's a 1k score lead, maybe he's explored 50\% more of the map, maybe he has some floating resources or something."                                                                 \\
    & Game Discussion       &                      & Experiential & why did           & Nili    & It is just very easy for us to see. Villese is still fighting in the biggest tournament spot he's ever had.                                                                                                              \\
    & Game Discussion       &                      & Background  &                   & T90     & Yeah, exactly. Yeah, he's a player who has been around, but to make it this far, I think is, uh, it might be the biggest tournament spot he's ever had, right?                                                           \\
    & Game Discussion       &                      & Background  &                   & T90     & Yeah, that and NAC, which was like, kind of group stage and whatnot. But yeah, I see some raids from both players.                                                                                                       \\
    & Event Description     & Eco State            & Descriptive   & what              & T90     & 90 vills against 48 right now. Um, DauT's monks are actually going home, right?                                                                                                                                          \\
    & Event Description     & Movement             & Descriptive   & what              & T90     & Yep, his four monks are going home to clear that up and maybe send the pikemen that way.                                                                                                                                 \\
    & Strategy Discussion   &                      & Domain      & judgement         & T90     & I can't say that DauT has done the best job of closing this out, but I also don't feel like these raids really matter all that much.                                                                                     \\
    & Strategy Discussion   &                      & Descriptive   & what              & Nili    & Yeah, villager advantage now down to only 35 again after he was already at a 45 lead.                                                                                                                                    \\
    & Strategy Discussion   &                      & Domain      & judgement         & Nili    & Well, DauT isn't playing perfectly at the moment, but it can't matter.                                                                                                                                                   \\
    & Joke                  &                      & Domain      & what              & Nili    & Like, yeah, he actively needs to leave his PC to not win this game.                                                                                                                                                      \\
    & Player Analysis       &                      & Experiential &                   & T90     & I feel like that's a very DauT thing, though.                                                                                                                                                                            \\
    & Player Analysis       &                      & Experiential &                   & T90     & Like, what DauT does well is not necessarily... um, how do I phrase this?                                                                                                                                                \\
    & Player Analysis       &                      & Background  &                   & T90     & DauT isn't like... think back to how Lierrey played against Nicov yesterday. It was brutal, non-stop pressure, just suffocating.                                                                                         \\
    & Player Analysis       &                      & Experiential & what              & T90     & DauT's not really like that.                                                                                                                                                                                             \\
    & Player Analysis       &                      & Experiential &                   & T90     & DauT is like, "Oh, I see that. I'm gonna take control of that area or that position," and he slowly gets ahead in games.                                                                                                 \\
    & Player Analysis       &                      & Experiential &                   & T90     & So he's officially ahead now, and so he's like, "All right, I'm just gonna take my time," and he might not be as quick to close out the game as some of the faster players.                                              \\
  \end{longtable}
} % end scriptsize

{\scriptsize

  \setcounter{AppRow}{0}
  \begin{longtable}{@{}>{\stepcounter{AppRow}\theAppRow}r|L{2.3cm}L{1.0cm}U{11.55cm}@{}}
    \caption[Co-cast coding Villese vs DauT G2]{Coding of Lidakor casting Villese vs DauT -- King of the Desert 3 quarterfinals Game 2}\label{tab:lidakor-coding}\\
    \multicolumn{4}{@{}l}{\footnotesize Source: \url{https://www.youtube.com/watch?v=f11soGHv8H0&t=1923s}}\\[\AfterSourceLinksVSpace]
    \toprule
    \multicolumn{1}{r}{\textbf{\#}} &
    \makecell[l]{\textbf{Type}} &
    \makecell[l]{\textbf{Speaker}} &
    \multicolumn{1}{l}{\makecell[l]{\textbf{Utterance}}} \\
    \midrule
    \endfirsthead
    \multicolumn{4}{@{}l}{\textit{Table \thetable\ (continued)}}\\[0.3em]
    \toprule
    \multicolumn{1}{r}{\textbf{\#}} &
    \makecell[l]{\textbf{Type}} &
    \makecell[l]{\textbf{Speaker}} &
    \multicolumn{1}{l}{\makecell[l]{\textbf{Utterance}}} \\
    \midrule
    \endhead

    \midrule
    \multicolumn{4}{r}{\small Continued on next page} \\
    \endfoot

    \bottomrule
    \endlastfoot
    & Match Introduction    & Lidakor & A Malians versus Khmer game over here.                                                                                                                                                                                                                                                          \\
    & Match Introduction    & Lidakor & Interestingly, it's somewhat unexpected to see Villese picking Malians here whereas, um, I think that DauT's Khmer here is just pretty straightforward.                                                                                                                                         \\
    & Match Introduction    & Lidakor & Let's see what DauT plans here.                                                                                                                                                                                                                                                                 \\
    & Civilization Analysis & Lidakor & There are absolutely no guarantees he is going to play with, um, scouts here.                                                                                                                                                                                                                   \\
    & Civilization Analysis & Lidakor & In fact, he could definitely go for the surprise archers play here.                                                                                                                                                                                                                             \\
    & Civilization Analysis & Lidakor & DauT on the right side in red as the Khmer; left side is going to be Villese as Malians in blue. So let's take a look.                                                                                                                                                                          \\
    & Map Analysis          & Lidakor & These maps are definitely a lot more open than the previous map generation.                                                                                                                                                                                                                     \\
    & Map Analysis          & Lidakor & So technically, this is defendable for Villese, and I would make an argument for walling here in front, here, and here.                                                                                                                                                                         \\
    & Map Analysis          & Lidakor & Not an easy wall-off, but technically this is a defendable base for Villese.                                                                                                                                                                                                                    \\
    & Map Analysis          & Lidakor & On the other side, um, for DauT, this base isn't that great, and the main reason for that is because there is no absolutely safe gold here for DauT.                                                                                                                                            \\
    & Map Analysis          & Lidakor & So, an easy wall-off on the north, that's nice.                                                                                                                                                                                                                                                 \\
    & Map Analysis          & Lidakor & Stone at the back, woodlines at the back; however, there is no safe gold mine.                                                                                                                                                                                                                  \\
    & Map Analysis          & Lidakor & I think the best thing that you could do if you are DauT is, uh, wall like this and maybe use this gold mine.                                                                                                                                                                                   \\
    & Map Analysis          & Lidakor & The problem with this gold mine is that it takes a lot of time to walk from the TC to this gold mine, and that's something that DauT might not want to do.                                                                                                                                      \\
    & Map Analysis          & Lidakor & So there's a chance DauT will try to take this gold mine, but Villese, seeing that gold mine, could heavily consider playing archers over here.                                                                                                                                                 \\
    & Event Description     & Lidakor & He's got four on wood;                                                                                                                                                                                                                                                                          \\
    & Strategy Prediction   & Lidakor & that is definitely an archer play here,                                                                                                                                                                                                                                                         \\
    & Event Description     & Lidakor & as we are going to have a barracks coming up rather fast here from Villese.                                                                                                                                                                                                                     \\
    & Event Description     & Lidakor & DauT is going to deny that or at least delay the timing, but this was trying to be a fast barracks even before mill.                                                                                                                                                                            \\
    & Strategy Discussion   & Lidakor & If you take a look at Villese's base, he tried to surprise DauT, and that's one of the reasons why he went for four on wood at the beginning.                                                                                                                                                   \\
    & Strategy Discussion   & Lidakor & Because, uh, he wanted to have that wood in, even with the cheaper Malian buildings, for a fast barracks,                                                                                                                                                                                       \\
    & Event Description     & Lidakor & but DauT sees the barracks and he also managed to delay that a tiny bit.                                                                                                                                                                                                                        \\
    & Civilization Analysis & Lidakor & And we didn't even talk about the fact that DauT can jump into houses, so that rush should not really do a lot.                                                                                                                                                                                 \\
    & State Update          & Lidakor & As DauT is adding villagers to gold here quite fast,                                                                                                                                                                                                                                            \\
    & Strategy Prediction   & Lidakor & this looks like, um, maybe a men-at-arms build from DauT here.                                                                                                                                                                                                                                  \\
    & Strategy Discussion   & Lidakor & A 20-pop men-at-arms would be a little bit too tight of a build order                                                                                                                                                                                                                           \\
    & State Update          & Lidakor & but that's already three villagers added to gold here.                                                                                                                                                                                                                                          \\
    & State Update          & Lidakor & Although on the other hand, he's not dropping a barracks; he doesn't even have the wood for it.                                                                                                                                                                                                 \\
    & Strategy Prediction   & Lidakor & So this might be double-range straight archers from DauT because, um, if he wanted to play men-at-arms, he would be dropping the barracks by now, which is something he's not doing.                                                                                                            \\
    & Civilization Analysis & Lidakor & And since Khmer do not need the barracks to build an archery range,                                                                                                                                                                                                                              \\
    &                      & Lidakor & you just, uh, are completely fine playing with the archery ranges to start with.                                                                                                                                                                                                                 \\
    & Event Description     & Lidakor & So with all of this said and done, DauT also pushes in some deer.                                                                                                                                                                                                                               \\
    & Civilization Analysis & Lidakor & I like that he walled in his woodline, and one of the reasons why Khmer is great is because you can hop into the house from one side and hop out on the other side.                                                                                                                             \\
    & Civilization Analysis & Lidakor & And the reason why it's great is because you don't need to open these walls to add more villagers to the woodline.                                                                                                                                                                              \\
    & Event Description     & Lidakor & DauT is adding, uh, just one range right now, whereas,                                                                                                                                                                                                                                          \\
    & Strategy Discussion   & Lidakor & so that is a men-at-arms towers play from Villese,                                                                                                                                                                                                                                              \\
    & Event Description     & Lidakor & in fact, a forward archery range.                                                                                                                                                                                                                                                                \\
    & Strategy Prediction   & Lidakor & So probably a few skirmishers will be out right now.                                                                                                                                                                                                                                            \\
    & Civilization Analysis & Lidakor & Malian men-at-arms will have, uh, extra pierce armor,                                                                                                                                                                                                                                           \\
    & Event Description     & Lidakor & and it appears that Villese is going to be able to force DauT off from gold.                                                                                                                                                                                                                    \\
    & Civilization Analysis & Lidakor & Not so great for DauT over here, and that extra pierce armor really helps here with the two armor already.                                                                                                                                                                                      \\
    & Strategy Prediction   & Lidakor & If Villese mixes in a few skirmishers, DauT could be in trouble because he doesn't really have a lot to deal with those men-at-arms.                                                                                                                                                            \\
    & Event Description     & Lidakor & So, an immediate Blacksmith from DauT.                                                                                                                                                                                                                                                          \\
    & Strategy Prediction   & Lidakor & He wants to get the fast Fletching.                                                                                                                                                                                                                                                             \\
    & Strategy Prediction   & Lidakor & He has enough gold to get Fletching and a few more archers, but he needs to micro this precisely.                                                                                                                                                                                               \\
    & Strategy Prediction   & Lidakor & If he is able to get rid of those men-at-arms, I think he's going to be fine because then he can use his scout to just push away the initial skirmishers from his opponent.                                                                                                                     \\
    & Strategy Discussion   & Lidakor & But he needs to deal with those men-at-arms,                                                                                                                                                                                                                                                    \\
    & Strategy Prediction   & Lidakor & and Villese will probably try to tower this gold mine.                                                                                                                                                                                                                                          \\
    & State Update          & Lidakor & Villese has very detailed scouting on DauT's base. He sees the single woodline,                                                                                                                                                                                                                 \\
    & Event Description     & Lidakor & and that's an absolutely excellent tower.                                                                                                                                                                                                                                                       \\
    & Strategy Discussion   & Lidakor & It denies the gold, denies the woodline over here. This is going to be a very, very hard tower to deal with if you are DauT.                                                                                                                                                                    \\
    & Event Description     & Lidakor & And DauT is going to try and counter-tower that.                                                                                                                                                                                                                                                \\
    & Strategy Discussion   & Lidakor & DauT does have safe stone in the back, but no villagers are on it just yet.                                                                                                                                                                                                                     \\
    & Event Description     & Lidakor & Scout micro happening here.                                                                                                                                                                                                                                                                     \\
    & Strategy Discussion   & Lidakor & It is extremely important for Villese to kill DauT's scout because that is the main danger to his skirmishers.                                                                                                                                                                                  \\
    & Strategy Discussion   & Lidakor & And this is looking very good for Villese here. This is looking absolutely amazing, to be honest.                                                                                                                                                                                               \\
    & Event Description     & Lidakor & So on the other hand, DauT is going to have the tower up as well,                                                                                                                                                                                                                               \\
    & Event Description     & Lidakor & and both players have Fletching                                                                                                                                                                                                                                                                 \\
    & Strategy Discussion   & Lidakor & so the towers are completely even to each other here.                                                                                                                                                                                                                                           \\
    & Strategy Prediction   & Lidakor & And what DauT wants is a counter-attack.                                                                                                                                                                                                                                                        \\
    & Strategy Prediction   & Lidakor & I like this. For a moment, at least, he thought about a chance of a counter-attack because if you are fully walled, well at least, um, not really fully walled, but if you are in a defendable position, you could actually argue for a counter-attack to try and hit your forwarding opponent. \\
    & Strategy Prediction   & Lidakor & Because one of the best ways to deal with a forward is to counter-attack.                                                                                                                                                                                                                       \\
    & Event Description     & Lidakor & On the other hand, Villese is committing a lot of effort into walling his base here                                                                                                                                                                                                             \\
    & Strategy Discussion   & Lidakor & , but the right side is wide open, so that might not matter.                                                                                                                                                                                                                                    \\
    & Strategy Prediction   & Lidakor & The problem for DauT is that he's going to have to take this fight. He needs to snipe that scout first, I believe, but this is not looking great for DauT by any means.                                                                                                                         \\
    & Strategy Prediction   & Lidakor & He is going to lose all those forces.                                                                                                                                                                                                                                                           \\
    & State Update          & Lidakor & On the other hand, the men-at-arms are away from DauT's base,                                                                                                                                                                                                                                   \\
    & Strategy Prediction   & Lidakor & so this is the moment when DauT just wants to go out with the villagers and rush down the tower and take his gold mine back.                                                                                                                                                                    \\
    & Strategy Prediction   & Lidakor & So for now, DauT, I think, is still fine. He can go for a few skirmishers and just go for the gold mine now.                                                                                                                                                                                    \\
    & Strategy Prediction   & Lidakor & Why isn't DauT going for a gold mine? It's a mystery to me.                                                                                                                                                                                                                                     \\
    & Strategy Prediction   & Lidakor & He's not even going for stone and, like, a Market to sell stone.                                                                                                                                                                                                                                \\
    & State Update          & Lidakor & So it appears that he just wants to play with full skirmishers,                                                                                                                                                                                                                                 \\
    & Strategy Discussion   & Lidakor & but the problem with the full skirmisher play here is that the men-at-arms, even without the extra pierce armor of Malians, are just virtually resistant to skirmisher fire.                                                                                                                    \\
    & Strategy Prediction   & Lidakor & You need one or two archers to deal with them,                                                                                                                                                                                                                                                  \\
    & State Update          & Lidakor & and that's not something that DauT has now.                                                                                                                                                                                                                                                     \\
    & Event Description     & Lidakor & He's adding villagers to gold.                                                                                                                                                                                                                                                                  \\
    & State Update          & Lidakor & Remember, he's got 12 on food, so despite the fact that                                                                                                                                                                                                                                         \\
    & Strategy Discussion   & Lidakor & , um, he got pressured quite hard over here, we shouldn't overlook, um,                                                                                                                                                                                                                         \\
    & State Update          & Lidakor & the situation of DauT having 13 on food with Horse Collar.                                                                                                                                                                                                                                      \\
    & Strategy Prediction   & Lidakor & This means that he could have a very, very good Castle Age time over here                                                                                                                                                                                                                       \\
    & Strategy Discussion   & Lidakor & even without a huge amount of gold in the bank.                                                                                                                                                                                                                                                 \\
    & Event Description     & Lidakor & Now, he's trying to snipe down the ranged units,                                                                                                                                                                                                                                                \\
    & Strategy Prediction   & Lidakor & which I once again really like, because if he can just wall the villagers in here on the gold mine, then the men-at-arms aren't really going to be efficient.                                                                                                                                   \\
    & Event Description     & Lidakor & He has a stable over here.                                                                                                                                                                                                                                                                      \\
    & Event Description     & Lidakor & He's going to try and use the villagers to force away the men-at-arms.                                                                                                                                                                                                                          \\
    & Strategy Discussion   & Lidakor & DauT's defense here is looking pretty good.                                                                                                                                                                                                                                                     \\
    & Strategy Prediction   & Lidakor & He will be able to add one or two more archers,                                                                                                                                                                                                                                                 \\
    & Strategy Discussion   & Lidakor & and skirmishers aren't really threatening to the villagers of DauT.                                                                                                                                                                                                                             \\
    & Event Description     & Lidakor & On the other hand, there are the men-at-arms coming in.                                                                                                                                                                                                                                         \\
    & Event Description     & Lidakor & Very nice micro from DauT picking off the archers.                                                                                                                                                                                                                                              \\
    & Event Description     & Lidakor & And now an archer comes in.                                                                                                                                                                                                                                                                     \\
    & Event Description     & Lidakor & There is nothing for Villese to kill those archers with right now.                                                                                                                                                                                                                              \\
    & Event Description     & Lidakor & So the men-at-arms are dropping in HP, and once the men-at-arms are dead, I think DauT is fine.                                                                                                                                                                                                 \\
    & Event Description     & Lidakor & He is, I think, going to lose one villager over here.                                                                                                                                                                                                                                           \\
    & Event Description     & Lidakor & Sloppy micro over there from DauT                                                                                                                                                                                                                                                               \\
    & Technical Issues      & Lidakor & and immediately after, we will have a pause.                                                                                                                                                                                                                                                    \\
    & Event Description     & Lidakor & So the thing is that now Villese is about to lose all three of those men-at-arms,                                                                                                                                                                                                               \\
    & Strategy Discussion   & Lidakor & and if that happens, Villese only has skirmishers and archers.                                                                                                                                                                                                                                  \\
    & Strategy Discussion   & Lidakor & DauT can defend against that one very nicely,                                                                                                                                                                                                                                                   \\
    & Strategy Prediction   & Lidakor & and, uh, if that happens, I think DauT can just rush down this tower as well.                                                                                                                                                                                                                   \\
    & Strategy Prediction   & Lidakor & He could just send five villagers to take this down.                                                                                                                                                                                                                                            \\
    & Event Description     & Lidakor & Indeed, Villese realizes that and he's just going to wall in.                                                                                                                                                                                                                                   \\
    & Strategy Prediction   & Lidakor & On the other hand, DauT can just go for this gold mine if he really wants to, or go for stone and sell stone.                                                                                                                                                                                   \\
    & State Update          & Lidakor & He doesn't need a lot of gold, um, to get into Castle Age,                                                                                                                                                                                                                                      \\
    & Strategy Prediction   & Lidakor & and once he gets into Castle Age, he can just go for knights or mangonels and deal with this tower aggression.                                                                                                                                                                                  \\
    & Recap                 & Lidakor & So at this point, I think this is still looking better for DauT.                                                                                                                                                                                                                                \\
    & Event Description     & Lidakor & He's going to have a counter-attack. It's just a few skirmishers, but it has two archers and the scout inside, so this could be dangerous.                                                                                                                                                      \\
    & Strategy Prediction   & Lidakor & I think the best it can achieve is potentially pushing away the villagers of Villese from the woodline.                                                                                                                                                                                         \\
    & Strategy Prediction   & Lidakor & Let's see if DauT is going to spot that stable. That is a hidden stable from Villese, and Villese is massing scouts.                                                                                                                                                                            \\
    & Strategy Prediction   & Lidakor & So this could still be dangerous for DauT because he's almost exclusively skirmishers and he doesn't even have a barracks for spearmen.                                                                                                                                                         \\
    & Strategy Prediction   & Lidakor & DauT, on the other hand, is going to see the stable with the flags.\footnote{Flags indicated garissoned units}                                                                                                                                                                                                                             \\
    & Strategy Prediction   & Lidakor & Let's see if he's going to react with a barracks immediately.                                                                                                                                                                                                                                   \\
    & Strategy Prediction   & Lidakor & He could also react with his own scouts, which could work just as fine.                                                                                                                                                                                                                         \\
    & State Update          & Lidakor & Still 16 on food versus 15 on food,                                                                                                                                                                                                                                                             \\
    & Civilization Analysis & Lidakor & but it's Khmer farming versus Malian farming, so DauT's eco is still way better over here.                                                                                                                                                                                                      \\
    & Event Description     & Lidakor & So yeah, a little bit of idle time here as the woodline of Villese is getting harassed.                                                                                                                                                                                                         \\
    & Strategy Prediction   & Lidakor & And let's see if DauT is planning to sell stone here to go up.                                                                                                                                                                                                                                  \\
    & Event Description     & Lidakor & Indeed, there is the Market.                                                                                                                                                                                                                                                                    \\
    & Event Description     & Lidakor & One of the very common ways, or he can just go here. Oh, I love this from DauT!                                                                                                                                                                                                                 \\
    & Player Analysis       & Lidakor & This map awareness is so good.                                                                                                                                                                                                                                                                  \\
    & State Update          & Lidakor & I didn't even see this third gold, but now he has safe gold. And if you look at his eco, 17 on food, seven on gold...                                                                                                                                                                           \\
    & Strategy Prediction   & Lidakor & This is pretty much a straight-to-knights build's eco distribution at this stage of the game.                                                                                                                                                                                                   \\
    & Recap                 & Lidakor & So at the end of the day, DauT got pressured by towers, but it barely affected his eco.                                                                                                                                                                                                         \\
    & Event Description     & Lidakor & In fact, he's going to idle his TC a tiny bit                                                                                                                                                                                                                                                   \\
    & Strategy Prediction   & Lidakor & and potentially click up to Castle Age, whereas Villese is not even close.                                                                                                                                                                                                                      \\
    & Event Description     & Lidakor & Villese is massing a lot of scouts. This is super dangerous.                                                                                                                                                                                                                                    \\
    & Strategy Prediction   & Lidakor & This is extremely dangerous over here for DauT, because if Villese goes for Bloodline scouts here, this could be a disaster.                                                                                                                                                                    \\
    & Event Description     & Lidakor & But DauT breaks in,                                                                                                                                                                                                                                                                             \\
    & State Update          & Lidakor & and Villese is still without Bloodlines.                                                                                                                                                                                                                                                        \\
    & Strategy Prediction   & Lidakor & So if DauT just trades off his scouts with Villese's scouts, he's absolutely fine.                                                                                                                                                                                                              \\
    & Event Description     & Lidakor & There is Bloodlines coming in for Villese.                                                                                                                                                                                                                                                      \\
    & Strategy Prediction   & Lidakor & I think Villese will try to take this fight at some point.                                                                                                                                                                                                                                      \\
    & Strategy Prediction   & Lidakor & DauT might need a few spearmen. Apparently, he does not see these skirmishers killing his units.                                                                                                                                                                                                \\
    & Event Description     & Lidakor & But there is going to be Forging for DauT,                                                                                                                                                                                                                                                      \\
    & Strategy Prediction   & Lidakor & and, uh, I think by the time Villese really gets to DauT's base,                                                                                                                                                                                                                                \\
    & Strategy Prediction   & Lidakor & DauT will be up to Castle Age and he can make a few knights to clean that one up.                                                                                                                                                                                                               \\
    & Civilization Analysis & Lidakor & And also, they can jump into houses to prevent the villagers from getting killed by the scouts.                                                                                                                                                                                                 \\
    & Strategy Prediction   & Lidakor & So I don't think this scout counter-attack is going to work for Villese.                                                                                                                                                                                                                        \\
    & Strategy Discussion   & Lidakor & He's waiting too long for his Scale Barding and Bloodlines.                                                                                                                                                                                                                                     \\
    & Strategy Discussion   & Lidakor & I understand the concept of Villese here, but DauT might just break into Villese's base, which is even worse.                                                                                                                                                                                   \\
    & Event Description     & Lidakor & Because okay, now Villese needs to fight this one.                                                                                                                                                                                                                                              \\
    & State Update          & Lidakor & DauT is going to see the upgrades, that's super important.                                                                                                                                                                                                                                      \\
    & Event Description     & Lidakor & He tries to get the archers and skirmishers behind.                                                                                                                                                                                                                                             \\
    & Event Description     & Lidakor & This is going to be a cleanup, make no mistake, but is it going to be a cleanup?                                                                                                                                                                                                                \\
    & Event Description     & Lidakor & Yeah, it is going to be... okay, so it is a cleanup from Villese, but at the end of the day,                                                                                                                                                                                                    \\
    & Event Description     & Lidakor & if you look at what is left, not a lot, because DauT also had Forging on his scouts.                                                                                                                                                                                                            \\
    & Event Description     & Lidakor & It's two scouts left for Villese, and that is basically all.                                                                                                                                                                                                                                    \\
    & Event Description     & Lidakor & DauT kills a villager here on the tower of Villese                                                                                                                                                                                                                                              \\
    & Event Description     & Lidakor & , and that is a huge Castle Age difference.                                                                                                                                                                                                                                                     \\
    & State Update          & Lidakor & So, double stable...                                                                                                                                                                                                                                                                            \\
    & Strategy Prediction   & Lidakor & it could be elephants as well, indeed. So this can be nice. This can be elephants.                                                                                                                                                                                                              \\
    & Event Description     & Lidakor & It's going to be knights from DauT in this one.                                                                                                                                                                                                                                                 \\
    & Strategy Prediction   & Lidakor & He can go for Bodkin Arrow without any further issues, and, uh, he could even go for a siege workshop and mangonel these down.                                                                                                                                                                  \\
    & Strategy Discussion   & Lidakor & This is looking extremely rough right now for Villese because Villese is just about to click up to Castle Age.                                                                                                                                                                                  \\
    & Strategy Prediction   & Lidakor & He can make a few spearmen to defend against the knights, but DauT could decide on a forward siege workshop and just go for rams and knights and take down these buildings.                                                                                                                     \\
    & State Update          & Lidakor & Let's see if DauT is gonna do that. Um, there is no forward villager from DauT just yet.                                                                                                                                                                                                        \\
    & Event Description     & Lidakor & Brings in a knight and for whatever reason, there is a hole in this wall, so this knight is just going to kill the tower with no chance of Villese repairing.                                                                                                                                   \\
    & Strategy Prediction   & Lidakor & DauT is perfectly fine here. I think the only ingredient missing for him is a forward siege workshop.                                                                                                                                                                                           \\
    & Joke                  & Lidakor & Yeah, who is doubting DauT? What did he do to earn doubt?                                                                                                                                                                                                                                       \\
    & Player Analysis       & Lidakor & DauT has always been a very, very good player in terms of mindset and everything.                                                                                                                                                                                                               \\
    & Player Analysis       & Lidakor & I think what makes DauT very good here is the map awareness that he's having.                                                                                                                                                                                                                   \\
    & Player Analysis       & Lidakor & And overall, I think one thing that really distinguishes DauT from many players is how calm he is.                                                                                                                                                                                              \\
    & Player Analysis       & Lidakor & Just go to a DauT stream. Has anyone seen DauT being annoyed at all?                                                                                                                                                                                                                            \\
    & Player Analysis       & Lidakor & I don't think so. I've never seen DauT annoyed.                                                                                                                                                                                                                                                 \\
    & Player Analysis       & Lidakor & There is a villager dying for him, but he's always just chilling and whatnot.                                                                                                                                                                                                                   \\
    & Event Description     & Lidakor & The Khmer cavalry archer strategy is a good idea.                                                                                                                                                                                                                                               \\
    & Civilization Analysis & Lidakor & Khmer cavalry archers are kind of legit, and you can get to your opponent's base quite easily and deny the second layers, so it's not a bad idea by any means.                                                                                                                                  \\
    & Player Analysis       & Lidakor & So the fact that DauT is super calm helps very, very massively when you're this deep in a tournament like this.                                                                                                                                                                                 \\
    & Event Description     & Lidakor & I don't know what this cavalry archer is doing and how he teleported inside that house foundation.                                                                                                                                                                                              \\
    & Technical Issues      & Lidakor & Do you see this?                                                                                                                                                                                                                                                                                \\
    & Technical Issues      & Lidakor & Do you see this? That he's actually inside the house foundation.                                                                                                                                                                                                                                \\
    & State Update          & Lidakor & Anyways, um, at this point, it's a three-villager advantage for DauT.                                                                                                                                                                                                                           \\
    & Event Description     & Lidakor & Villese is going to be Castle Age,                                                                                                                                                                                                                                                              \\
    & Strategy Prediction   & Lidakor & but Castle Age with what?                                                                                                                                                                                                                                                                       \\
    & Historical Account    & Lidakor & And after Villese 3-0'd TaToH, a lot of people said that Villese could actually be, maybe even in the semi-finals.                                                                                                                                                                              \\
    & Historical Account    & Lidakor & It looks like he's going to be down by two games against DauT.                                                                                                                                                                                                                                  \\
    & Strategy Prediction   & Lidakor & He needs to go for camels because he's behind in knight numbers,                                                                                                                                                                                                                                \\
    & Strategy Prediction   & Lidakor & but he's not going to have the food eco.                                                                                                                                                                                                                                                        \\
    & State Update          & Lidakor & Once again, look at all these idles. There's absolutely nothing to do.                                                                                                                                                                                                                          \\
    & Event Description     & Lidakor & DauT can just sacrifice his knights beneath the TC and just keep killing villagers.                                                                                                                                                                                                             \\
    & State Update          & Lidakor & Villese is down to 35 with no gold income.                                                                                                                                                                                                                                                      \\
    & Strategy Prediction   & Lidakor & There are going to be a few camels in here, but at this point, I don't think this is gonna work.                                                                                                                                                                                                \\
    & State Update          & Lidakor & I mean, five on food, zero on gold,                                                                                                                                                                                                                                                             \\
    & Strategy Discussion   & Lidakor & you cannot replenish those camels.                                                                                                                                                                                                                                                              \\
    & State Update          & Lidakor & And those camels are +1/+1.                                                                                                                                                                                                                                                                     \\
    & State Update          & Lidakor & DauT has a lot more knights,                                                                                                                                                                                                                                                                    \\
    & Strategy Prediction   & Lidakor & and as I said, Villese cannot make more camels here. He can get this batch out and then basically, that is it.                                                                                                                                                                                  \\
    & Strategy Prediction   & Lidakor & Yeah, indeed, exactly. Like, DauT has nothing to lose, but even if he had a lot of things to lose,                                                                                                                                                                                              \\
    & Player Analysis       & Lidakor & DauT is definitely one of the more calm players out there, and that helps immensely.                                                                                                                                                                                                            \\
    & Player Analysis       & Lidakor & As much as we are talking about the fact that, let's say for Nicov, he often gets nervous in tournaments and that can hurt him, we also think about the fact that DauT barely gets very, very nervous.                                                                                          \\
    & Player Analysis       & Lidakor & And being this calm helps so, so much when you are trying to just finish off your opponent with surgical precision.                                                                                                                                                                             \\
    & Event Description     & Lidakor & He picks off one villager on the farm. The cavalry archer is going to finish off the vills, and, uh, basically, this is going to be almost impossible for Villese now.                                                                                                                          \\
    & State Update          & Lidakor & He's behind by 24 vills at this stage of the game against a Khmer player that also has a great farming eco.                                                                                                                                                                                     \\
    & Event Description     & Lidakor & Goes for a monastery,                                                                                                                                                                                                                                                                           \\
    & Strategy Prediction   & Lidakor & so we are going to see probably a few monks to convert the camels over here.                                                                                                                                                                                                                    \\
    & State Update          & Lidakor & Villese is still trying, but he's got basically no gold income                                                                                                                                                                                                                                  \\
    & State Update          & Lidakor & , so if you look at these stables, all of them are idle.                                                                                                                                                                                                                                        \\
    & Strategy Prediction   & Lidakor & If Villese had a Market, he could sell some resources to get some more gold in, but this is never going to be efficient, and DauT would just bring in the monks.                                                                                                                                \\
    & Event Description     & Lidakor & TCing that other gold is DauT.                                                                                                                                                                                                                                                                  \\
    & State Update          & Lidakor & He has reclaimed control over most of his base, and DauT also has two villagers on stone and almost 400 stone.                                                                                                                                                                                  \\
    & Strategy Prediction   & Lidakor & I think that if DauT sends two or three more villagers to stone, he could easily get a castle up in the face of his opponent.                                                                                                                                                                   \\
    & Skill Discussion      & Lidakor & Um, to answer that question on the chat, the reason why he still has two villagers on stone is because he was getting towered and he just forgot about the villagers on the stone.                                                                                                              \\
    & Strategy Prediction   & Lidakor & But right now, he has enough stone to just send two or three more vills and pile up resources for a castle.                                                                                                                                                                                     \\
    & Strategy Prediction   & Lidakor & And if that happens, he can just Castle-drop Villese over here because he is going to have the military advantage.                                                                                                                                                                              \\
    & Strategy Prediction   & Lidakor & The knights and monks over here will just be better for DauT, although he still needs to watch out against the camels of Villese.                                                                                                                                                               \\
    & Strategy Prediction   & Lidakor & I could see pikemen as well from DauT, like tech-switching to pikemen.                                                                                                                                                                                                                          \\
    & Strategy Discussion   & Lidakor & Because, uh, the thing is that Villese doesn't have the eco to tech-switch into archers to deal with the pikemen.                                                                                                                                                                               \\
    & Strategy Prediction   & Lidakor & So because of that, I think a pikemen switch here from DauT could be justified. Or he could just boom.                                                                                                                                                                                          \\
    & Strategy Discussion   & Lidakor & He knows he did a lot of damage.                                                                                                                                                                                                                                                                \\
    & State Update          & Lidakor & It's 64 against 37 eco.                                                                                                                                                                                                                                                                         \\
    & Strategy Discussion   & Lidakor & He doesn't need to take any risks.                                                                                                                                                                                                                                                              \\
    & Strategy Discussion   & Lidakor & When you are ahead by this much, you can do two things.                                                                                                                                                                                                                                         \\
    & Strategy Prediction   & Lidakor & You either, uh, are trying to finish your opponent off and, like, drop a forward castle, or you're just gonna sit back and boom, knowing that you already have a huge advantage.                                                                                                                \\
    & Strategy Discussion   & Lidakor & So there is almost, um, no chance for your opponent to catch up, and that's exactly what we're seeing.                                                                                                                                                                                          \\
    & State Update          & Lidakor & DauT is at 69 villagers with three TCs; Villese is at one TC with 40 villagers.                                                                                                                                                                                                                 \\
    & State Update          & Lidakor & Even though Villese has Wheelbarrow and DauT does not, that isn't enough to compensate for this big of a deficit.                                                                                                                                                                               \\
    & Event Description     & Lidakor & So at this point, DauT is just in a way better spot, and he's gonna start picking up relics as well.                                                                                                                                                                                            \\
    & Event Description     & Lidakor & I think he just is about to get the first one.                                                                                                                                                                                                                                                  \\
    & Event Description     & Lidakor & Looking at his scouting, he doesn't see any more.                                                                                                                                                                                                                                               \\
    & Event Description     & Lidakor & But that's the second one he's gonna grab                                                                                                                                                                                                                                                       \\
    & Strategy Prediction   & Lidakor & and I think he can scout the third one very soon.                                                                                                                                                                                                                                               \\
    & Strategy Discussion   & Lidakor & So he's gonna have a lot of relics here. It's all about just extending that lead as much as possible.                                                                                                                                                                                           \\
    & Strategy Discussion   & Lidakor & So, um, like, as much as it is, or as much as it seems okay for Villese that he has a decent amount of camels to deal with the knights,                                                                                                                                                         \\
    & Strategy Discussion   & Lidakor & he just doesn't have the eco to replenish them.                                                                                                                                                                                                                                                 \\
    & Event Description     & Lidakor & And that is indeed a double barrack switch from DauT.                                                                                                                                                                                                                                           \\
    & Strategy Prediction   & Lidakor & He's going to add pikemen, and Villese is not going to have the eco to tech-switch into crossbows or swordsmen to deal with them.                                                                                                                                                               \\
    & Strategy Prediction   & Lidakor & Maybe he can add one or two scorpions, but I don't think he can really mass anything that can deal with pikemen.                                                                                                                                                                                \\
    & State Update          & Lidakor & Plus, um, DauT also is going to have four monks.                                                                                                                                                                                                                                                \\
    & Strategy Prediction   & Lidakor & If you get four conversions on camels, that's perfect.                                                                                                                                                                                                                                          \\
    & Strategy Prediction   & Lidakor & It appears Villese is going to try and flank, but this just gives a chance for DauT to potentially push the front line over here.                                                                                                                                                               \\
    & Strategy Prediction   & Lidakor & And let's see if Villese is going to bail out and go back home here.                                                                                                                                                                                                                            \\
    & Strategy Prediction   & Lidakor & I think he needs to because he cannot afford that many knights to get into his eco.                                                                                                                                                                                                             \\
    & Strategy Prediction   & Lidakor & Whereas for DauT, he can defend with monks. He still can jump into houses, and camels aren't great as a raiding unit.                                                                                                                                                                           \\
    & Strategy Discussion   & Lidakor & The low base attack of camels makes them kind of poor as a raiding unit.                                                                                                                                                                                                                        \\
    & Event Description     & Lidakor & DauT even tries to convert the villagers to prevent the second layering over here.                                                                                                                                                                                                              \\
    & Event Description     & Lidakor & The barracks falls, the knights are going to swarm inside.                                                                                                                                                                                                                                      \\
    & Event Description     & Lidakor & Villese is trying to get some damage done within DauT's base, but that seems impossible.                                                                                                                                                                                                        \\
    & Event Description     & Lidakor & Some idle time is forced on DauT, but nothing significant, and DauT is gonna have pikemen out in a brief moment so he can deal with those.                                                                                                                                                      \\
    & Event Description     & Lidakor & And that is a huge amount of knights in there.                                                                                                                                                                                                                                                  \\
    & Event Description     & Lidakor & Now DauT sends the monks back,                                                                                                                                                                                                                                                                  \\
    & Event Description     & Lidakor & and there is still a decent amount of camels to deal with,                                                                                                                                                                                                                                      \\
    & Strategy Discussion   & Lidakor & so this is not as perfect as it looks like for DauT.                                                                                                                                                                                                                                            \\
    & State Update          & Lidakor & But he still has a nearly insurmountable lead over here.                                                                                                                                                                                                                                        \\
    & Strategy Discussion   & Lidakor & Again, not the most efficient fight here for DauT.                                                                                                                                                                                                                                              \\
    & Strategy Discussion   & Lidakor & Indeed, he's going to have to bail out of this one for now, but Villese is losing a lot of his army, and replenishing that is going to be super, super hard.                                                                                                                                    \\
    & State Update          & Lidakor & So even with all these idles for DauT, he still has a better eco than Villese does.                                                                                                                                                                                                             \\
    & Strategy Prediction   & Lidakor & And he could also convert a few camels here.                                                                                                                                                                                                                                                    \\
    & State Update          & Lidakor & No conversions just yet.                                                                                                                                                                                                                                                                        \\
    & Event Description     & Lidakor & There is one 4-HP camel that is being converted over here,                                                                                                                                                                                                                                      \\
    & Event Description     & Lidakor & and apparently, the knights from DauT are taking this fight,                                                                                                                                                                                                                                    \\
    & Event Description     & Lidakor & and this is an awful fight for DauT.                                                                                                                                                                                                                                                            \\
    & Strategy Discussion   & Lidakor & DauT can be very happy that he did so much damage to Villese, because the last few minutes were rather disappointing from DauT.                                                                                                                                                                 \\
    & Strategy Discussion   & Lidakor & He took some very, very disappointing engagements.                                                                                                                                                                                                                                              \\
    & Strategy Prediction   & Lidakor & Eventually, he's gonna have the pikemen, and I think his lead is big enough to just be able to make those big mistakes.                                                                                                                                                                         \\
    & State Update          & Lidakor & He's got 48 on wood with 1.5k wood in the bank.                                                                                                                                                                                                                                                 \\
    & Strategy Prediction   & Lidakor & He needs to take 10-15 villagers from wood and add them to food here.                                                                                                                                                                                                                           \\
    & Strategy Discussion   & Lidakor & So yeah, DauT isn't really having, um, the best game of his life here, but he did enough damage on Villese to be able to just, um, get by with some big mistakes that he's making.                                                                                                              \\
    & Strategy Discussion   & Lidakor & Like, the eco isn't great, and it appears that Villese is just gonna tap out.                                                                                                                                                                                                                   \\
    & Strategy Discussion   & Lidakor & As I said, DauT had too big of a lead here.                                                                                                                                                                                                                                                     \\
    & Strategy Prediction   & Lidakor & So if DauT just makes the pikemen here, he can defend against this one.                                                                                                                                                                                                                         \\
    & Strategy Prediction   & Lidakor & He is going to send a few more villagers to food and he can be into Imperial very, very soon.                                                                                                                                                                                                   \\
    & Strategy Discussion   & Lidakor & Not the best game from DauT, I think his previous game was better, but the defense was exceptional from DauT, I believe.                                                                                                                                                                        \\
    & Strategy Discussion   & Lidakor & That basically, pretty much won the game alone for him.                                                                                                                                                                                                                                         \\
    & Strategy Discussion   & Lidakor & So with that, that is game number two in the hands of DauT as well.                                                                                                                                                                                                                             \\
    & Historical Account    & Lidakor & And just as Villese 3-0'd TaToH, DauT could 3-0 Villese over here, and Villese could be the first player out in the quarterfinals in a rather short series.                                                                                                                                     \\
    & Historical Account    & Lidakor & This is just one hour in, and, uh, we already have two games done.
  \end{longtable}
} % end scriptsize

{\scriptsize
  % Use custom tight utterance column type Ut for last column (defined in preamble)
  % Slightly adjusted widths to reduce overfull hbox warnings (widen final column)
  % Use single-letter custom column type U (defined in preamble) instead of multi-char Ut

  \setcounter{AppRow}{0}
  \begin{longtable}{@{}>{\stepcounter{AppRow}\theAppRow}r|L{2.3cm}L{1.0cm}U{11.55cm}@{}}
    \caption[Co-cast coding Villese vs DauT G2]{Coding of MembTV and Chris co-casting TheViper vs TheMax -- King of the Desert 3 Group Stage, Round 2, Group D, Game 2}\label{tab:memb-chris-coding}\\
    \multicolumn{4}{@{}l}{\footnotesize Source: \url{https://www.youtube.com/watch?v=1P9KoseYvsc&t=1194s}}\\[\AfterSourceLinksVSpace]
    \toprule
    \multicolumn{1}{r}{\textbf{\#}} &
    \makecell[l]{\textbf{Type}} &
    \makecell[l]{\textbf{Speaker}} &
    \multicolumn{1}{l}{\makecell[l]{\textbf{Utterance}}} \\
    \midrule
    \endfirsthead
    \multicolumn{4}{@{}l}{\textit{Table \thetable\ (continued)}}\\[0.3em]
    \toprule
    \multicolumn{1}{r}{\textbf{\#}} &
    \makecell[l]{\textbf{Type}} &
    \makecell[l]{\textbf{Speaker}} &
    \multicolumn{1}{l}{\makecell[l]{\textbf{Utterance}}} \\
    \midrule
    \endhead

    \midrule
    \multicolumn{4}{r}{\small Continued on next page} \\
    \endfoot

    \bottomrule
    \endlastfoot
    & Off-Topic             & MembTV       & Mr. Chris, thank you for joining me.                                                                                                                                                                                                                             \\
    & Match Introduction          & MembTV       & As TheMax in blue and, in yellow, TheViper on the left.                                                                                                                                                                                                          \\
    & Match Introduction          & MembTV       & And we have Byzantines versus China.                                                                                                                                                                                                                             \\
    & Civilization Analysis & Chris        & Wow, yeah, that's a matchup I wasn't expecting.                                                                                                                                                                                                                  \\
    & Civilization Analysis & Chris        & Um, Byzantines this early are definitely an interesting pick.                                                                                                                                                                                                    \\
    & Civilization Analysis & Chris        & Um, the maps are very open, so if you can get a lot of trash units out, they are a strong civ.                                                                                                                                                                   \\
    & Civilization Analysis & Chris        & But they don't really have the early economy bonuses, so it's tough to say.                                                                                                                                                                                      \\
    & Civilization Analysis & Chris        & Um, TheMax with Chinese, I mean, that's a pretty common pick just because you have the extra villagers.                                                                                                                                                          \\
    & Map Analysis          & Chris        & Um, if you look at this map, it seems like TheMax at least has one back gold, but that main gold is pretty forward.                                                                                                                                              \\
    & Map Analysis          & Chris        & At least it's on flat this time so you can TC it.                                                                                                                                                                                                                \\
    & Map Analysis          & Chris        & And if you look at TheViper, yes, his main gold is pretty forward too, and it's on an awkward hill, so I think it's pretty similar for the maps here.                                                                                                            \\
    & Map Analysis          & Chris        & The only one thing I'll say about TheMax's map is that the wood is a little far away. The same with TheViper, I guess. They both have far away wood.                                                                                                             \\
    & Game Discussion      & MembTV       & The wood on this custom map by Chrazini, just to let you know, we do this on purpose, it's farther all the time.                                                                                                                                                 \\
    & Game Discussion      & MembTV       & You know, a little bit smaller, and, uh, the elephants or rhinos have to be on a side or at the back.                                                                                                                                                            \\
    & Game Discussion      & MembTV       & It's like a moose (?), you know? You are almost never going to see a woodline really close to the Town Center.                                                                                                                                                   \\
    & Game Discussion      & Chris        & And that's good. I mean, it looks like, uh, this one looks like a, I would call it a pretty fair map generation.                                                                                                                                                 \\
    & Game Discussion      & Chris        & Like, you've got a very equal {[}map{]}. Some of your resources are kind of forward, some of them are a little open.                                                                                                                                             \\
    & Map Analysis          & Chris        & It's not too easy for either player to fully wall; maybe they can do some small walls.                                                                                                                                                                           \\
    & Map Analysis          & Chris        & And I'm happy to see that, yeah, a game like this.                                                                                                                                                                                                               \\
    & Historical Account    & MembTV       & Well, do you remember back in the days?                                                                                                                                                                                                                          \\
    & Historical Account    & MembTV       & Those typical Arabias where you check the wood and when you're going to make the lumber camp... the oasis,                                                                                                                                                       \\
    & Historical Account    & Chris        & that's right!                                                                                                                                                                                                                                                    \\
    & Historical Account    & Chris        & You remember those times, huh?                                                                                                                                                                                                                                   \\
    & Historical Account    & Chris        & Of course, of course, you could never forget. I mean, you sometimes on occasion would win or lose games because of that, so I'm happy to see it's closer for sure.                                                                                               \\
    & Joke \footnote{This is historical account, a joke and an event description in one. (Describing that something works now, but he lost due to it last game, so describing this minimal achievement is mostly said as a joke)}     & Chris        & And TheMax had a successful zebra lure this time, which is good.                                                                                                                                                                                                 \\
    & State Update          & Chris        & Right off the first try. And, uh, looking at TheViper's base here, he's got four on wood.                                                                                                                                                                        \\
    & Event Description     & Chris        & Um, what's he doing here? Is that a house?                                                                                                                                                                                                                       \\
    & Strategy Discussion   & Chris        & That house... so, kind of the small walls, using the house for that.                                                                                                                                                                                             \\
    & Event Description     & Chris        & Um, is he going to build the mill yet? He is.                                                                                                                                                                                                                    \\
    & Strategy Prediction   & Chris        & So he's not going for the fast rush.                                                                                                                                                                                                                             \\
    & Event Description     & MembTV       & And if we look at TheMax here, same idea. He's starting to wall there with that one villager for houses.                                                                                                                                                         \\
    & Event Description     & MembTV       & If you check, Chris... it's killing, it's killing!  it's killing! Oh! He got it!                                                                                                                                                                                 \\
    & Historical Account    & MembTV       & I have to say, he's happy, man. Because he's always saving the villagers! I mean, he's not human, he's safe all the time, man, come on!                                                                                                                          \\
    & Event Description     & MembTV       & I mean, you have to lose some, and now he's laming two goats also.                                                                                                                                                                                               \\
    & Event Description     & MembTV       & Oh, this is good now for TheMax, really good.                                                                                                                                                                                                                    \\
    & Event Description     & Chris        & Yeah, that was very well played there.                                                                                                                                                                                                                           \\
    & Historical Account    & MembTV       & In the previous series, in the previous game, sorry... if you think everything went completely terrible for TheMax... pushing the zebras (?), deleting the mill (?), I mean, everything went wrong for the first minute.                                         \\
    & Historical Account    & MembTV       & Incredible. Now it seems to be different, which is good because we want a longer series for sure.                                                                                                                                                                \\
    & Historical Account    & MembTV       & I want to see great games because the previous one was unfortunate, to be honest, and we couldn't see a really potential huge battle, right?                                                                                                                     \\
    & Historical Account    & Chris        & Yeah, it was definitely a bit lopsided in the first game, and I think it's always tough to deal with Aztecs no matter what civ you are.                                                                                                                          \\
    & Historical Account    & Chris        & But in a situation like that, and playing against TheViper, I mean, those are things that make that even tougher, I would say.                                                                                                                                   \\
    & Historical Account    & Chris        & But it looks like TheMax is doing really well here, I would say.                                                                                                                                                                                                 \\
    & Strategy Prediction   & Chris        & And, uh, see what strategy are they doing here? No barracks yet.                                                                                                                                                                                                 \\
    & Strategy Prediction   & Chris        & I don't know, what do you think? I mean, I'm looking at Feudal Age.                                                                                                                                                                                              \\
    & Strategy Prediction   & Chris        & Yeah, I think they will go all-in in the Feudal Age.                                                                                                                                                                                                             \\
    & Civilization Analysis & MembTV       & Okay, overall, what civilization you prefer, if you're looking at Feudal, Castle, and late Imp: Byzantines or Chinese?                                                                                                                                           \\
    & Civilization Analysis & Chris        & So, it's one of those things where Chinese is a really good civ, but, uh, if you can get that early Feudal, you know, early Castle Age trash units out, like camels or pikemen or skirms or whatever it is...                                                    \\
    & Civilization Analysis & Chris        & Um, and if it's an open map, you can really do a lot of damage.                                                                                                                                                                                                  \\
    & Civilization Analysis & Chris        & So you might take that early edge, and you also can go to Imp earlier as Byzantines.                                                                                                                                                                             \\
    & Civilization Analysis & Chris        & So, you know, I'm going to say it's not that clear-cut. It's definitely not that clear-cut.                                                                                                                                                                      \\
    & Civilization Analysis & Chris        & If it was a 100\% closed map and you're booming to Imp or something, I would prefer Chinese for sure. Um, but on an open Arabia, you know, we'll see.                                                                                                            \\
    & Civilization Analysis & MembTV       & Yeah, Byzantines have so many possibilities. The camels are not the best, but they are super cheap as well. The skirms are just cheap completely too.                                                                                                            \\
    & Civilization Analysis & MembTV       & The transition to Imperial Age is really cheap also. So I don't know, I think people underestimate Byzantines.                                                                                                                                                   \\
    & Civilization Analysis & Chris        & Absolutely, they're very strong. And I think the other thing too is that we had a lot of closed Arabia maps for a while.                                                                                                                                         \\
    & Game Discussion      & Chris        & You know, from what I've heard about the ladder is that it was a much more closed map.                                                                                                                                                                           \\
    & Civilization Analysis & Chris        & And if you have a closed map, you don't necessarily want Byzantines because you're not doing mass spears, you know, pikes, or camels that often.                                                                                                                 \\
    & Civilization Analysis & Chris        & But in a map like this, if you're fighting your opponent and if it's open enough to do damage, those units are awesome, especially because they're so cheap.                                                                                                     \\
    & Strategy Prediction   & Chris        & And we can see, look here, TheMax is doing a men-at-arms rush.                                                                                                                                                                                                   \\
    & Strategy Prediction   & Chris        & And it looks like TheViper's just going straight to Feudal, kind of trying to wall the left side here.                                                                                                                                                           \\
    & Strategy Discussion   & Chris        & We'll see if that men-at-arms rush can get in and do some damage.                                                                                                                                                                                                \\
    & Event Description     & Both         & Okay, he's taking this and, uh, yeah, scout war there in the middle, yep.                                                                                                                                                                                        \\
    & Event Description     & MembTV       & And the damage he did with this scout is worthy for sure for TheMax. Gathering information, delaying a little bit TheViper.                                                                                                                                      \\
    & State Update          & MembTV       & And TheViper, well, TheViper didn't explore anything at all. He knows where the main gold is, but nothing else.                                                                                                                                                  \\
    & State Update          & MembTV       & And he doesn't know anything about plans from TheMax, and TheMax is with Chinese, with not really a common strategy: men-at-arms.                                                                                                                                \\
    & Strategy Discussion   & Chris        & Yeah, definitely. And I'm really surprised to see TheViper's economy here. That's a lot of food for a skirm build.                                                                                                                                               \\
    & Strategy Discussion   & Chris        & Um, definitely more than I expected.                                                                                                                                                                                                                             \\
    & Strategy Discussion   & Chris        & Usually you go up a little earlier and you have a little more wood.                                                                                                                                                                                              \\
    & Event Description     & Chris        & Um, but he is going to kill TheMax's scout. That's a good thing to do.                                                                                                                                                                                           \\
    & Event Description     & Chris        & Okay, I'm not saying it's not necessarily bad. It's, uh, I think he gathered the zebras and the ostriches.                                                                                                                                                       \\
    & Strategy Discussion   & Chris        & Um, which helps you get quicker to Castle Age, but you're probably not making any more units because...                                                                                                                                                          \\
    & Event Description     & Chris        & oh, he just spotted the men-at-arms!                                                                                                                                                                                                                             \\
    & Strategy Prediction   & Chris        & Uh, did he see it? Is he going to quick-wall his wood?                                                                                                                                                                                                           \\
    & Skill Discussion      & Chris        & We're going to see here in a second... is that he didn't notice that because he was focused on the scout war?                                                                                                                                                    \\
    & Skill Discussion      & Chris        & I think he did notice. No, yeah, he didn't.                                                                                                                                                                                                                      \\
    & Civilization Analysis & Chris        & But he's got Town Watch as Byzantines, so he might see this.                                                                                                                                                                                                     \\
    & Strategy Discussion   & MembTV       & Reaction... that's true. I think there's a hole.                                                                                                                                                                                                                 \\
    & Strategy Discussion   & MembTV       & Oh no, it's not a hole. He never has holes in this situation.                                                                                                                                                                                                    \\
    & Civilization Analysis & MembTV       & Town Watch for free, I mean, Byzantines is amazing in this kind of map, honestly. I really like them.                                                                                                                                                            \\
    & Strategy Discussion   & Chris        & The Town Watch definitely helped there.                                                                                                                                                                                                                          \\
    & Skill Discussion      & Chris        & I mean, it might have been tougher to react in time. He still might have, but, uh, definitely helped. Okay.                                                                                                                                                      \\
    & Strategy Discussion   & MembTV       & Well, going for the men-at-arms.                                                                                                                                                                                                                                 \\
    & Off-Topic & MembTV       & Thank you everyone for the support! Go, go, go, amigos, go, go, go!                                                                                                                                                                                              \\
    & Off-Topic & MembTV       & Let's make Age of Empires the best game ever! Well, I think it is already, you know.                                                                                                                                                                             \\
    & Off-Topic & MembTV       & They don't need to make it, but still, we need to make it even better, Mr. Chris.                                                                                                                                                                                \\
    & Off-Topic & MembTV       & With almost, well, on the way to 7,000 viewers, Chris.                                                                                                                                                                                                           \\
    & Off-Topic & MembTV       & So not bad, man, not bad. We're getting there.                                                                                                                                                                                                                   \\
    & Strategy Discussion   & MembTV       & But now, TheViper is not fully walled, but most likely... what would be the next step from TheMax to do the damage?                                                                                                                                              \\
    & State Update          & MembTV       & He needed an archery range, and he has one, wow.                                                                                                                                                                                                                 \\
    & State Update          & MembTV       & His economy is looking great.                                                                                                                                                                                                                                    \\
    & Civilization Analysis & Chris        & It is, and that's the extra Chinese villagers.                                                                                                                                                                                                                   \\
    & Strategy Prediction   & Chris        & I mean, and I really hope he's walling that bottom part.                                                                                                                                                                                                         \\
    & Strategy Prediction   & Chris        & And that's what I was about to say, finish the wall and also make sure you check your wood pile, TheMax.                                                                                                                                                         \\
    & Strategy Prediction   & Chris        & But if he goes up quicker and goes ranged units, that's a very powerful...                                                                                                                                                                                       \\
    & State Update          & Chris        & I mean, you see TheViper's economy, he had to build all those houses all over the place, which slowed down his farms.                                                                                                                                            \\
    & Event Description     & Chris        & And, uh, he's making skirmishers,                                                                                                                                                                                                                                \\
    & Strategy Prediction   & Chris        & so he's probably going to hit Castle Age later than TheMax will, I think.                                                                                                                                                                                        \\
    & State Update          & MembTV       & Because look, Chris, what TheMax is doing. He has one archer \footnote{Garrisoned, so TheViper knows that there is one garrisoned, but not how many},                                                                                                                                                                                                    \\
    & Strategy Discussion   & MembTV       & so in case TheViper is coming... mind games!                                                                                                                                                                                                                     \\
    & Strategy Discussion   & MembTV       & Thinking, "Okay, he has a lot of villagers inside," what... no, is that the purpose? I'm asking.                                                                                                                                                                 \\
    & Strategy Discussion   & Chris        & Well, if I were him, I'd probably make a couple more archers. But, uh, he might be doing a trick.                                                                                                                                                                \\
    & Strategy Discussion   & Chris        & So this could be a complete trick where you pretend you're going archers, TheViper sees that and goes skirms, and maybe you just build two stables and go mass knights.                                                                                          \\
    & Strategy Prediction   & Chris        & and just surprise the crap out of him. It could be, let's see.                                                                                                                                                                                                   \\
    & Historical Account    & MembTV       & Did you watch yesterday, did you watch yesterday, DauT versus Lierrey, uh, last game?                                                                                                                                                                            \\
    & Historical Account    & Chris        & I watched some of the games, not all of them, but yeah, I watched some of them.                                                                                                                                                                                  \\
    & Historical Account    & MembTV       & I mean, he did two stables, two stables inside behind the walls.                                                                                                                                                                                                 \\
    & Historical Account    & MembTV       & Full scouts with Bloodlines.                                                                                                                                                                                                                                     \\
    & Event Description     & MembTV       & Exciting! But look at the transition for TheMax to Castle Age, Chris, 18 minutes! Not bad!                                                                                                                                                                       \\
    & Strategy Discussion   & Chris        & Yeah, that's huge! That's what I mean, that was a very well-planned-out play.                                                                                                                                                                                    \\
    & Strategy Discussion   & Chris        & Men-at-arms keep them busy, and I liked how you ran around with them too. It's just... distract as long as possible and then just go to Castle Age.                                                                                                              \\
    & Event Description     & MembTV       & He's going forward with a villager. He's going to go forward to push with his siege.                                                                                                                                                                             \\
    & Event Description     & MembTV       & Is he? Where? Look on the right side.                                                                                                                                                                                                                            \\
    & Event Description     & MembTV       & Look at that villager. Oh, jeez, yes, with the villager here, man!                                                                                                                                                                                               \\
    & State Update          & Chris        & These are the stables behind, and the stables are behind, so TheViper won't see them!                                                                                                                                                                            \\
    & Strategy Discussion   & Chris        & Oh wow, I like that, that's a really smart strategy.                                                                                                                                                                                                             \\
    & State Update          & Chris        & TheViper's got nine skirmishers.                                                                                                                                                                                                                                 \\
    & Event Description     & MembTV       & Well, and he's mining gold now, TheViper.                                                                                                                                                                                                                        \\
    & Strategy Discussion   & MembTV       & Honestly, I mean, what can TheViper do in this situation?                                                                                                                                                                                                        \\
    & Strategy Discussion   & Chris        & Seriously, that's the thing. If TheViper doesn't see it, he could be in for a very big surprise when TheMax gets up way earlier.                                                                                                                                 \\
    & Strategy Discussion   & Chris        & And these nine skirmishers are not really worth anything because there are no archers around.                                                                                                                                                                    \\
    & Strategy Discussion   & Chris        & So, very sneaky strategy by TheMax here. I really like the foresight in his planning here.                                                                                                                                                                       \\
    & Strategy Prediction   & Chris        & Now I'm just hoping that villager doesn't go too close to get spotted.                                                                                                                                                                                           \\
    & Strategy Prediction   & Chris        & Just, if you hide that villager and hide the stable... I mean, there's only one stable, I think he needs a second stable, but, uh, if you hide that long enough... wait till TheViper attacks you with his skirms and just clean them up, and then go offensive. \\
    & Strategy Discussion   & Chris        & Uh, he could do really well here.                                                                                                                                                                                                                                \\
    & State Update          & MembTV       & Scale Barding Armor.                                                                                                                                                                                                                                             \\
    & Strategy Prediction   & MembTV       & And will you go now for a one-TC push? Because in this situation...                                                                                                                                                                                              \\
    & Strategy Prediction   & Chris        & oh yeah, if you're surprising someone like this, you go all out.                                                                                                                                                                                                 \\
    & Strategy Prediction   & Chris        & That's the plan, right? You don't leave any stone unturned.                                                                                                                                                                                                      \\
    & State Update          & MembTV       & Chris, check, check TheMax, check TheMax! He's still open at the back, be careful! Oh no.                                                                                                                                                                        \\
    & Strategy Discussion   & Chris        & He didn't connect the wall, yeah. So you always gotta check.                                                                                                                                                                                                     \\
    & Strategy Discussion   & Chris        & You always gotta check when you have a full wall. Try and send a villager outside. If it has a way out, it'll try and go there.                                                                                                                                  \\
    & Strategy Discussion   & Chris        & Oh, he's checking it right now. Look at that, TheMax! (What the hell,) TheMax is checking! Okay.                                                                                                                                                                 \\
    & Strategy Discussion   & Chris        & Good timing. I'm happy to see that.                                                                                                                                                                                                                              \\
    & Strategy Discussion   & Chris        & Good timing, yes, yeah, yeah. So you always want to do that, and most top players will do it. It's just you gotta have enough time to think of it.                                                                                                               \\
    & Strategy Discussion   & Chris        & So that's good. He did delete his wall here, I think he's just built a gate.                                                                                                                                                                                     \\
    & Event Description     & Chris        & Yeah... oh, he might have just seen the knight though, that's unfortunate.                                                                                                                                                                                       \\
    & Event Description     & MembTV       & Ah, but why did he delete it? Just make it... well, anyway, he was already doing a double stable, monastery siege workshop.                                                                                                                                      \\
    & Event Description     & MembTV       & And then my question is... TheViper is on the way to Castle Age. He was looking so far away!                                                                                                                                                                     \\
    & Strategy Discussion   & MembTV       & And look how he managed the economy! Because I'm not going to tell you... I mean, look at how he just caught up, and now a 21-minute Castle Age! Unbelievable!                                                                                                   \\
    & Strategy Discussion   & Chris        & Yeah, and I'm actually a little concerned for TheMax now. I was hoping for the surprise.                                                                                                                                                                         \\
    & State Update          & Chris        & But two stables going camels against one knight, one stable, one monk, you know, one monastery, and one siege workshop...                                                                                                                                        \\
    & Civilization Analysis & Chris        & Um, it's probably easy to defend with Byzantines because those units are so cheap.                                                                                                                                                                               \\
    & Strategy Prediction   & Chris        & So, I mean, it still might cause some damage. It might cause some headaches for TheViper here, and maybe TheMax can also, um, you know, do some damage and put some TCs behind it, who knows.                                                                    \\
    & Event Description     & Chris        & You see that first scorpion too. Scorpion, mangonel, knight up at the top there.                                                                                                                                                                                 \\
    & Event Description     & Chris        & And he's running around with that one knight, just being annoying. It's good.                                                                                                                                                                                    \\
    & Event Description     & MembTV       & But he's keeping all that army for one knight out of position.                                                                                                                                                                                                   \\
    & Event Description     & MembTV       & Now he's there with the siege.                                                                                                                                                                                                                                   \\
    & Event Description     & MembTV       & And like you said, he can still add TCs. He's doing a Town Center on the gold, Chinese economy.                                                                                                                                                                  \\
    & Strategy Discussion   & MembTV       & And I mean, if he's going to be able to defend, I don't know, but he should really be able to defend properly.                                                                                                                                                   \\
    & State Update          & MembTV       & He's got three crossbows just to deal with the spearman,                                                                                                                                                                                                         \\
    & Strategy Prediction   & MembTV       & and now the woodline is really, really, really exposed, so he has to move.                                                                                                                                                                                       \\
    & Event Description     & MembTV       & Making another lumber camp that is also exposed.                                                                                                                                                                                                                 \\
    & Strategy Discussion   & MembTV       & Still, this is a good position for TheMax.                                                                                                                                                                                                                       \\
    & Event Description     & MembTV       & Doing Bowsaw and gold mining upgrade, wow.                                                                                                                                                                                                                       \\
    & Skill Discussion      & MembTV       & It's looking like smooth gameplay.                                                                                                                                                                                                                               \\
    & Event Description     & Chris        & He's going to get in there, and that little hill might help quite a bit.                                                                                                                                                                                         \\
    & Event Description     & Chris        & Um, and see, TheViper is immediately making camels.                                                                                                                                                                                                              \\
    & Strategy Prediction   & Chris        & And that's the smart play right here because, you know, if you get six camels or ten camels, it's very easy to wipe that army.                                                                                                                                   \\
    & Strategy Prediction   & Chris        & But, um, it takes a while to build those too, so we'll see.                                                                                                                                                                                                      \\
    & Event Description     & Chris        & He does have monks with this army                                                                                                                                                                                                                                \\
    & Strategy Discussion   & Chris        & which is good                                                                                                                                                                                                                                                    \\
    & Event Description     & Chris        & and he did upgrade to crossbows just now as well.                                                                                                                                                                                                                \\
    & Strategy Discussion   & Chris        & A little surprised to see that.                                                                                                                                                                                                                                  \\
    & Event Description     & MembTV       & Look at TheViper, he's going forward.                                                                                                                                                                                                                            \\
    & Strategy Prediction   & MembTV       & He's going forward to make probably a TC on the gold.                                                                                                                                                                                                            \\
    & State Update          & MembTV       & He has the resources,                                                                                                                                                                                                                                            \\
    & Event Description     & MembTV       & but he's pushing the Town Center.                                                                                                                                                                                                                                \\
    & State Update          & MembTV       & He's got crossbows, he's got knights, siege, monks...                                                                                                                                                                                                            \\
    & Strategy Prediction   & MembTV       & well, if he plays with good micro, he's got the military to counter pretty much everything that TheViper has at the moment.                                                                                                                                      \\
    & Strategy Discussion   & Chris        & Absolutely. It's going to be a tough call here.                                                                                                                                                                                                                  \\
    & State Update          & Chris        & With these two stables, he's got garrisoned camels and garrisoned scouts.                                                                                                                                                                                        \\
    & State Update          & Chris        & And there are two monks there and a couple of crossbows and knights.                                                                                                                                                                                             \\
    & Strategy Discussion   & Chris        & Um, and very, very fortunately,                                                                                                                                                                                                                                  \\
    & Event Description     & Chris        & TheViper didn't get spotted with that front TC.                                                                                                                                                                                                                  \\
    & Strategy Prediction   & Chris        & Those two knights could have done some damage there for sure with just two spearmen,                                                                                                                                                                             \\
    & Event Description     & Chris        & but, um, he is pushing quite well here.                                                                                                                                                                                                                          \\
    & Strategy Discussion   & Chris        & Definitely causing some headaches for TheViper.                                                                                                                                                                                                                  \\
    & Event Description     & Chris        & TheMax adding the third TC at home.                                                                                                                                                                                                                              \\
    & State Update          & MembTV       & Okay, I wanted to ask you, how many villagers does TheViper have idle right now?                                                                                                                                                                                 \\
    & State Update          & MembTV       & Well, now six because he is with the Town Center, but for a while without any Town Center, without any idles.                                                                                                                                                    \\
    & Event Description     & MembTV       & And now be careful with the fight! Big fight coming! He's going to snipe the monks!                                                                                                                                                                              \\
    & Event Description     & MembTV       & TheViper is going to snipe the monks,                                                                                                                                                                                                                            \\
    & Strategy Discussion   & MembTV       & and that's terrible because the camels are there.                                                                                                                                                                                                                \\
    & Event Description     & MembTV       & The scorp is going to help. Look at the scout!                                                                                                                                                                                                                   \\
    & Skill Discussion      & MembTV       & The micro from TheViper is being on point.                                                                                                                                                                                                                       \\
    & Strategy Prediction   & MembTV       & And he probably is going to clean up everything, absolutely.                                                                                                                                                                                                     \\
    & Strategy Prediction   & Chris        & He will, absolutely.                                                                                                                                                                                                                                             \\
    & Strategy Discussion   & Chris        & It's, uh, those camels, those cheap camels are just a great counter for anything you come out with like knights and siege and that type of stuff.                                                                                                                \\
    & Strategy Discussion   & Chris        & So TheMax is doing the right thing here by retreating because you want to save as many units as you can.                                                                                                                                                         \\
    & Strategy Discussion   & Chris        & But, uh, it's going to be tough here.                                                                                                                                                                                                                            \\
    & State Update          & Chris        & So all of a sudden, TheMax has a bit better economy but a bit less military.                                                                                                                                                                                     \\
    & Strategy Discussion   & Chris        & So we'll see how it goes.                                                                                                                                                                                                                                        \\
    & State Update          & MembTV       & Yeah, and the monastery and the siege workshop... bye-bye.                                                                                                                                                                                                       \\
    & Event Description     & MembTV       & He's researching Redemption and Husbandry. TheViper is even getting Redemption                                                                                                                                                                                   \\
    & Strategy Discussion   & MembTV       & in case TheMax comes with some siege. What the hell!                                                                                                                                                                                                             \\
    & Event Description     & MembTV       & And now the third Town Center is nice, going home.                                                                                                                                                                                                               \\
    & Event Description     & MembTV       & But TheMax is instantly building two extra barracks.                                                                                                                                                                                                             \\
    & Strategy Prediction   & MembTV       & He's going to go full pikes?                                                                                                                                                                                                                                     \\
    & State Update          & Chris        & I don't know, but it's funny, you see that siege workshop from TheViper on the left of TheMax's base there?                                                                                                                                                      \\
    & Strategy Discussion   & Chris        & That's going to be hard for TheMax to deal with because it's right on the hill next to his main gold, um, and his secondary gold and secondary stone.                                                                                                            \\
    & Strategy Discussion   & Chris        & So there are tons of resources here for TheMax that are at play.                                                                                                                                                                                                 \\
    & State Update          & Chris        & And, uh, he doesn't see it yet either.                                                                                                                                                                                                                           \\
    & Game Discussion      & MembTV       & For a long time, Chinese were great to defend this because you could get Redemption as well. Now Chinese don't have Redemption anymore.                                                                                                                          \\
    & Event Description     & MembTV       & So he's going to go pikes.                                                                                                                                                                                                                                       \\
    & Strategy Prediction   & MembTV       & He will need a siege workshop at home on his own, right?                                                                                                                                                                                                         \\
    & State Update          & Chris        & Yes, absolutely. It looks like he's adding some crossbows.                                                                                                                                                                                                       \\
    & State Update          & Chris        & He's getting Husbandry,                                                                                                                                                                                                                                          \\
    & State Update          & Chris        & three barracks, and a second archery range.                                                                                                                                                                                                                      \\
    & Strategy Discussion   & Chris        & That's a very big mix of units, eh?                                                                                                                                                                                                                              \\
    & Strategy Prediction   & Chris        & Um, archers, probably pikes with these three barracks.                                                                                                                                                                                                           \\
    & Strategy Discussion   & MembTV       & But this boom is good,                                                                                                                                                                                                                                           \\
    & State Update          & MembTV       & he's still 10 villagers ahead                                                                                                                                                                                                                                    \\
    & Civilization Analysis & MembTV       & with the Chinese economy,                                                                                                                                                                                                                                        \\
    & Strategy Discussion   & MembTV       & which is solid.                                                                                                                                                                                                                                                  \\
    & State Update          & MembTV       & But still, uh, he's not mining stone.                                                                                                                                                                                                                            \\
    & Strategy Prediction   & MembTV       & A castle would help here a lot.                                                                                                                                                                                                                                  \\
    & Civilization Analysis & MembTV       & And we should not forget, we should not forget that the Byzantines can go up to Imp super cheap, and his units are super cheap as well.                                                                                                                          \\
    & Strategy Discussion   & MembTV       & So be careful with this push.                                                                                                                                                                                                                                    \\
    & Strategy Discussion   & Chris        & Yeah, I'm a bit worried about this push, to be honest.                                                                                                                                                                                                           \\
    & Event Description     & Chris        & It's on a hill, um, on top of that, there's no access to stone for TheMax inside his walls.                                                                                                                                                                      \\
    & State Update          & Chris        & And, uh, there's not a lot of wood left, to be honest.                                                                                                                                                                                                           \\
    & Strategy Discussion   & Chris        & Like that wood pile next to his gold was the main one; the other one's kind of out in the open.                                                                                                                                                                  \\
    & Event Description     & Chris        & Um, he's building his own siege workshop there.                                                                                                                                                                                                                  \\
    & Strategy Discussion   & Chris        & He doesn't want to lose his main gold,                                                                                                                                                                                                                           \\
    & Strategy Prediction   & Chris        & and I don't know if he really has a choice to keep his main gold here.                                                                                                                                                                                           \\
    & Strategy Discussion      & Chris        & It's going to be tough to hold, I think.                                                                                                                                                                                                                         \\
    & State Update          & MembTV       & Yeah, look at the amount of camels he has, Chris!                                                                                                                                                                                                                \\
    & State Update          & MembTV       & Just a lot of camels, but he has crossbows now with +1.                                                                                                                                                                                                          \\
    & Strategy Prediction   & MembTV       & And he's going to start to spawn pikes probably now.                                                                                                                                                                                                             \\
    & Strategy Prediction   & MembTV       & He will need sick micro with a future mangonel.                                                                                                                                                                                                                  \\
    & Strategy Discussion   & MembTV       & For now, he doesn't have anything because he doesn't have a lot of gold.                                                                                                                                                                                         \\
    & State Update          & MembTV       & Luckily, he has gold at the back.                                                                                                                                                                                                                                \\
    & Event Description     & MembTV       & Oh, he's going for a counter-attack with the knights. It seems so, looks like it, yeah, I think so.                                                                                                                                                              \\
    & Event Description     & Chris        & But that TC is going down. It's losing a lot of HP. He's repairing it, very expensive. Uh, he's got a few crossbows here.                                                                                                                                        \\
    & Strategy Discussion   & Chris        & Kind of in trouble.                                                                                                                                                                                                                                              \\
    & Strategy Discussion   & MembTV       & Okay, this is nice.                                                                                                                                                                                                                                              \\
    & Event Description     & Chris        & More economy, though.                                                                                                                                                                                                                                            \\
    & Event Description     & MembTV       & Oh, four Town Centers!                                                                                                                                                                                                                                           \\
    & Strategy Discussion   & MembTV       & I mean, how is it possible after this push he gets back now?                                                                                                                                                                                                     \\
    & Recap                 & MembTV       & Chain Barding Armor, siege workshop, monastery, and three mangonels pushing, and he's still able to make four Town Centers, Chris!                                                                                                                               \\
    & Strategy Discussion      & Chris        & Yeah, it looks like right now he's focused on his economy.                                                                                                                                                                                                       \\
    & State Update          & Chris        & He's not building any new military at the moment.                                                                                                                                                                                                                \\
    & Strategy Discussion   & Chris        & Maybe one new monk, but, uh, he's just focusing on booming with all these four TCs, which is a Viper classic.                                                                                                                                                    \\
    & Strategy Discussion   & Chris        & You know, you make enough army to cause your enemy some heartache with mangonels and monks, and then you boom behind it.                                                                                                                                         \\
    & Strategy Discussion   & Chris        & So, definitely a useful strategy.                                                                                                                                                                                                                                \\
    & Strategy Prediction   & MembTV       & He needs economy because look at the amount of farming he has. I mean, he needs more army!                                                                                                                                                                       \\
    & Event Description     & MembTV       & He has a lot of army already. He's pushing you with the camels.                                                                                                                                                                                                  \\
    & Strategy Discussion   & MembTV       & It's a crazy situation here.                                                                                                                                                                                                                                     \\
    & Strategy Discussion   & MembTV       & I don't know what TheMax can do in this situation. It's incredible how TheViper has changed the game completely.                                                                                                                                                 \\
    & Strategy Discussion   & Chris        & Yeah, that's a very aggressive push in a very sensitive spot for TheMax.                                                                                                                                                                                         \\
    & Strategy Discussion   & Chris        & Um, I think TheMax is going to have a tough time. This is because of that hill, too, right?                                                                                                                                                                      \\
    & State Update          & Chris        & And, uh, those monks got Redemption,                                                                                                                                                                                                                             \\
    & Strategy Discussion   & Chris        & so it's kind of like, whatever you build, you're in a tight spot.                                                                                                                                                                                                \\
    & Event Description     & Chris        & So TheMax is trying to go for the stone on the right there,                                                                                                                                                                                                      \\
    & Strategy Discussion   & Chris        & and I think if he had had that earlier, like if it was somewhere in the back of his base, um, he would have had a chance to maybe build a castle by now.                                                                                                         \\
    & Strategy Discussion   & Chris        & But with where it is, it's not really been an option for him.                                                                                                                                                                                                    \\
    & Strategy Prediction   & MembTV       & Yeah, but you know what time it is? What I smell is a Viper castle on his face.                                                                                                                                                                                  \\
    & Strategy Discussion   & MembTV       & And if he's building a castle there, what can TheMax do? It's going to be terrible.                                                                                                                                                                              \\
    & Event Description     & MembTV       & He's coming with a lot of villagers.                                                                                                                                                                                                                             \\
    & State Update          & MembTV       & He's got a good amount of pikes, yes. He's also got two mangonels                                                                                                                                                                                                \\
    & Strategy Discussion   & MembTV       & , but a castle there is deadly.                                                                                                                                                                                                                                  \\
    & Strategy Discussion   & MembTV       & This can be definitive, man,                                                                                                                                                                                                                                     \\
    & Strategy Discussion   & Chris        & I agree. That's a brutal castle because that covers his main gold, your secondary gold, and your secondary stone, all three at once.                                                                                                                             \\
    & Strategy Discussion   & Chris        & And you're on a hill that... he tried to steal the mangonel there.                                                                                                                                                                                               \\
    & State Update          & Chris        & So I mean, I know the score is close and the vills are close,                                                                                                                                                                                                    \\
    & Strategy Discussion   & Chris        & but, uh, that's a very, very difficult-to-deal-with castle.                                                                                                                                                                                                      \\
    & Strategy Discussion   & MembTV       & Really complicated.                                                                                                                                                                                                                                              \\
    & State Update          & MembTV       & It's also true that TheMax, for now, has gold at the back.                                                                                                                                                                                                       \\
    & State Update          & MembTV       & He has stone at the back as well, and he's taking it.                                                                                                                                                                                                            \\
    & State Update          & MembTV       & He still has good population overall, almost the same if you see it.                                                                                                                                                                                             \\
    & State Update          & MembTV       & Vills are one ahead now. He lost the Town Center. The resources are looking good as well.                                                                                                                                                                        \\
    & State Update          & MembTV       & But the map control is complete domination right now for TheViper.                                                                                                                                                                                               \\
    & Strategy Discussion   & Chris        & Yep, and that's the tough thing about camels, right? Camels and monks.                                                                                                                                                                                           \\
    & Strategy Discussion   & Chris        & Like, you know, if you're just making knights and that type of stuff, you can get chased down by camels.                                                                                                                                                         \\
    & Strategy Discussion   & Chris        & And even a small amount of crossbows like this, they don't really stand a chance against, you know, 10 cheap camels from Byzantines.                                                                                                                             \\
    & Event Description     & MembTV       & Wow. Look on the right side, he's not walling that. Okay, he walled it now.                                                                                                                                                                                      \\
    & Event Description     & MembTV       & And he's building the Town Center.                                                                                                                                                                                                                               \\
    & Event Description     & MembTV       & He needs to wall more because he's getting inside.                                                                                                                                                                                                               \\
    & Strategy Discussion   & MembTV       & He's sending out the pikes, but the Town Center is up so he's okay in the north.                                                                                                                                                                                 \\
    & Event Description     & MembTV       & Look what, uh, TheMax is doing. I mean, he's trying to go and push with rams and probably, I don't know, light cavalry?                                                                                                                                          \\
    & Event Description     & MembTV       & He's building a watch tower,                                                                                                                                                                                                                                     \\
    & Event Description     & MembTV       & but TheViper is on the way to Imperial Age.                                                                                                                                                                                                                      \\
    & Event Description     & Chris        & Yeah, and TheViper just scouted that right wood pile and how open it is.                                                                                                                                                                                         \\
    & State Update          & Chris        & Like, there are 50 vills there on wood for TheMax because he doesn't have any other wood piles.                                                                                                                                                                  \\
    & Event Description     & Chris        & So if he goes over there with those camels, which he's trying to do with those six camels,                                                                                                                                                                       \\
    & Strategy Prediction   & Chris        & um, that would be very, very sketchy for TheMax.                                                                                                                                                                                                                 \\
    & Event Description     & Both         & Right now, he's raiding there. He's raiding there! Oh, it's a brutal spot and no army.                                                                                                                                                                           \\
    & Strategy Discussion   & MembTV       & But that spot was very easy to wall. It wasn't like a really big gap. He doesn't have anything there.                                                                                                                                                            \\
    & State Update          & MembTV       & Population is still... well, now TheViper is getting ahead.                                                                                                                                                                                                      \\
    & Event Description     & MembTV       & Another TC.                                                                                                                                                                                                                                                      \\
    & Strategy Discussion   & MembTV       & It's, uh, it's impressive how he just changed the game.                                                                                                                                                                                                          \\
    & Recap                 & MembTV       & Waiting a little bit and, uh, it was one attack and he came back completely in the game when it was looking good for TheMax.                                                                                                                                     \\
    & Strategy Discussion      & Chris        & He did, absolutely. Those camels were just a godsend at that point in time.                                                                                                                                                                                      \\
    & Recap                 & Chris        & Um, he had enough resources for those ten camels, those two stables, and that definitely made it easier to defend.                                                                                                                                               \\
    & Strategy Discussion      & Chris        & I was going to say that gets really tough for him right now because he's kind of stuck in Castle Age with his current resources, and TheViper's almost in Imp.                                                                                                   \\
    & Strategy Prediction   & Chris        & So, I mean, even if he can kill a TC or two here, um, he's going to see trebuchets on his main TC very soon.                                                                                                                                                     \\
    & Strategy Prediction   & Chris        & And Byzantines, you know, an early Imp gets really good with, um, Chemistry and these archers that are coming out now from TheViper.                                                                                                                             \\
    & Strategy Prediction   & Chris        & And it looks like he's almost able... I think he'll be able to defend the TC and repair it to survive.                                                                                                                                                           \\
    & Strategy Discussion   & MembTV       & Yeah, he's crazy, man, it's crazy.                                                                                                                                                                                                                               \\
    & Event Description     & MembTV       & Yeah, he killed the ram, but the thing is                                                                                                                                                                                                                        \\
    & State Update          & MembTV       & , look at the resources for TheViper.                                                                                                                                                                                                                            \\
    & Event Description     & MembTV       & Look at his economy, look at the economy! He's getting all the upgrades he wants.                                                                                                                                                                                \\
    & Strategy Discussion   & MembTV       & It's just crazy.                                                                                                                                                                                                                                                 \\
    & Event Description     & MembTV       & Another castle on his face.                                                                                                                                                                                                                                      \\
    & Strategy Discussion   & MembTV       & This is brutal domination in this game by TheViper with Byzantines versus Chinese.                                                                                                                                                                               \\
    & Strategy Discussion   & MembTV       & As simple as that. And TheMax was playing really, really nice, but TheViper stopped his push, let's say, quite easily.                                                                                                                                           \\
    & Strategy Discussion   & Chris        & Wow, yeah, I think that push from TheMax was tough. He almost needed to commit more and build two stables and more monks.                                                                                                                                        \\
    & Strategy Prediction   & Chris        & And, you know, maybe keep the same amount of siege and get in and try and go for TheViper's gold and try and keep that main gold away from him.                                                                                                                  \\
    & Strategy Discussion   & Chris        & But, uh, without that, it kind of snowballed in the other direction.                                                                                                                                                                                             \\
    & Technical Issues      & MembTV       & I don't know if he did resign or not, but the game for me says... end of game data. It says to me...                                                                                                                                                             \\
    & Technical Issues      & Chris        & Yeah, TheMax resigned. He resigned, right? Okay.                                                                                                                                                                                                                 \\
    & Event Description     & MembTV       & The game is over. That was the case. He resigned. Two - Zero man.
  \end{longtable}
}
\section{Supplementary Tables}
\label{appendix:domain-analysis-supplementary-tables}

Table~\ref{tab:captureage-states} (sub-type descriptions and PEAS mapping) appears in Chapter~\ref{chap:prelim} because it serves as a core scaffold for the coding scheme.

\section{User Study Operationalization Notes}
\label{appendix:prelim-operationalization}

The following notes summarize how we operationalized event and content sub-types when selecting questionnaire items.

We distinguish between the following types of Battles:
\begin{itemize}
  \item \textbf{Eco Delay:} Battles that aim to delay an opponent's economy or building production, often early game. These don't usually lead to deaths, but can dictate where battles take place.
  \item \textbf{Raids:} Battles that lead to villager deaths, often early game. This can also happen late game, aiming at killing villagers.
  \item \textbf{Regular Battles:} Battles that lead to unit deaths.
  \item \textbf{Breaking through Defenses:} Battles where one player tries to break through the opponent's walls or other defenses, often leading to a decisive victory or great economic damage.
  \item \textbf{Overwhelming Opponent Base:} Battles where one player floods the opponent's base with units, often leading to a decisive victory or great economic damage. This aims to further push an advantage, rather than causing economic damage.
\end{itemize}

For Production we distinguish between:
\begin{itemize}
  \item \textbf{Military Production:} Production of military units, often indicating an upcoming battle.
  \item \textbf{Economic Production:} Production of economic units, often indicating a focus on economy.
  \item \textbf{Forward Construction:} Construction of forward military buildings and fortifications, aiming to establish a strategic position closer to the opponent or deny access to critical resources.
  \item \textbf{Defensive Construction:} Construction of military buildings and fortifications at home, often indicating a shift in strategy.
  \item \textbf{Economy Construction:} Construction of buildings that focus on economy, such as Town Centers or Mills, often indicating an economic boom strategy.
\end{itemize}
We leave out technologies from events since they are more important as part of state updates.

In our coding we didn't distinguish between different types of scouting, but based on our experience we distinguish between:
\begin{itemize}
  \item \textbf{Scout movements:} Movements of scout units to gather information about the opponent's base and army composition and movements.
  \item \textbf{Exploration:} Exploring the map to find resources, enemy positions, and potential expansion sites, this is permanently revealed in the Fog of War.
  \item \textbf{Player gaze:} The areas of the map that players are focusing on, as indicated by their camera movements and unit selections. While technically the player can see something either in the viewport or minimap, it doesn't mean that they are actually paying attention to it.
\end{itemize}

Finally for Moving we distinguish between:
\begin{itemize}
  \item \textbf{Army Movements:} Movements of military units, often indicating an upcoming battle or strategic positioning.
  \item \textbf{Economy Movements:} Movements of economic units, often indicating resource gathering or expansion, or a forward building.
  \item \textbf{Monk Movements:} Movements of monk units, often indicating relic collection or strategic positioning.
\end{itemize}

For additional context, such as state updates, we identify the following aspects:
\begin{itemize}
  \item \textbf{Economy distribution:} Per-player allocation of villagers across tasks (including idle Villagers)
  \item \textbf{Army composition:} The counts and types of military units within each player's army.
  \item \textbf{Economy balance:} The total count of economic units (and idles) compared between players.
  \item \textbf{Army balance:} The total count of military units compared between players.
  \item \textbf{Economic Technologies:} Technologies that enhance economic capabilities.
  \item \textbf{Military Technologies:} Technologies that enhance military capabilities.
  \item \textbf{Age-up Timings:} The timings of advancing to the next age.
\end{itemize}

For high-level understanding of the narrative content we select three distinct categories. Although we didn't explicitly use these categories in our coding, we find them useful for understanding the importance of different content types. We distinguish between:
\begin{itemize}
  \item \textbf{Action-heavy sections:} Sections that focus on battles and other high-intensity events.
  \item \textbf{Strategic sections:} Sections that focus on strategic decisions and analysis, such as technology choices and uncommon plays.
  \item \textbf{Skill-based sections:} Sections that focus on explaining player skill and mastery, such as specific unit control (micromanagement).
\end{itemize}

\section{Preliminary Study Questionnaire}
\label{appendix:prelim-questionnaire}

This appendix lists the full wording of the questionnaire items used in the preliminary user study.

\subsection{Commentary Content Items (C-set)}
Items were answered on a 5-point Likert scale: Not Relevant, Slightly Relevant, Relevant, Very Relevant, Required.

\textit{For each of the following described events please indicate their relevance for you to understand the course of the game\footnote{"Game" refers to match among viewers of Age of Empires II.} from the video. If you don't know what a term means, simply select "I don't know". If you think the relevance of an event is dependent on other factors, assume it to be the most common case. Optionally, you can rewatch sections of the video below at your convenience.}
\begin{enumerate}[label={C\arabic*}]
  \item \textbf{IntroLayout} Introduction of map layout - including resource locations and possible wall locations
  \item \textbf{IntroPlayers} Introduction of players and related trivia
  \item \textbf{IntroCivs} Introduction of civilizations and their strengths
  \item \textbf{IntroEconomy} Early economy decisions (e.g. pushing deer)
  \item \textbf{BattleDelay} Aggressive maneuvers: Not leading to kills, but delaying a player
  \item \textbf{BattleMilitary} Aggressive maneuvers: Leading to kills or losses on military units
  \item \textbf{BattleEcoKills} Aggressive maneuvers: Leading to kills on villagers and other economy units
  \item \textbf{BattleBreakIn} Breaking through the wall, over-chops and possible defenses.
  \item \textbf{BattleOverwhelm} Overwhelming the opponent's base with a large army
  \item \textbf{EventScouting} Scouting information
  \item \textbf{EventUncommon} Uncommon events like: Daut Castles and Siege Tower hops
  \item \textbf{StateAgeUp} Technology Status: Age-up timings
  \item \textbf{StateMilTech} Technology Status: Improving unit strength leading up to a fight, e.g. "+4 knights"
  \item \textbf{Predictions} Predictions (e.g. Villese will resign when Imperial Age is reached)
  \item \textbf{Knowledge} Specific information to the current situation (e.g. You don't want to lose your berries as Franks)
\end{enumerate}

\subsection{General Preference Items (G-set)}
Items were answered on a 5-point Likert scale: "Strongly Disagree", "Disagree", "Neutral", "Agree", "Strongly Agree". "Does not apply" was added as an option to indicate they do not watch other content; when selected the question was ignored in the analysis.

\begin{enumerate}[label={G\arabic*}]
  \item \textbf{Understanding} Watching the video summary gave me a good idea how the game played out.
  \item \textbf{SummaryOpinion} In general, I think game summaries can help me consume competitive AoE2 content more effectively.
  \item \textbf{PreferFullGames} Generally, I prefer watching full games, even if that means that some of the spoken content is not relevant to the game.
  \item \textbf{LimitedTime} I only watch a few competitive games from a tournament because I don't have time to watch all the games I would want to watch.
  \item \textbf{InterestBattles} When watching games I'm generally interested in the large battles.
  \item \textbf{InterestStrategy} When watching games I'm generally interested in the strategic choices made by each player.
  \item \textbf{InterestSkill} When watching games I am generally interested in watching the skill level of the players.
  \item \textbf{TechnicalIssues} It was hard to ignore video cutting issues and casting style as they negatively impacted my watching experience.
\end{enumerate}

\subsection{Other and Open-ended Items (O-set)}
\begin{enumerate}[label={O\arabic*}]
  \item \textbf{CompetenceSelf} How competent do you think you are in understanding Age of Empires II: Definitive Edition competitive 1v1 Random Map games?  \textit{(Not Competent, Slightly Competent, Competent, Very Competent, Expert)}
  \item \textbf{CompetenceELO} What is your current ELO\footnote{Elo is a rating system use in the game to calculate the relative skill levels of players in games, to match opponents of similar strength} on the 1v1 ladder for AoE2: Definitive Edition?  \textit{(800 or less, 801-1000, 1001-1200, 1201-1500, 1501-1700, 1700 or higher, I don't have a rating)}
  \item \textbf{WatchFrequency} How often do you watch competitive AoE2 videos or live streams?  \textit{(Daily, Weekly, Monthly, Never)}
  \item \textbf{WatchPreviously} Did you see this game before? This can be on any channel or in spectator mode. \textit{(Yes/No, Unsure)}
\end{enumerate}

\begin{enumerate}[label={O5}]
  \item \textbf{Feedback} Did you feel there was anything missing to follow "the story" of this match? If so, please indicate what you were missing below.
\end{enumerate}

\begin{table}[htbp]
  \caption{Frequency of Event Description sub-types reported in the sample commentary, repeated tags for the same event have been deduplicated. State information that is part of an event description (e.g. contrasting economic strength) is excluded.}
  \centering
  \setlength{\tabcolsep}{8pt}
  \renewcommand{\arraystretch}{1.3}
  \begin{tabularx}{0.6\textwidth}{@{}l>{\centering\arraybackslash}X@{}}
    \toprule
    \textbf{Event} & \textbf{Count} \\
    \midrule
    Battle (Incl. Eco Delay \& Raids)              & 16 \\
    Building Production & 9  \\
    Scouting            & 4  \\
    Movement            & 3  \\
    Technology          & 2  \\
    Economy             & 2  \\
    Laming              & 1  \\
    Age-up              & 1  \\
    Unit Production     & 1  \\
    \bottomrule
  \end{tabularx}
  \label{tab:event-frequencies}
\end{table}

\begin{table}[htbp]
  \caption{Frequency of Information sources used by T90 \& Nili in their commentary.}
  \centering
  \setlength{\tabcolsep}{8pt}
  \renewcommand{\arraystretch}{1.3}
  \begin{tabularx}{0.6\textwidth}{@{}l>{\centering\arraybackslash}X@{}}
    \toprule
    \textbf{Event} & \textbf{Count} \\
    \midrule
    Descriptive         & 92 \\
    Domain              & 55  \\
    Analytical          & 46  \\
    Background          & 21  \\
    Experiential        & 19  \\
    \bottomrule
  \end{tabularx}
  \label{tab:info-sources-frequencies}
\end{table}

\begin{table}[htbp]
  \caption{Co-occurrence of Categories and Information Sources used by T90 \& Nili. Values indicate the proportion of each Information source within each Commentary Category. The rightmost column indicates the total count of utterances for each Commentary Category.}
  \label{tab:co-occurrence-type-source}
  \begin{tabular}{lrrrrrrrr}
    & Descr. & Domain & Analy. & Backgr. & Exp. & \textless{}untagged\textgreater{} & Count &  \\
    Event Description                 & \textbf{78\%}                                                                           & 10\%                                                                      & 8\%                                                                           & 0\%                                                                           & 0\%                                                                             & 3\%                                                                                                  & 86                                                                       \\
    Strategy Discussion               & 5\%                                                                            & \textbf{47\%}                                                                      & \textbf{39\%}                                                                          & 5\%                                                                           & 0\%                                                                             & 3\%                                                                                                  & 38                                                                       \\
    Strategy Prediction               & 12\%                                                                           & \textbf{24\%}                                                                      & \textbf{53\%}                                                                          & 3\%                                                                           & 0\%                                                                             & 9\%                                                                                                  & 34                                                                       \\
    Game Discussion                   & 0\%                                                                            & \textbf{21\%}                                                                      & 13\%                                                                          & \textbf{17\%}                                                                          & \textbf{17\%}                                                                            & \textbf{33\%}                                                                                                 & 24                                                                       \\
    Historical Account                & 0\%                                                                            & 0\%                                                                       & 0\%                                                                           & \textbf{50\%}                                                                          & \textbf{25\%}                                                                            & \textbf{25\%}                                                                                                 & 16                                                                       \\
    Off-Topic                         & 0\%                                                                            & 0\%                                                                       & 0\%                                                                           & 0\%                                                                           & \textbf{38\%}                                                                            & \textbf{63\%}                                                                                                 & 16                                                                       \\
    Civilization Analysis             & \textbf{21\%}                                                                           & \textbf{43\%}                                                                      & 7\%                                                                           & 14\%                                                                          & 0\%                                                                             & 14\%                                                                                                 & 14                                                                       \\
    State Update                      & \textbf{83\%}                                                                           & 17\%                                                                      & 0\%                                                                           & 0\%                                                                           & 0\%                                                                             & 0\%                                                                                                  & 12                                                                       \\
    Player Analysis                   & 0\%                                                                            & 10\%                                                                      & 0\%                                                                           & \textbf{30\%}                                                                          & \textbf{50\%}                                                                            & 10\%                                                                                                 & 10                                                                       \\
    Technical Issues                  & 17\%                                                                           & \textbf{33\%}                                                                      & 0\%                                                                           & 0\%                                                                           & 0\%                                                                             & \textbf{50\%}                                                                                                 & 6                                                                        \\
    Joke                              & 0\%                                                                            & \textbf{33\%}                                                                      & 0\%                                                                           & 17\%                                                                          & 0\%                                                                             & \textbf{50\%}                                                                                                 & 6                                                                        \\
    Recap                             & \textbf{80\%}                                                                           & 0\%                                                                       & \textbf{20\%}                                                                          & 0\%                                                                           & 0\%                                                                             & 0\%                                                                                                  & 5                                                                        \\
    Skill Discussion                  & 0\%                                                                            & \textbf{67\%}                                                                      & \textbf{33\%}                                                                          & 0\%                                                                           & 0\%                                                                             & 0\%                                                                                                  & 3                                                                        \\
    Match Introduction                & \textbf{100\%}                                                                           & 0\%                                                                       & 0\%                                                                           & 0\%                                                                           & 0\%                                                                             & 0\%                                                                                                  & 1                                                                        \\
    Map Analysis                      & \textbf{100\%}                                                                           & 0\%                                                                       & 0\%                                                                           & 0\%                                                                           & 0\%                                                                             & 0\%                                                                                                  & 1                                                                        \\
    \textless{}untagged\textgreater{} & 0\%                                                                            & 0\%                                                                       & 0\%                                                                           & 0\%                                                                           & 0\%                                                                             & \textbf{100\%}                                                                                                & 2                                                                        \\
  \end{tabular}
\end{table}

\FloatBarrier

\chapter{Prototype}
\label{appendix:prototype}

\section{Data Schema}
\label{appendix:data-schema}
Each recorded game is exported as a set of Apache Arrow files, one per data type, along with a descriptor file that indexes the files and provides metadata.

\subsection{Game Lobby Settings}
\label{appendix:lobby-settings}
The recorded games in our KotD IV dataset were played with the following lobby configuration:
\begin{itemize}[noitemsep,topsep=0pt]
  \item \textbf{Data Mod:} Definitive Set
  \item \textbf{Game Mode:} Random Map
  \item \textbf{Location:} Arabia
  \item \textbf{Map Size:} Tiny
  \item \textbf{AI Difficulty:} Standard
  \item \textbf{Resources:} Standard
  \item \textbf{Population:} 200
  \item \textbf{Game Speed:} Normal
  \item \textbf{Reveal Map:} Normal
  \item \textbf{Starting Age:} Standard
  \item \textbf{Ending Age:} Standard
  \item \textbf{Treaty Length:} None
  \item \textbf{Victory:} Conquest
  \item \textbf{Lock Teams:} On
  \item \textbf{Team Together:} On
  \item \textbf{Team Positions:} Off
  \item \textbf{Shared Exploration:} Off
  \item \textbf{Lock Speed:} On
  \item \textbf{Allow Cheats:} Off
  \item \textbf{Turbo Mode:} Off
  \item \textbf{Full Tech Tree:} Off
  \item \textbf{Record Game:} On
\end{itemize}

\subsection{Descriptor}
The export process generates a \texttt{descriptor.json} file that serves as an index for the generated Arrow files. It categorizes the data files into dictionaries, binned tables, series, and single records. It also contains metadata such as string offsets and the time boundaries for each bin.

The fields are defined as follows:
\begin{itemize}
  \item \textbf{name}: The name of the recorded game.
  \item \textbf{recordedGame}: The filename of the source CaptureAge recording.
  \item \textbf{dictionaries}: Static key-value pairs used for lookup (e.g. strings).
  \item \textbf{binned}: Windowed values providing snapshots of the game state at a fixed frequency (defined by \texttt{binWidth} in milliseconds).
  \item \textbf{series}: Time-stamped streams of symbolic data (events).
  \item \textbf{records}: One-off values that apply to the entire game and do not evolve over time (e.g. map terrain).
  \item \textbf{metadata}: Contains auxiliary information:
    \begin{itemize}
      \item \texttt{civStringIdOffset} / \texttt{colorStringIdOffset}: Offsets added to IDs to retrieve the corresponding name from the string dictionary.
      \item \texttt{binTimes}: An array of inclusive start/end times (in milliseconds) corresponding to each bin index. Gaps between windows exist due to the discrete nature of the simulation time steps.
    \end{itemize}
\end{itemize}

\begin{verbatim}
{
  "name": "Kotd4-Groups S06G2-Mbl vs Vinchester",
  "recordedGame": ".../Kotd4-Groups S06G2-Mbl vs Vinchester.carz",
  "dictionaries": [
    { "type": "strings", "path": "./strings.arrow" }
  ],
  "binned": [
    { "type": "players", "binWidth": 60000, "path": "./players.arrow" },
    { "type": "researchStates", "binWidth": 60000,
      "path": "./research-states.arrow" },
    { "type": "entities", "binWidth": 60000, "path": "./entities.arrow" },
    { "type": "masters", "binWidth": 60000, "path": "./masters.arrow" },
    { "type": "world", "binWidth": 60000, "path": "./world.arrow" }
  ],
  "series": [
    { "type": "events", "path": "./events.arrow" }
  ],
  "records": [
    { "type": "settings", "path": "./settings.arrow" },
    { "type": "map", "path": "./map.arrow" }
  ],
  "metadata": {
    "civStringIdOffset": 10270,
    "colorStringIdOffset": 10631,
    "binTimes": [
      [ 4841, 60014 ],
      [ 60039, 120021 ],
      ...
      [ 3480496, 3531404 ]
    ]
  }
}
\end{verbatim}

\subsection{Entry Shapes}

Each Arrow file contains rows ("entries") built from the corresponding schema.
Columns marked (dict) are Arrow dictionary encoded.
Optional columns (nullable in schema) are noted in the \textbf{Condition} column, specifying the \texttt{modelType} or logical condition that enables them.

\subsubsection*{Notes}
\begin{itemize}
  \item We could not use Arrow's Unions for polymorphic types in Arrow due to the level of support in the JavaScript implementation. Instead we indicate the type using a \texttt{modelType} dictionary column, and optional columns for each type. We use the optional column to indicate for which type the column is valid.
  \item Some values are not well described below (in particular id-mapping), please refer to Advanced Genie Editor (included with the AoE2 DE tools) for further details. If some values are unclear, please contact the author for clarification.
  \item Long identifiers in the tables below may be hyphenated (e.g. with \texttt{-}) for display purposes; these hyphens are not part of the actual field names or types.
\end{itemize}
\par\noindent\rule{\textwidth}{0.4pt}

\subsubsection*{Common Values}

\begin{itemize}
  \item \texttt{sage:modelType} (dict \texttt{Uint8}) — present in multi‑model tables (\texttt{entities}, \texttt{masters}, \texttt{events}).
  \item \texttt{sage:binId} (\texttt{Uint32}) — present in binned tables, increasing id, can be used with descriptor's \texttt{binTimes} array to retrieve window start/end times.
  \item DUCT lifecycle (only in \texttt{entities}):
    \begin{itemize}
      \item \texttt{sage:ductType} (dict \texttt{Int8}) values: Deleted, Updated, Created, Transient.
      \item \texttt{sage:creationTime} (\texttt{Uint32}, optional), kept at all times to keep track on the first time this entry was seen
      \item \texttt{sage:deletionTime} (\texttt{Uint32}, optional), kept in the last window where it was present to indicate the last time it was seen
    \end{itemize}
\end{itemize}

List element wrapping:

\begin{itemize}
  \item \texttt{List<Struct>} columns (e.g. production queues, attribute lists) store an empty list \texttt{[]} when logically present but source array is null, \texttt{null} when not present.
\end{itemize}

\par\noindent\rule{\textwidth}{0.4pt}

\subsubsection*{\texttt{settings.arrow} (record)}

This contains only three fields, which support the map.arrow interpretation.
\begin{longtable}{@{}p{0.24\textwidth}p{0.14\textwidth}p{0.22\textwidth}p{0.30\textwidth}@{}}
  \toprule
  Column & Type & Condition & Notes \\
  \midrule
  \endhead
  mapWidth & \texttt{Uint16} & & Map width in tile count \\
  mapHeight & \texttt{Uint16} & & Map height in tile count \\
  mapName & \texttt{Utf8} & & Always ``Arabia'' for our data set \\
  \bottomrule
\end{longtable}

\subsubsection*{\texttt{map.arrow} (record)}

One row per tile.

\begin{longtable}{@{}p{0.24\textwidth}p{0.14\textwidth}p{0.22\textwidth}p{0.30\textwidth}@{}}
  \toprule
  Column & Type & Condition & Notes \\
  \midrule
  \endhead
  x & \texttt{Int32} & & Tile x-position \\
  y & \texttt{Int32} & & Tile y-position \\
  elevationLevel & \texttt{Uint8} & & Tile height level \\
  terrainType & \texttt{Uint8} & & Tile terrain type (id refers to in-game terrain id e.g. ``Grass 2'') \\
  \bottomrule
\end{longtable}

\subsubsection*{\texttt{world.arrow} (binned)}

\begin{longtable}{@{}p{0.24\textwidth}p{0.14\textwidth}p{0.22\textwidth}p{0.30\textwidth}@{}}
  \toprule
  Column & Type & Condition & Notes \\
  \midrule
  \endhead
  \texttt{sage:binId} & \texttt{Uint32} & & Bin identifier \\
  time & \texttt{Uint32} & & World time (ms) \\
  woodPrice & \texttt{Float32} & & Base price in gold to buy or sell wood \\
  foodPrice & \texttt{Float32} & & Base price in gold to buy or sell food \\
  stonePrice & \texttt{Float32} & & Base price in gold to buy or sell stone \\
  \bottomrule
\end{longtable}
Note: Market events report actual values themselves.

\subsubsection*{\texttt{strings.arrow} (dictionary)}

\begin{longtable}{@{}p{0.24\textwidth}p{0.14\textwidth}p{0.22\textwidth}p{0.30\textwidth}@{}}
  \toprule
  Column & Type & Condition & Notes \\
  \midrule
  \endhead
  id & \texttt{Int32} & & String ID \\
  string & \texttt{Utf8} & & schema non‑nullable \\
  \bottomrule
\end{longtable}

For our practical use we only used English strings, and filtered to those referenced by any string id, pluse colors and civilizations.

\par\noindent\rule{\textwidth}{0.4pt}

\subsubsection*{\texttt{players.arrow} (binned)}

Basic (non-spatial) player state snapshot.

\begin{longtable}{@{}p{0.24\textwidth}p{0.14\textwidth}p{0.22\textwidth}p{0.30\textwidth}@{}}
  \toprule
  Column & Type & Condition & Notes \\
  \midrule
  \endhead
  \texttt{sage:binId} & \texttt{Uint32} & & Bin identifier \\
  id & \texttt{Int8} & & Player id, typically 0 is Gaia (computer controlled, e.g. resources, props, and animals) \\
  color & \texttt{Int32} & & Color, can be used with the descriptor's \texttt{colorStringIdOffset} to get the name from \texttt{strings.arrow} \\
  name & \texttt{Utf8} & & Player in-game name \\
  civilization & \texttt{Uint32} & & Player civilization id, can be used with the descriptor's \texttt{civStringIdOffset} to get the name from \texttt{strings.arrow} \\
  victoryPoints & \texttt{Int32} & & Score as shown in game \\
  startX & \texttt{Float32} & & x-position of the starting units (typically a Town Center) \\
  startY & \texttt{Float32} & & y-position of the starting units (typically a Town Center) \\
  resigned & \texttt{Bool} & & Has the player resigned? \\
  gameStatus & \texttt{Uint8} & & Player status (e.g. Human, AI) \\
  \bottomrule
\end{longtable}

Dynamic (technically optional, but commonly available):

\begin{longtable}{@{}p{0.24\textwidth}p{0.14\textwidth}p{0.22\textwidth}p{0.30\textwidth}@{}}
  \toprule
  Column & Type & Condition & Notes \\
  \midrule
  \endhead
  food & \texttt{Float32} & Dynamic & Food stockpile \\
  wood & \texttt{Float32} & Dynamic & Wood stockpile \\
  stone & \texttt{Float32} & Dynamic & Stone stockpile \\
  gold & \texttt{Float32} & Dynamic & Gold stockpile \\
  maxPop & \texttt{Float32} & Dynamic & Population Capacity \\
  currentAge & \texttt{Float32} & Dynamic & 0 = Dark Age, etc. \\
  population & \texttt{Float32} & Dynamic & Current Population \\
  razings & \texttt{Float32} & Dynamic & Cumulative number of buildings destroyed \\
  foodTotal & \texttt{Float32} & Dynamic & Cumulative food collected throughout the game \\
  woodTotal & \texttt{Float32} & Dynamic & Cumulative wood collected throughout the game \\
  stoneTotal & \texttt{Float32} & Dynamic & Cumulative stone collected throughout the game \\
  goldTotal & \texttt{Float32} & Dynamic & Cumulative gold collected throughout the game \\
  relicGold & \texttt{Float32} & Dynamic & Cumulative gold collected from holding relics \\
  relicFood & \texttt{Float32} & Dynamic & Cumulative food collected from holding relics \\
  relicWood & \texttt{Float32} & Dynamic & Cumulative wood collected from holding relics \\
  relicStone & \texttt{Float32} & Dynamic & Cumulative stone collected from holding relics \\
  civilianPop & \texttt{Float32} & Dynamic & Number of civilian units (villagers, fishing ships, trade carts, etc.) \\
  militaryPop & \texttt{Float32} & Dynamic & Number of military units \\
  \bottomrule
\end{longtable}

\par\noindent\rule{\textwidth}{0.4pt}

\subsubsection*{\texttt{research-states.arrow} (binned)}

States of all research include:
\begin{itemize}
  \item \texttt{CanResearchLater} (0)
  \item \texttt{CanResearch} (1)
  \item \texttt{Researching} (2)
  \item \texttt{Researched} (3)
  \item \texttt{Queued} (4)
\end{itemize}

Excludes (for storage space reasons):
\begin{itemize}
  \item \texttt{CantResearch} (-1) — technologies unavailable to the player.
\end{itemize}

\begin{longtable}{@{}p{0.24\textwidth}p{0.14\textwidth}p{0.22\textwidth}p{0.30\textwidth}@{}}
  \toprule
  Column & Type & Condition & Notes \\
  \midrule
  \endhead
  \texttt{sage:binId} & \texttt{Uint32} & & Bin identifier \\
  playerId & \texttt{Int8} & & Player id \\
  researchId & \texttt{Int16} & & Unique id for the technology \\
  stringId & \texttt{Int32} & Named Tech & If technology is named (some are not UI facing, such as civilization bonuses) \\
  researchDone & \texttt{Float32} & & Progress \\
  state & \texttt{Uint8} & & Enum (e.g. \texttt{CanResearch}, ...) \\
  buildObjId & \texttt{Int16} &  & Building being researched at \\
  timesResearched & \texttt{Int16} & Repeatable & For repeatable techs \\
  allowMultipleResearch & \texttt{Bool} & Multi-research & For techs allowing multiple \\
  \bottomrule
\end{longtable}

\par\noindent\rule{\textwidth}{0.4pt}

\subsubsection*{\texttt{events.arrow} (series)}

Per event; \texttt{sage:modelType} dictionary over:
\begin{itemize}
  \item \texttt{damage}: An entity deals damage to another.
  \item \texttt{death}: An entity kills another (in addition to damage).
  \item \texttt{monkConversion}: An entity converts another (unless heresy activated, in which case it died).
  \item \texttt{marketTransaction}: A player bought or sold resources.
  \item \texttt{techStatus}: Tech status changed. (tags: \texttt{researching}, \texttt{completed}, \texttt{canceled}, \texttt{forceCanceled})
  \item \texttt{garrison}: An entity enters or leaves another. (tags: \texttt{garrison}, \texttt{ungarrison})
  \item \texttt{unitCreated}: A unit is created.
  \item \texttt{lifeCycle}: An entity changed its life cycle (tags: \texttt{created}, \texttt{constructed}, \texttt{underConstruction}, \texttt{dead}).
\end{itemize}

\begin{longtable}{@{}p{0.24\textwidth}p{0.14\textwidth}p{0.22\textwidth}p{0.30\textwidth}@{}}
  \toprule
  Column & Type & Condition & Notes \\
  \midrule
  \endhead
  \texttt{sage:modelType} & dict \texttt{Uint8} & & Event kind \\
  worldTime & \texttt{Uint32} & & World time this event happened (ms) \\
  worldX, worldY, worldZ & \texttt{Float32} & Spatial & Spatial events (unit / position relevant) \\
  objectEntity & \texttt{Int32} & & Target / object ID \\
  subjectEntity & \texttt{Int32} & & Source / or actor ID \\
  objectPlayerId & \texttt{Uint8} & & Player ID of the target \\
  subjectPlayerId & \texttt{Uint8} & & Player ID of the source \\
  resourceType & \texttt{Uint8} & \texttt{marketTransaction} & Resource type ID (e.g. Food, Wood) \\
  resourceAmount & \texttt{Uint16} & \texttt{marketTransaction} & Amount of resource bought/sold \\
  goldAmount & \texttt{Uint16} & \texttt{marketTransaction} & Gold cost or revenue \\
  techId & \texttt{Uint16} & \texttt{techStatus} & Technology ID \\
  heresyActivated & \texttt{Bool} & \texttt{monkConversion} & True if Heresy tech prevented conversion \\
  tag & dict \texttt{Uint8} & (various) & See list above \\
  \bottomrule
\end{longtable}

\par\noindent\rule{\textwidth}{0.4pt}

\subsubsection*{\texttt{entities.arrow} (binned)}

\texttt{sage:modelType} dictionary over:
\texttt{Entity}, \texttt{DoppleEntity}, \texttt{ActionEntity}, \texttt{CombatEntity}, \texttt{MissileEntity}, \texttt{BuildingEntity}.

This contains all entities present in the game at each time window, including units, buildings, resources, and projectiles, as well as doppleganger entities used to represent Fog of War states. Notably, it also includes entities that are in different states of their lifecycle:

\begin{itemize}
  \item \texttt{Building} (0)
  \item \texttt{AlmostAlive} (1)
  \item \texttt{Alive} (2)
  \item \texttt{Dead} (3)
  \item \texttt{AlmostDead} (4)
  \item \texttt{ReallyDead} (5)
  \item \texttt{Undead} (6)
  \item \texttt{Remove} (7)
  \item \texttt{Going} (8)
\end{itemize}

Some of these states are transient within the game loop and will likely not appear in the binned snapshots, but they are included for completeness.

Base (always present):

\begin{longtable}{@{}p{0.24\textwidth}p{0.14\textwidth}p{0.22\textwidth}p{0.30\textwidth}@{}}
  \toprule
  Column & Type & Condition & Notes \\
  \midrule
  \endhead
  \texttt{sage:modelType} (dict) & \texttt{Uint8} & & Entity model type \\
  \texttt{sage:binId} & \texttt{Uint32} & & Bin identifier \\
  \texttt{sage:ductType} (dict) & \texttt{Int8} & & Data update type (Created, Updated, etc.) \\
  \texttt{sage:creationTime} & \texttt{Uint32} & & Time entity was created \\
  \texttt{sage:deletionTime} & \texttt{Uint32} & & Time entity was deleted \\
  id & \texttt{Int32} & & Entity ID \\
  masterId & \texttt{Int16} & & When negative it uses a custom master field\footnote{For example this can happen when a unit is converted from another player, e.g. with more advanced upgrades.}, which is added to the \texttt{masters.arrow} \\
  refMasterId & \texttt{Int16} & & Used for special cases, e.g., all villagers types (male farmer, female forager) have their own masterIds, refMasterId refers the same id for all. In most cases this is the recommended identifier. \\
  ownerId & \texttt{Uint8} & & Player ID owning the entity \\
  worldX & \texttt{Float32} & & X coordinate \\
  worldY & \texttt{Float32} & & Y coordinate \\
  worldZ & \texttt{Float32} & & Z coordinate (height) \\
  heldAttributeAmount & \texttt{Float32} & & Amount of resource held (e.g. by villager) \\
  heldAttributeType & \texttt{Int16} & & Type of resource held \\
  state & \texttt{Uint8} & & Entity state (e.g. Alive, Dead) \\
  doppleFlag & \texttt{Bool} & & Is a virtual clone of another player's building or resource node, to show the last seen variation in the Fog of War \\
  type & \texttt{Uint8} & & Unit Class Type \\
  hp & \texttt{Float32} & & Current hit points \\
  insideObjId & \texttt{Int32} & & ID of object this entity is inside (garrisoned) \\
  groupId & \texttt{Int32} & & Group ID for formation/selection \\
  workerNum & \texttt{Int32} & & Number of workers (e.g. on a building) \\
  lastSeen & \texttt{List\-<Uint32>} (or \texttt{[undefined]}) & & Time entry for lastSeen for each player (Gaia = 0) \\
  \bottomrule
\end{longtable}

Optional / modelType‑specific:

\begin{longtable}{@{}p{0.24\textwidth}p{0.14\textwidth}p{0.22\textwidth}p{0.30\textwidth}@{}}
  \toprule
  Column & Type & Condition & Notes \\
  \midrule
  \endhead
  doppledObjectId & \texttt{Int32} & \texttt{DoppleEntity} & value = entity.id \\
  doppledPlayerId & \texttt{Int8} & \texttt{DoppleEntity} & value = entity.ownerId \\
  currentAction & Action & \texttt{ActionEntity}, \texttt{CombatEntity} & Current action being performed \\
  formationType & \texttt{Uint8} & \texttt{ActionEntity}, \texttt{CombatEntity} & Current formation ID \\
  attackStance & \texttt{Uint8} & \texttt{ActionEntity}, \texttt{CombatEntity} & Stance (Aggressive, Defensive, etc.) \\
  firedFromId & \texttt{Int32} & \texttt{MissileEntity} & ID of unit that fired this missile \\
  linkedOwnerId & \texttt{Int32} & \texttt{BuildingEntity} & ID of owner for linked buildings (e.g. farms) \\
  current\-Production\-Queue\-Action & Action & \texttt{BuildingEntity} & Action currently at head of queue \\
  gatherPointX / gatherPointY / gatherPointZ & \texttt{Float32} & \texttt{BuildingEntity} & Rally point coordinates \\
  gatherPointTargetId & \texttt{Int32} & \texttt{BuildingEntity} & ID of target for rally point \\
  buildPts & \texttt{Float32} & \texttt{BuildingEntity} & Construction progress points \\
  originalOwnerId & \texttt{Int8} & \texttt{BuildingEntity} & Player ID who originally built this \\
  built & \texttt{Bool} & \texttt{BuildingEntity} & True if construction is complete \\
  relicCount & \texttt{Int32} & \texttt{BuildingEntity} & Number of relics garrisoned \\
  productionQueue & \texttt{List\-<Production}\-\texttt{Queue}\-\texttt{Record>} & \texttt{BuildingEntity} & List of units/techs queued \\
  \bottomrule
\end{longtable}

\par\noindent\rule{\textwidth}{0.4pt}

\subsubsection*{\texttt{masters.arrow} (binned)}

Binned; \texttt{sage:modelType} dictionary over:
\texttt{MasterEntity}, \texttt{MasterMovingEntity}, \texttt{MasterActionEntity}, \texttt{MasterCombatEntity}, \texttt{MasterMissileEntity}, \texttt{MasterBuildingEntity}.

Base (always):

\begin{longtable}{@{}p{0.24\textwidth}p{0.14\textwidth}p{0.22\textwidth}p{0.30\textwidth}@{}}
  \toprule
  Column & Type & Condition & Notes \\
  \midrule
  \endhead
  \texttt{sage:binId} & \texttt{Uint32} & & Bin identifier \\
  \texttt{sage:modelType} (dict) & \texttt{Uint8} & & Master model type \\
  id & \texttt{Int16} & & Master ID (Unit/Building Type ID) \\
  playerId & \texttt{Uint8} & & Player ID \\
  type & \texttt{Uint8} & & Unit class type \\
  stringId & \texttt{Int32} & & Name string ID \\
  objectGroup & \texttt{Int16} & & Object group ID \\
  hp & \texttt{Int16} & & Max Hit Points \\
  objCapacity & \texttt{Uint8} & & Garrison capacity \\
  radiusX & \texttt{Float32} & & Collision radius X \\
  radiusY & \texttt{Float32} & & Collision radius Y \\
  radiusZ & \texttt{Float32} & & Collision radius Z \\
  available & \texttt{Bool} & & True if available to player \\
  disabled & \texttt{Bool} & & True if disabled \\
  constructionRadiusX & \texttt{Float32} & & Foundation size X \\
  constructionRadiusY & \texttt{Float32} & & Foundation size Y \\
  resourceGroup & \texttt{Uint8} & & Resource group ID \\
  attributesHeld & \texttt{List\-<Attribute}\-\texttt{Value>} & & Resources carried capacity/type \\
  \bottomrule
\end{longtable}

Optional / modelType mapping (prefixes \texttt{Master} and suffixes \texttt{Entity} omitted in Condition for brevity):

\begin{longtable}{@{}p{0.24\textwidth}p{0.14\textwidth}p{0.22\textwidth}p{0.30\textwidth}@{}}
  \toprule
  Column & Type & Condition & Notes \\
  \midrule
  \endhead
  speed & \texttt{Float32} & \texttt{Action}, \texttt{Combat}, \texttt{Moving} & Movement speed \\
  baseArmor & \texttt{Int16} & \texttt{Combat} & Base armor value \\
  armor & \texttt{List\-<Armor}\-\texttt{Weapon}\-\texttt{Info>} & \texttt{Combat} & Armor classes and values \\
  weapon & \texttt{List\-<Armor}\-\texttt{Weapon}\-\texttt{Info>} & \texttt{Combat} & Attacks and damage values \\
  weaponRange & \texttt{Float32} & \texttt{Combat} & Attack range \\
  areaDamage & \texttt{Float32} & \texttt{Combat} & Blast radius \\
  speedOfAttack & \texttt{Float32} & \texttt{Combat} & Reload time \\
  combatAbility & \texttt{Uint8} & \texttt{Combat} & Special combat flags \\
  origArmor & \texttt{Int16} & \texttt{Combat} & Original armor (before buffs) \\
  origWeapon & \texttt{Int16} & \texttt{Combat} & Original weapon \\
  origWeaponRange & \texttt{Float32} & \texttt{Combat} & Original range \\
  origSpeedOfAttack & \texttt{Float32} & \texttt{Combat} & Original reload time \\
  buildPtsRequired & \texttt{Float32} & \texttt{Combat}, \texttt{Building} & Total build points needed \\
  minimumWeaponRange & \texttt{Float32} & \texttt{Combat} & Minimum range \\
  buildInventory & \texttt{List\-<Attribute}\-\texttt{Value>} & \texttt{Combat} & Cost to build \\
  \bottomrule
\end{longtable}

\par\noindent\rule{\textwidth}{0.4pt}

\subsubsection*{Nested Struct: AttributeValue}

\begin{longtable}{@{}p{0.24\textwidth}p{0.14\textwidth}p{0.22\textwidth}p{0.30\textwidth}@{}}
  \toprule
  Field & Type & Condition & Notes \\
  \midrule
  \endhead
  type & \texttt{Int16} & & Attribute type ID (e.g. Gold, Wood) \\
  amount & \texttt{Float32} & & Quantity \\
  flag & \texttt{Uint8} & & Usage flag (e.g. some resources are return upon death) \\
  \bottomrule
\end{longtable}

\par\noindent\rule{\textwidth}{0.4pt}

\subsubsection*{Nested Struct: ProductionQueueRecord}

\begin{longtable}{@{}p{0.24\textwidth}p{0.14\textwidth}p{0.22\textwidth}p{0.30\textwidth}@{}}
  \toprule
  Field & Type & Condition & Notes \\
  \midrule
  \endhead
  unitId & \texttt{Int16} & & ID of unit being produced \\
  refUnitId & \texttt{Int16} & & Reference ID for unit type \\
  techId & \texttt{Int16} & & ID of technology being researched \\
  unitCount & \texttt{Int16} & & Number of units in this batch \\
  \bottomrule
\end{longtable}

\par\noindent\rule{\textwidth}{0.4pt}

\subsubsection*{Nested Struct: Action}

\begin{longtable}{@{}p{0.24\textwidth}p{0.14\textwidth}p{0.22\textwidth}p{0.30\textwidth}@{}}
  \toprule
  Field & Type & Condition & Notes \\
  \midrule
  \endhead
  \texttt{sage:modelType} & dict \texttt{Uint8} & & Action model type \\
  type & \texttt{Int16} & & Action Type ID \\
  state & \texttt{Int8} & & Action State \\
  targetId & \texttt{Int32} & & Target Entity ID \\
  targetX & \texttt{Float32} & & Target X coordinate \\
  targetY & \texttt{Float32} & & Target Y coordinate \\
  targetZ & \texttt{Float32} & & Target Z coordinate \\
  timer & \texttt{Float32} & & Duration \\
  objId & \texttt{Int16} & \texttt{MakeObject} & Object ID being created \\
  workDone & \texttt{Float32} & \texttt{MakeObject} & Progress \\
  techId & \texttt{Int16} & \texttt{MakeTech} & Tech ID being researched \\
  researchProgress & \texttt{Float32} & \texttt{MakeTech} & Research progress \\
  startTime & \texttt{Float32} & \texttt{MakeTech} & Start time \\
  targetType & \texttt{Int32} & \texttt{Gather}, \texttt{Hunt} & Target type \\
  wasSameOwner & \texttt{Bool} & \texttt{Convert} & Was same owner? \\
  requiredRange & \texttt{Float32} & \texttt{Convert} & Required range \\
  totalTimer & \texttt{Float32} & \texttt{Convert} & Total timer \\
  taskWorkVal1 & \texttt{Float32} & \texttt{Convert} & Task work value 1\footnote{Related to complex conversion mechanics: see: \url{https://ageofempires.fandom.com/wiki/Conversion\#Conversion_mechanics}} \\
  taskWorkVal2 & \texttt{Float32} & \texttt{Convert} & Task work value 2 \\
  wonderTime & \texttt{Float32} & \texttt{Wonder} & Wonder time \\
  \bottomrule
\end{longtable}

\par\noindent\rule{\textwidth}{0.4pt}

\subsubsection*{Nested Struct: ArmorWeaponInfo}

\begin{longtable}{@{}p{0.24\textwidth}p{0.14\textwidth}p{0.22\textwidth}p{0.30\textwidth}@{}}
  \toprule
  Field & Type & Condition & Notes \\
  \midrule
  \endhead
  type & \texttt{Int16} & & Armor or Attack class ID \\
  value & \texttt{Int16} & & Value (armor or damage) \\
  \bottomrule
\end{longtable}

\par\noindent\rule{\textwidth}{0.4pt}

\section{Data Transformations}\label{appendix:examples}

This appendix contains the data structures output by each stage of the pipeline. These examples illustrate the transformation of data from raw ingestion to the final preprocessed format used for realization.

\subsection{Data Ingestion}
\label{sec:example-ingestion}
\begin{lstlisting}[language=JavaScript, caption={Raw event input stream to Signal Analysis.}, label={lst:signal-input}]
{
    type: "damage" | "death",
    worldX: number,
    worldY: number,
    worldZ: number,
    subjectEntity: number,      // Unit or Building causing the event
    subjectPlayerId: number,
    objectEntity: number,       // Unit or Building receiving the event
    objectPlayerId: number
}

{
    type: "lifeCycle",
    tag: "created" | "constructed" | "underConstruction" | "dead",
    subjectEntity: number,
    subjectPlayerId: number,
    worldX: number,
    worldY: number,
    worldZ: number
}
\end{lstlisting}

\subsection{Signal Analysis}
\label{sec:example-signal}

\begin{lstlisting}[language=JavaScript, caption={Clustered and filtered output from Signal Analysis.}, label={lst:signal-output}]
{
    type: "battle",
    startTimeMs: number,    // Interval event: start time
    endTimeMs: number,      // Interval event: end time
    worldX: number,
    worldY: number,
    data: InstantEvent[]    // Array of clustered "damage" and "death" events
}

{
    type: "buildingPlaced" | "buildingConstructed",
    playerId: number,
    worldTime: number,      // Instant event: time in ms
    worldX: number,
    worldY: number,
    worldZ: number,
    entityId: number,
    refMasterId: number     // Used to look up building type
}

{
    type: "minuteTick",
    worldTime: number,      // Time in ms
    minute: number          // Minute number
}
\end{lstlisting}

\subsubsection{Signal Analysis Implementation Details}
\label{appendix:signal-analysis-details}

\paragraph{Spatiotemporal Clustering}
We use a spatiotemporal variant of DBSCAN (ST-DBSCAN)\cite{shi2019, ansari2020} to cluster damage and kill events into fights. The main text (Section~\ref{sec:pipeline-implementation}; Figure~\ref{fig:online-temporal-spatial-clustering}) describes the method choice, requirements, and online behavior. This appendix focuses on implementation specifics: Figure~\ref{lst:online-dbscan} provides the per-event procedure, and the following section lists the concrete parameter values used in the prototype.

\begin{figure}[H]
\begin{verbatim}
Definitions:
    e : The current incoming event from the stream
    C_open : Set of currently open clusters
    N : Set of noise points
    dist(a, b) : Spatial distance between events a and b
    time(a, b) : Temporal difference between events a and b
    eps : Maximum spatial distance threshold
    tau : Maximum temporal difference threshold
    minPts : Minimum number of points to form a cluster

Algorithm ProcessEvent(e):
    // 1. Identify and close timed-out clusters
    closed_clusters <- {}
    for each cluster c in C_open:
        if time(e, c.last_event) > tau:
            move c from C_open to closed_clusters

    // 2. Find neighbors for the current event
    neighbors <- {}
    for each cluster c in C_open:
        // Check if e is density-reachable from the cluster
        if exists p in c such that dist(e, p) <= eps:
            merge c into neighbors
            remove c from C_open

    noise_neighbors <- { n in N | dist(e, n) <= eps and time(e, n) <= tau }

    // 3. Update clusters or noise
    if neighbors is empty and size(noise_neighbors) < minPts:
        add e to N
    else:
        // new cluster from union of neighbors,
        //   noise_neighbors, and incoming event {e}
        new_cluster <- neighbors U noise_neighbors U {e}
        add new_cluster to C_open
        remove noise_neighbors from N

    // 4. Clean up old noise
    remove { n in N | time(e, n) > tau } from N

    // 5. Return completed clusters
    return closed_clusters
\end{verbatim}
  \caption{Online spatiotemporal clustering of damage events - it operates on a per incoming event basis, and can close off old clusters as new events arrive.}
  \label{lst:online-dbscan}
\end{figure}

Originally our implementation could buffer closed clusters as long as ongoing battles exist. But we disabled this behavior in the final prototype to ensure Interval Events were emitted as soon as possible and as single battles, supporting the baseline evaluation.

\paragraph{Clustering Parameters}
The following values were found to perform well within our iteration goals for Age of Empires II battle clusters:
\begin{itemize}
  \item \textbf{Spatial Distance (\texttt{eps}):} 14 square tiles --- Maximum distance in tiles between events to consider them part of the same cluster. (In some cases, e.g., Trebuchets fighting each other from a distance, this may be too small, but generally works well for most battles.)
  \item \textbf{Temporal Distance (\texttt{tau}):} 60 seconds --- Maximum game time between events to consider them part of the same cluster. This is per local cluster, so clusters can close off while others are still ongoing.
  \item \textbf{Minimum Points (\texttt{minPts}):} 10 events --- Minimum of 10 damage or kill events to form a cluster. (For every kill event, a damage event is also emitted, so implicitly battles with kills are more likely to be detected, but this is a rare case.)
\end{itemize}

We also experimented with a $maxSilenceTime$, ensuring that clusters were returned only if no events occurred for a certain duration globally, but found the $\tau$ parameter sufficient.




\subsection{Data Interpretation}
\label{sec:example-interpretation}
\begin{lstlisting}[language=JavaScript, caption={Enriched fight event with contextual information.}, label={lst:enriched-output}]
{
  type: "battle";
  startTimeMs: number;
  endTimeMs: number;
  location: string; // named location
  worldX: number;
  worldY: number;
  data: StudyInstant[];
  player1: { // player1 is used over arrayer index for easier LLM parsing
    participants: string[]; // all participating entities (as receiver or dealer of damage)
    casualties: string[]; // resulting casualties
    survivors: string[]; // resulting survivors
  };
  player2: {...}
  fightType: string; // as classified (e.g., "raid")
  beforeState: {...} // before and after player states are added here (same as matchIntro below)
  afterState: {...}
}

{
  type: "buildingPlaced" | "buildingConstructed";
  playerId: number;
  entityId: number;
  refMasterId: number;
  buildingName: string; // used entityId, then refMasterId to look up building type
  worldX: number;
  worldY: number;
  worldZ: number;
  worldTime: number;
  location: string; // named location based on query
  isForward: boolean;
  type: "castle" | "townCenter" | "tower"; // building type category, castle and tower have subtypes
};

{
    type: "mapIntro"
    timeMs: number;
    minute: number;
    mapName: string; // e.g., "Arabia"
    player1:{
      resourceData: {
        ["forageBushes" | "huntables" | "herdables" | "lurables" | "stoneMines" | "goldMines" | "hills" | "woodlines" ]: {
            size: number; // number of entities or tiles in cluster
            forwardedness: number; /
            playerRelativePosition: string;
            playerId: number;
            height: number; // only for hills
        }
      }
    }
    player2:{...}
}

{
    type: "matchIntro"
    timeMs: number;
    minute: number;
    player1:{
        name: string;
        color: string;
        id: number;
        civilization: string;

        unitCount: {
            unitName: string;
            count: number;
            forwardCount: number;
        };
        buildingCount: {
            unitName: string;
            count: number;
            forwardCount: number;
        };
        technologies: [{name: string, state: string}]; //
        score: number;
        totalMilitaryUnits: {[name: string]: number;}[];
        totalEcoUnits: {[name: string]: number;}[];
        expertInsights: { // If Expert Insights module is active
            looksLikeStrategy: string;
            playerTrivia: string;
            civilizationInfo: string;
        };
    }
    player2:{...}
}

{ // minute tick event are propagated
    type: "minuteTick",
    worldTime: number,
    minute: number
}
\end{lstlisting}

\subsection{Document Planning}
\label{sec:example-planning}
\begin{lstlisting}[language=JavaScript, caption={Document Planner output: A content unit containing grouped events.}, label={lst:planner-output}]
{
    type: "multipleEvents"
    startTimeMs: number;
    endTimeMs: number;
    battles: [...],
    constructions: [...],
    strategyUpdate: {...} // if Expert Insights module is active only
}

{
    type: "mapIntro" // as is
    ...
}

{
    type: "matchIntro" // as is
    ...
}
\end{lstlisting}


\begin{figure}[htbp]
\begin{lstlisting}[
  language={},
  basicstyle=\small\ttfamily,
  backgroundcolor=\color{white},
  numbers=none,
  frame=single,
  xleftmargin=1em,
  xrightmargin=1em,
  breaklines=true,
  label=alg:doc-planning,
  caption={Document Planning flush logic for the \textit{Event Structuring} module. MatchIntro and MapQuality events bypass the buffer and are realized immediately. All other events are deferred and flushed as a grouped content unit when one of three conditions is met.}
]
battlesInflationCounter <- 1   // cooldown: suppress rapid battle reports
deferredQueue <- []
lastFlushTime <- 0

for each event in stream:

  // Intro and Map events: bypass buffer, realize immediately
  if event.type in {intro, mapQuality}:
    realize(event)
    continue

  // All other events: interpret and push to deferred queue
  interpreted <- interpret(event)
  if interpreted != null:
    deferredQueue.push({type, time: event.worldTime, data: interpreted})
    // return null -- no immediate realization

  // On each minute tick: decrement battle cooldown
  if event.type == minuteTick:
    battlesInflationCounter <- max(0, battlesInflationCounter - 1)

  // Flush conditions (checked only on minute ticks)
  if event.type == minuteTick:
    hasBattles  <- deferredQueue.any(type == "battles")
    hasEndGame  <- deferredQueue.any(type == "endGame")
    timeSinceFlush <- event.worldTime - lastFlushTime

    shouldFlush <- (hasBattles AND battlesInflationCounter == 0)
                OR (timeSinceFlush > 180000)   // 3-minute timeout
                OR hasEndGame

    if shouldFlush:
      if hasBattles:
        battlesInflationCounter <- 5  // 5-tick cooldown after battle report

      grouped <- {
        battles:         deferredQueue.filter(type == "battles"),
        technologies:    deferredQueue.filter(type == "technologies"),
        buildings:       deferredQueue.filter(type == "buildingEvent"),
        strategyUpdates: deferredQueue.filter(type == "strategyUpdate")
      }
      deferredQueue   <- []
      lastFlushTime   <- event.worldTime

      realize(grouped)   // send structured group to LLM Realizer
\end{lstlisting}
\end{figure}

\subsection{Data Preprocessing}
\label{sec:example-preprocessing}
\begin{lstlisting}[language=JavaScript, caption={Document Preprocessing output: Including readable timestamps and full text player names.}, label={lst:preprocessing-output}]
  [previousRealisations],
  {
    "type": "gameEnded",
    "winner": 2,
    "winnerName": "WallMaster77",
    "worldTime": 987845,
    "timestamp": "16:27"
  }
\end{lstlisting}
\subsection{Realization}
\label{sec:example-realization}
An example without Expert Insights or Event Structuring modules active, but the Writing Style module is active.
\begin{quote}
  Welcome to Arabia! In this matchup, we have CastleDropper, playing as the Yellow Bulgarians. Facing them is WallMaster77, representing the Blue Indians. The action is just getting started!
\end{quote}

\begin{quote}
  Good forward main gold and stone for CastleDropper, and their woodlines are close. WallMaster77 also has close main gold, stone, and readily available woodlines. Standard Arabia spawns for both!
\end{quote}

\begin{quote}
  CastleDropper is pushing forward! A Watch Tower is now going up right on the doorstep of WallMaster77! This could be a very aggressive play early on!
\end{quote}

\begin{quote}
  \dots
\end{quote}

\begin{quote}
  A huge battle just concluded at WallMaster77's forward gold! CastleDropper's entire military of Man-at-Arms and Skirmishers was wiped out, losing 9 units and a Villager! WallMaster77 only lost 2 Archers in a crushing defense!
\end{quote}

\begin{quote}
  And that's it! The game has ended! WallMaster77, playing as the Blue Indians, takes the victory! What a comeback by WallMaster77 after early aggression!
\end{quote}

You can see that sometimes it decides still use the past tense for recaps of earlier events, depending on the context.

\section{Fuzzy Logic Strategy Detection}
\label{appendix:fuzzy-logic}

This appendix details the implementation of the fuzzy logic system used in the Expert Insights module for strategy detection. The system allows for the definition of strategies using declarative templates, which are then evaluated against the game state using a fuzzy logic engine implemented in TypeScript.

\subsection*{Strategy Template Example}
Strategies are defined as templates specifying the required conditions in terms of unit counts, building placements, and resource distributions. Listing~\ref{lst:fuzzy-template} shows an example configuration for detecting a "Tower Rush" strategy.

Note that the configuration uses helper functions to generate the underlying fuzzy set definitions. For example, \texttt{atLeast(n)} creates a trapezoidal set that functions as a crisp requirement: it has a vertical "cliff" at \texttt{n}, meaning any value below \texttt{n} has 0 membership, and any value \texttt{n} or above has 1 membership. This allows for hard constraints within the fuzzy logic system.

\begin{lstlisting}[language=JavaScript,caption={Configuration template for Tower Rush strategy},label={lst:fuzzy-template}]
 {
    strategyName: "Tower Rush",
    timeRegion: between(6, 11),
    criteria: {
      unitCriteria: [
        {
          unitType: "Villager",
          count: { trapezoid: [18, 20, 25, 28] },
        },
      ],
      villagerDistribution: {
        stone: atLeast(2), // must have at least 2 on stone (crisp)
      },
      buildingCriteria: [
        {
          unitType: "Watch Tower",
          count: atLeast(1),
          position: "forward",
        },
      ],
      technologyCriteria: [
        {
          technology: "Loom",
          required: true, // Boolean crisp requirement
        },
      ],
      constraints: {
        age: between(0, 1), // Dark to Feudal Age
      },
    },
  }
\end{lstlisting}

\section{Prompt \& Context}
\label{appendix:prompts}

\subsection{Base Prompts}

\subsubsection{Base instruction}
\begin{verbatim}
You are a live esports commentator for Age of Empires 2 matches.
Write engaging, short recaps based on JSON game data provided in each request.

STRICT RULES:
- Only use information from the provided data set - never add external knowledge
  or make assumptions
- Write as if commenting live on the action
- Focus on general trends rather than specific numbers (except occasionally for
  score, economy, army size), you can mention specifics if it support the
  reader's mental image
- If data is missing, omit that information entirely
- Output only your commentary response

Your previous responses are included for context,
but only provide your next commentary segment.
\end{verbatim}

This prompt sets the role and constraints for the LLM agent:

\begin{itemize}
  \item Behave as a live Age of Empires 2 esports commentator.
  \item Write engaging, concise recaps based on provided JSON game data.
  \item Use only the information from the provided data; do not add external knowledge or assumptions.
  \item Write as if commenting live on the action.
  \item Focus on general trends rather than specific numbers, unless specifics help the reader visualize the action.
  \item Omit any information that is missing from the data.
  \item Output only the commentary response, with no additional context or requests.
\end{itemize}

\subsubsection{Output format}

>outputFormat: plain text, **no markup** (i.e. no HTML, Markdown, etc.), **use real newlines** not escaped/stringified newlines
This prompt ensures the output is in plain text format without any markup, which is required for the evaluation setup. The timestamp will be added automatically; do not include it in your response.

\subsubsection{Output Length}
\begin{verbatim}
Feel free to use any number of words, but your entire response has a
{{targetLengthInWords}} words limit.
\end{verbatim}
We use this prompt to set a target number of words, replacing the $targetLengthInWords$ with the given numerically represented number. We do not set an additional explicit hard output token limit in the LLM settings, since this can lead to abrupt cutoffs. The use of bold markdown is to emphasise the importance of this constraint to the LLM. We allow it to generate less words if the content does not require the full amount.

\subsubsection{Output Language}
>outputLanguage: English
This prompt sets the output language for the commentary, which is fixed to English for the current evaluation.

\subsection{Communicative Goal Prompts}
For each event type, we provide a specific communicative goal prompt to guide the LLM in generating relevant commentary. These prompts are designed to focus the commentary on the key aspects of each event type.

\subsubsection*{Match Introduction Prompt}
\begin{verbatim}
Write an engaging introduction including: map name, player names
civilizations, colors, and map positions.
\end{verbatim}
This prompt guides the LLM to generate an engaging introductory prompt.

\subsubsection*{Map Introduction Prompt}
>Analyze key resource locations (main/secondary gold, stone, woodlines) and their strategic implications. Use "mapQuality" data. Focus on gameplay impact rather than tile distances.
This prompt focuses the LLM on analyzing the strategic implications of resource placements on the map. It was prone to realizing exact tile distances, which was not desired, so we added the note to focus on gameplay impact, with moderate success.

\subsubsection*{Event Prompt}
We used a single prompt for all event types, since each has the same aims, and it adapted well to the information that was provided, not hallucinating information to adhere to the prompt. The base prompt is as follows:
\begin{verbatim}
Prioritize major events: highlight key battles, strategic moves, and
significant actions (forward buildings, tech switches). For multiple
important events, summarize each briefly and note simultaneous action.
Defensive and economic actions should only be mentioned briefly and
supportive of other events in the section.

Context metrics:
- Villager distribution: economic advantage
- Army size/value: military strength
- Victory points: overall game state

Report ongoing battles as in-progress; update when resolved. Focus on
engagement and clarity over exhaustive detail.
\end{verbatim}

This prompt aims to guide the LLM through different event types. It lists several important events to focus on, which are provided as events or contextual information. Event without event structuring module we asked it to note simultateous action, since this can happen (e.g., when a technology is researched at the same time as a battle is ongoing, or in the final moments of the game, where all information is gathered for a final prompt). It also frames contextual information and how to use it to inform the commentary. Finally, it instructs the LLM to report ongoing battles as in-progress. To further reduce the length of the output, it instructs the LLM to focus on engagement and clarity over exhaustive detail.

\subsubsection*{Writing style}
\begin{verbatim}
Write an engaging introduction including: map name, player names, civilizations,
colors, and map positions.
\end{verbatim}
This prompt was always included, and effectively operated as a Base Prompt, however by design it's a communicative goal prompt since tone can differ based on event type, but we did not explore this in depth after abandoning it as a ablative module, thus we reverted to a single prompt to reduce variables.

\subsection{Architectural Modules}
\label{appendix:arch-modules-prompts}
Finally, architectural modules modified prompts. This was done by adding "additional notes" to the prompt for each event. In addition to prompts a negative prompt was added to prevent the LLM from applying Analytical Knowledge when Expert Insights module was disabled, this prevented hallucinations and support evaluation of Expert Insights module.

\subsubsection{Expert Insights Module}
When the Expert Insights module was enabled, an additional note was added to the prompt for each event:
\paragraph{Match Introduction} This was a simple prompt addition, simply pointing it at the data to use in the prompt.
>Explain each civilization's strengths and include player trivia from "expertInsights" data.

\paragraph{Map Introduction} The LLM struggled making use of the Expert Insights data for map introductions, often ignoring it, possibly to not knowing enough about Age of Empires 2 intricacies. Thus we gave it explicit instructions to use the data for insights:
\begin{verbatim}
Analyze resource positioning impact: Forward main gold (risky), back main gold
(safe), and strategic implications. Use "expertInsights" data for map analysis.
Apply civilization knowledge:
- Against archer civs: need wide walls or be forced to make defensive
  towers when resources are pressured
- Against scout openings: smaller walls or spearmen may be sufficient
  to fend off early aggression
- Forward gold may force early aggression by the unfortunate player who
  has that, forcing them to pressure the opponent early
- Early pressure can be a Drush to get more time to secure forward resources
  or to enable an early Castle Age
- Stone position crucial for tower rush civs (Bulgarians, Incas, Koreans)
\end{verbatim}

With this is was effective at linking information from the civilization expert insights data to the map resource placements.

\paragraph{Event Realisation} The prompt was simple, as most was provided by the addition of the data for the insights:
\begin{verbatim}
Explain player actions, motivations, and potential game impact with strategic
analysis
\end{verbatim}

\paragraph{Negative Prompt when Disabled} When the Expert Insights module was disabled, the following negative prompt was added to prevent the LLM from applying Analytical Knowledge:
\begin{verbatim}
Report only events from data - no strategic commentary.
Essentially reiterating the base prompt's instruction to not add external
knowledge, but specifically targeting strategic commentary.
\end{verbatim}

\subsubsection{Event Structuring} When the Event Structuring module was enabled, an additional note was added to the prompt for each content unit. We applied this to both standard events and Introductions (which are treated as Synthetic Events). This was the hardest event to tweak, as the LLM would invent Information Asymmetry (e.g., claiming an army already in the opponent's base was yet to be discovered). And, hilariously, Domain Knowledge (civilization bonuses) as things that contributed to Information Asymetry. We focussed on  The final prompt used is as follows:
\begin{verbatim}
Select the most important events from the pre-ordered data to report.
Create suspense by clearly showing simultaneous events using
timestamps. Emphasize information asymmetry:
- What you know but players don't (opponent's actions that a player
  is about to discover)
- What players know but audience doesn't (strategies not yet clear
  to the audience)

For ongoing and parallel events, use transitional phrases such as:
"Meanwhile, on the other side", "But their opponent doesn't know",
or "What are they planning?"

Focus on: newly created armies, early resource locations, uncommon
strategies, unexpected tactical choices.
Note: General civilization knowledge is assumed for both players
and audience.
\end{verbatim}

For map introductions we used the following:
\begin{verbatim}
Maps are random - opponents must scout each other. Predict player
reactions to map layouts: forward gold creates pressure opportunities,
back gold enables safer play.
\end{verbatim}

\subsection{Experiment: Agentic Planning}
We also used an experiment to test in the LLM could take over content selection itself, rather than relying on the Document Planner module. This did not make it into the final evaluation, but is documented here for completeness. It was supplied as an "additional notes" prompt alongside the base prompts for each event passed to the LLM.

\begin{verbatim}
CONTENT FILTERING RULES:
- Skip unimportant or repetitive events to maintain engagement
- "deferredData" entries were deferred by you in previous responses
- NEVER defer if a winner is declared

DEFERRAL SYSTEM:
- To defer events: output only <DEFER>
- Before Castle Age: aim to report at least every minute
- After Castle Age: aim to report roughly every 3 minutes
- Incorporate "deferredData" entries to maintain narrative coherence

ENGAGEMENT PRIORITY: Keep readers engaged over completeness.
\end{verbatim}

This prompt was paired with a system that skips realization if the prompt output was exactly "<DEFER>", allowing the LLM to control which events to comment on. For the next prompt a "deferredData" field was added to the JSON context, containing all previously deferred events to allow the LLM to incorporate them later.
In practice, the LLM did not rely heavily on this ability, and the final evaluation used the Document Planner module for content selection.

\subsection{Target Length Selection}
For each combination of architectural modules, we selected a target word count to achieve a total commentary length of 500–800 words for an average match duration of 45 minutes.

\begin{itemize}
  \item \textbf{Base:}
    \begin{itemize}
      \item Match Introduction: 60 words
      \item Map Introduction: 40 words
      \item Event Report: 40 words
    \end{itemize}
  \item \textbf{Expert Insights:}
    \begin{itemize}
      \item Match Introduction: 60 words
      \item Map Introduction: 40 words
      \item Event Report: 40 words
    \end{itemize}
  \item \textbf{Event Structuring:}
    \begin{itemize}
      \item Match Introduction: 60 words
      \item Map Introduction: 60 words
      \item Event Report: 60 words
    \end{itemize}
  \item \textbf{Both:}
    \begin{itemize}
      \item Match Introduction: 60 words
      \item Map Introduction: 60 words
      \item Event Report: 60 words
    \end{itemize}
\end{itemize}

\subsection{Prompt Structure}
Finally the event data and all previous realization (but not their data), were provided to the LLM as Dynamic Content.
Reiterating on the full prompt structure. There are three types of prompt components:

\begin{itemize}
  \item \textbf{Base Prompts: }these are the same for each prompt, regardless of the request.
  \item \textbf{Communicative Goal Prompts: }these vary based on the type of event being commented on and the architectural modules enabled.
  \item \textbf{Dynamic Context: }this is the JSON content unit provided for each request.
\end{itemize}

\chapter{Evaluation Details}
\label{appendix:evaluation-details}

\begingroup

\section{Pre-Study Checklist \& Recap}
\label{appendix:pre-study-checklist}

Based on the best practice guidelines for human evaluation of Natural Language Generation (NLG) \cite{vanderlee2021} and reporting standards for Linear Mixed Models \cite{meteyard2020}, we established the following checklist for the study preparation and reporting.

\paragraph{Study Goal \& Design}
\begin{itemize}
  \item \textbf{Goal Definition:} This is a \textit{summative} evaluation. The primary goal is to judge the merit of the architectural modules (\textit{Expert Insights} and \textit{Event Structuring}) by testing specific hypotheses regarding their impact on viewer engagement.
  \item \textbf{Hypotheses:} We distinguish between our broader research questions (see Chapter \ref{chap:introduction}) and the specific confirmatory hypotheses listed in Section \ref{sec:hypotheses}.
  \item \textbf{Baselines:} We utilize the "No Modules" configuration as our baseline. This allows us to isolate the specific effects of our architectural contributions against a basic generation system.
  \item \textbf{Design \& Pre-testing:} A Latin Square design variation is used to mitigate order effects. A pilot study was conducted with two participants to verify the questionnaire flow and identify potential errors, though due to the number of permutations, not every condition was fully tested in the pilot.
\end{itemize}

\paragraph{Sampling \& Power Analysis}
\begin{itemize}
  \item \textbf{Target Audience:} As the texts entail complex domain-specific strategies, our target audience consists of existing Age of Empires II players. Recruitment was targeted via community hubs (i.e., Reddit, Discord).
  \item \textbf{Power Analysis:} A power analysis was conducted prior to data collection to determine the required sample size, utilizing effect size estimates derived from our preliminary study.
\end{itemize}

\paragraph{Analysis}
\begin{itemize}
  \item We employ Linear Mixed Models (LMM) and have verified the core assumptions: linearity, normality of residuals, and homoscedasticity were all checked prior to interpreting the results.
\end{itemize}

\paragraph{Ethics \& Transparency}
\begin{itemize}
  \item \textbf{Ethics:} Ethical clearance was obtained from the university ethics board.
  \item \textbf{Open Science:} We have made the anonymized dataset, analysis scripts, and stimuli available via a GitHub repository (same as Appendix~\ref{app:generated-outputs}).
  \item \textbf{Preregistration:} The study was not formally preregistered.
\end{itemize}

\section{Survey Questions}
\renewcommand{\arraystretch}{1.5}
\begin{table}[H]
  \centering\small\sloppy
  \begin{tabularx}{\textwidth}{>{\RaggedRight\arraybackslash}p{3.5cm} >{\RaggedRight\arraybackslash}X >{\RaggedRight\arraybackslash}X}
    \textbf{Code} & \textbf{Question} & \textbf{Options} \\
    \midrule
    AoEUnderstanding & How competent do you think you are at understanding AoE2: Definitive Edition competitive 1v1 Random Map matches? & "Extremely incompetent", "Moderately incompetent", "Slightly incompetent", "Neither competent nor incompetent", "Slightly competent", "Moderately competent", "Extremely competent" \\
    MeasurableSkill & What is your current Elo on the 1v1 ladder for AoE2: Definitive Edition? & 600 - 2000 at 400, 200 increments, "2000 or above" and "I don't have a 1v1 Elo" \\
    ContentEngagement & On average, how often do you watch competitive Age of Empires II videos or live streams? & "Several times per week", "Once a week", "Once a month", "Only during some large events", "Never" \\
    EnglishProficiency & How would you rate your English reading fluency? & "Not fluent at all", "Very limited fluency", "Basic fluency", "Intermediate fluency", "Good fluency", "Very good fluency", "Native-like fluency" \\
  \end{tabularx}
  \caption{Participant background questions.}
  \label{tab:intro-questions}
\end{table}

\renewcommand{\arraystretch}{1.5}
\begin{table}[H]
  \centering\small\sloppy
  \begin{tabularx}{\textwidth}{>{\RaggedRight\arraybackslash}p{4.5cm} >{\RaggedRight\arraybackslash}X}
    \textbf{Motivation/Code} & \textbf{Question} \\
    \midrule
    CompetitiveNature1 & I enjoyed reading about the competitive gameplay. \\
    CompetitiveNature2 & I enjoyed reading about two players seriously competing against each other. \\
    SkillImprovement1 & Reading match updates like this one can help me become a better player. \\
    SkillImprovement2 & Match updates like this one can help to improve my own play by getting ideas from professional players. \\
    GameKnowledge1 & I feel my understanding of Age of Empires 2 added to my enjoyment of reading the match updates. \\
    GameKnowledge2 & I liked reading the match updates because I understand the intricacies and strategies. \\
    SkillAppreciation1 & I liked reading about players going to their limits and show moves I could not typically think of. \\
    SkillAppreciation2 & I liked reading how others can do things in the game that I could never imagine. \\
    EntertainingNature1 & Reading the match updates was fun. \\
    EntertainingNature2 & Reading Age of Empires 2 match updates like this one could be fun to pass the time. \\
    DramaticNature1 & I found reading the match updates very exciting. \\
    DramaticNature2 & I felt hyped and excited when I read the match updates. \\
    VicariousSensation1 & I felt like I was in the game when it was close or coming down to the final moments. \\
    VicariousSensation2 & I could experience how professionals play without actually playing myself. \\
  \end{tabularx}
  \caption{Motivation-based engageability questions presented after each summary, adapted from the Motivation Scale of Esports Spectatorship (MSES)~\cite{qian2020}. Each question is answered on a 7-point Likert scale: "Strongly Disagree", "Disagree", "Somewhat disagree", "Neither agree nor disagree", "Somewhat agree", "Agree", "Strongly agree".}
  \label{tab:engageability-questions}
\end{table}

\renewcommand{\arraystretch}{1.5}
\begin{table}[H]
  \centering\small\sloppy
  \begin{tabularx}{\textwidth}{>{\RaggedRight\arraybackslash}p{2.5cm} >{\RaggedRight\arraybackslash}X >{\RaggedRight\arraybackslash}X}
    \textbf{Code} & \textbf{Question} & \textbf{Options} \\
    \midrule
    Overall & Overall, how would you rate the quality of the text? & "Very poor", "Poor", "Somewhat poor", "Neither good nor poor", "Somewhat good", "Good", "Very good" \\
    WordChoice & How appropriate were the word choices in the match updates? & "Very inappropriate", "Inappropriate", "Somewhat inappropriate", "Neither appropriate nor inappropriate", "Somewhat appropriate", "Appropriate", "Very appropriate" \\
    Readability & How hard was it to read the match updates? & "Extremely difficult", "Moderately difficult", "Slightly difficult", "Neither easy nor difficult", "Slightly easy", "Moderately easy", "Extremely easy" \\
    Logicality & How often did you feel the live updates were presented out of order? & "Always", "Often", "Somewhat often", "Sometimes", "Somewhat Rarely", "Rarely", "Never" \\
    Detail & How would you rate the level of detail in the match updates? & "Far too much", "Too much", "Slightly too much", "Just Right", "Slightly too little", "Too little", "Far too little" \\
  \end{tabularx}
  \caption{Linguistic quality questions presented after each summary, adapted from \cite{callaway2002}.}
  \label{tab:linguistic-questions}
\end{table}

\renewcommand{\arraystretch}{1.5}
\begin{table}[H]
  \centering\small\sloppy
  \begin{tabularx}{\textwidth}{>{\RaggedRight\arraybackslash}p{2.5cm} >{\RaggedRight\arraybackslash}X >{\RaggedRight\arraybackslash}X}
    \textbf{Code} & \textbf{Question} & \textbf{Options} \\
    \midrule
    Skipping & How much did you skip or skim in order to get to the parts you found interesting? & "Not at all", "Very little", "A little", "Some", "Quite a lot", "Most of it", "Nearly all of it" \\
    Understanding & How well do you understand how the match played out? & "Not at all", "Very little", "A little", "Somewhat", "Mostly", "Very well", "Completely" \\
    OpenText & Optionally, share any additional thoughts you have about the match updates you just read. & Open \\
  \end{tabularx}
  \caption{Additional questions presented after each summary.}
  \label{tab:additional-questions}
\end{table}

\renewcommand{\arraystretch}{1.5}
\begin{table}[H]
  \centering\small\sloppy
  \begin{tabularx}{\textwidth}{>{\RaggedRight\arraybackslash}p{2.5cm} >{\RaggedRight\arraybackslash}X >{\RaggedRight\arraybackslash}X}
    \textbf{Code} & \textbf{Question} & \textbf{Options} \\
    \midrule
    AICount & Out of the live match update texts you read, how many do you think were AI-generated? & 0-4 \\
    AISentiment & How did prior suspicion that the text was AI-generated affect your sentiment toward the live match updates? & "Much more negatively", "More negatively", "Somewhat more negatively", "No difference/No suspicion", "Somewhat more positively", "More positively", "Much more positively" \\
    AIUtility & How would you feel about an AI analyzing, summarizing, or shoutcasting matches (including your own) that would not normally get such attention? & "Very negative", "Negative", "Somewhat negative", "Neutral", "Somewhat positive", "Positive", "Very positive" \\
    OpenAll & Below you can optionally share any additional thoughts you have. & Open \\
  \end{tabularx}
  \caption{Final questions presented at the end of the questionnaire.}
  \label{tab:final-questions}
\end{table}

\section{Model Output Tables}

\scriptsize
\begin{longtable}{l l r r r r r r r}
  \caption{\label{app:fixed_effects}Fixed effects from linear mixed-effects models (95\% confidence intervals)}\\
  \toprule
  Outcome & Effect & $\beta$ & SE & df & t & p & 95\% CI low & 95\% CI high\\
  \midrule
  \endfirsthead
  \caption[]{Fixed effects from linear mixed-effects models (continued)}\\
  \toprule
  Outcome & Effect & $\beta$ & SE & df & t & p & 95\% CI low & 95\% CI high\\
  \midrule
  \endhead
  \bottomrule
  \endlastfoot
  \textit{Competitive Nature} & \textit{Intercept} & \textit{3.197} & \textit{0.267} & \textit{29.60} & \textit{11.98} & \textit{< .001} & \textit{2.65} & \textit{3.74}\\*
  Competitive Nature & Event Structuring & 0.047 & 0.074 & 95.99 & 0.64 & 0.522 & -0.10 & 0.19\\*
  \textbf{Competitive Nature} & \textbf{Expert Insights} & \textbf{0.199} & \textbf{0.074} & \textbf{96.49} & \textbf{2.70} & \textbf{0.008} & \textbf{0.05} & \textbf{0.35}\\*
  Competitive Nature & Order & -0.074 & 0.068 & 96.37 & -1.08 & 0.283 & -0.21 & 0.06\\*
  Competitive Nature & Interaction & 0.018 & 0.074 & 96.61 & 0.24 & 0.812 & -0.13 & 0.16\\
  \addlinespace
  \textit{Skill Improvement} & \textit{Intercept} & \textit{2.583} & \textit{0.242} & \textit{31.80} & \textit{10.70} & \textit{< .001} & \textit{2.09} & \textit{3.08}\\*
  Skill Improvement & Event Structuring & -0.049 & 0.060 & 95.96 & -0.82 & 0.413 & -0.17 & 0.07\\*
  \textbf{Skill Improvement} & \textbf{Expert Insights} & \textbf{0.186} & \textbf{0.060} & \textbf{96.67} & \textbf{3.12} & \textbf{0.002} & \textbf{0.07} & \textbf{0.30}\\*
  Skill Improvement & Order & -0.070 & 0.054 & 88.18 & -1.28 & 0.202 & -0.18 & 0.04\\*
  Skill Improvement & Interaction & 0.032 & 0.059 & 96.83 & 0.55 & 0.586 & -0.09 & 0.15\\
  \addlinespace
  \textit{Game Knowledge} & \textit{Intercept} & \textit{3.864} & \textit{0.315} & \textit{28.90} & \textit{12.28} & \textit{< .001} & \textit{3.22} & \textit{4.51}\\*
  Game Knowledge & Event Structuring & 0.016 & 0.088 & 96.00 & 0.18 & 0.854 & -0.16 & 0.19\\*
  Game Knowledge & Expert Insights & 0.143 & 0.088 & 96.47 & 1.62 & 0.109 & -0.03 & 0.32\\*
  Game Knowledge & Order & -0.064 & 0.082 & 97.03 & -0.78 & 0.437 & -0.23 & 0.10\\*
  Game Knowledge & Interaction & -0.081 & 0.088 & 96.58 & -0.92 & 0.359 & -0.26 & 0.09\\
  \addlinespace
  \textit{Skill Appreciation} & \textit{Intercept} & \textit{2.534} & \textit{0.223} & \textit{25.30} & \textit{11.37} & \textit{< .001} & \textit{2.08} & \textit{2.99}\\*
  Skill Appreciation & Event Structuring & -0.040 & 0.057 & 95.96 & -0.70 & 0.483 & -0.15 & 0.07\\*
  Skill Appreciation & Expert Insights & 0.101 & 0.057 & 96.26 & 1.76 & 0.082 & -0.01 & 0.22\\*
  Skill Appreciation & Order & -0.072 & 0.054 & 99.33 & -1.33 & 0.188 & -0.18 & 0.04\\*
  Skill Appreciation & Interaction & -0.107 & 0.057 & 96.32 & -1.86 & 0.066 & -0.22 & 0.01\\
  \addlinespace
  \textit{Entertaining Nature} & \textit{Intercept} & \textit{2.879} & \textit{0.247} & \textit{28.11} & \textit{11.65} & \textit{< .001} & \textit{2.37} & \textit{3.39}\\*
  Entertaining Nature & Event Structuring & 0.109 & 0.072 & 95.98 & 1.52 & 0.132 & -0.03 & 0.25\\*
  Entertaining Nature & Expert Insights & 0.112 & 0.072 & 96.46 & 1.55 & 0.124 & -0.03 & 0.26\\*
  Entertaining Nature & Order & 0.105 & 0.067 & 96.95 & 1.57 & 0.120 & -0.03 & 0.24\\*
  Entertaining Nature & Interaction & 0.001 & 0.072 & 96.57 & 0.02 & 0.987 & -0.14 & 0.14\\
  \addlinespace
  \textit{Dramatic Nature} & \textit{Intercept} & \textit{2.652} & \textit{0.249} & \textit{23.21} & \textit{10.65} & \textit{< .001} & \textit{2.14} & \textit{3.17}\\*
  Dramatic Nature & Event Structuring & 0.014 & 0.076 & 95.90 & 0.18 & 0.856 & -0.14 & 0.16\\*
  Dramatic Nature & Expert Insights & 0.103 & 0.076 & 96.28 & 1.36 & 0.176 & -0.05 & 0.25\\*
  Dramatic Nature & Order & 0.058 & 0.071 & 98.96 & 0.83 & 0.410 & -0.08 & 0.20\\*
  Dramatic Nature & Interaction & 0.087 & 0.076 & 96.37 & 1.15 & 0.252 & -0.06 & 0.24\\
  \addlinespace
  \textit{Vicarious Sensation} & \textit{Intercept} & \textit{2.792} & \textit{0.241} & \textit{25.56} & \textit{11.60} & \textit{< .001} & \textit{2.30} & \textit{3.29}\\*
  Vicarious Sensation & Event Structuring & 0.098 & 0.077 & 95.91 & 1.27 & 0.206 & -0.05 & 0.25\\*
  \textbf{Vicarious Sensation} & \textbf{Expert Insights} & \textbf{0.252} & \textbf{0.077} & \textbf{96.40} & \textbf{3.29} & \textbf{0.001} & \textbf{0.10} & \textbf{0.41}\\*
  Vicarious Sensation & Order & 0.045 & 0.071 & 97.12 & 0.63 & 0.530 & -0.10 & 0.19\\*
  Vicarious Sensation & Interaction & 0.042 & 0.077 & 96.51 & 0.54 & 0.588 & -0.11 & 0.19\\
  \addlinespace
  \textit{Overall Quality} & \textit{Intercept} & \textit{3.818} & \textit{0.201} & \textit{33.00} & \textit{19.02} & \textit{< .001} & \textit{3.41} & \textit{4.23}\\*
  Overall Quality & Event Structuring & 0.077 & 0.104 & 99.00 & 0.74 & 0.461 & -0.13 & 0.28\\*
  Overall Quality & Expert Insights & 0.080 & 0.104 & 99.00 & 0.77 & 0.444 & -0.13 & 0.29\\*
  \textbf{Overall Quality} & \textbf{Order} & \textbf{0.185} & \textbf{0.093} & \textbf{99.00} & \textbf{2.00} & \textbf{0.049} & \textbf{0.00} & \textbf{0.37}\\*
  Overall Quality & Interaction & -0.030 & 0.104 & 99.00 & -0.29 & 0.770 & -0.24 & 0.18\\
  \addlinespace
  \textit{Word Choice} & \textit{Intercept} & \textit{4.280} & \textit{0.188} & \textit{33.00} & \textit{22.77} & \textit{< .001} & \textit{3.90} & \textit{4.66}\\*
  Word Choice & Event Structuring & 0.146 & 0.105 & 99.00 & 1.40 & 0.166 & -0.06 & 0.35\\*
  Word Choice & Expert Insights & 0.012 & 0.105 & 99.00 & 0.12 & 0.905 & -0.20 & 0.22\\*
  Word Choice & Order & 0.169 & 0.094 & 99.00 & 1.81 & 0.074 & -0.02 & 0.36\\*
  Word Choice & Interaction & 0.023 & 0.104 & 99.00 & 0.22 & 0.828 & -0.19 & 0.23\\
  \addlinespace
  \textit{Readability} & \textit{Intercept} & \textit{4.621} & \textit{0.261} & \textit{10.54} & \textit{17.68} & \textit{< .001} & \textit{4.04} & \textit{5.20}\\*
  Readability & Event Structuring & -0.097 & 0.108 & 95.57 & -0.90 & 0.371 & -0.31 & 0.12\\*
  Readability & Expert Insights & -0.044 & 0.108 & 95.95 & -0.41 & 0.685 & -0.26 & 0.17\\*
  Readability & Order & 0.002 & 0.101 & 99.53 & 0.02 & 0.982 & -0.20 & 0.20\\*
  Readability & Interaction & -0.102 & 0.108 & 96.04 & -0.94 & 0.349 & -0.32 & 0.11\\
  \addlinespace
  \textit{Logicality} & \textit{Intercept} & \textit{5.606} & \textit{0.146} & \textit{11.07} & \textit{38.32} & \textit{< .001} & \textit{5.28} & \textit{5.93}\\*
  \textbf{Logicality} & \textbf{Event Structuring} & \textbf{0.216} & \textbf{0.086} & \textbf{94.84} & \textbf{2.52} & \textbf{0.013} & \textbf{0.05} & \textbf{0.39}\\*
  Logicality & Expert Insights & 0.038 & 0.086 & 95.74 & 0.45 & 0.656 & -0.13 & 0.21\\*
  Logicality & Order & 0.133 & 0.078 & 87.38 & 1.70 & 0.092 & -0.02 & 0.29\\*
  Logicality & Interaction & 0.125 & 0.086 & 95.94 & 1.45 & 0.149 & -0.05 & 0.30\\
  \addlinespace
  \textit{Detail Level} & \textit{Intercept} & \textit{4.409} & \textit{0.291} & \textit{10.76} & \textit{15.16} & \textit{< .001} & \textit{3.77} & \textit{5.05}\\*
  Detail Level & Event Structuring & 0.064 & 0.104 & 95.71 & 0.62 & 0.538 & -0.14 & 0.27\\*
  \textbf{Detail Level} & \textbf{Expert Insights} & \textbf{-0.243} & \textbf{0.104} & \textbf{96.00} & \textbf{-2.34} & \textbf{0.022} & \textbf{-0.45} & \textbf{-0.04}\\*
  Detail Level & Order & 0.013 & 0.098 & 99.56 & 0.13 & 0.896 & -0.18 & 0.21\\*
  Detail Level & Interaction & 0.051 & 0.104 & 96.06 & 0.49 & 0.628 & -0.16 & 0.26\\
  \addlinespace
  \textit{Reading Thoroughness} & \textit{Intercept} & \textit{4.561} & \textit{0.328} & \textit{6.31} & \textit{13.92} & \textit{< .001} & \textit{3.77} & \textit{5.35}\\*
  Reading Thoroughness & Event Structuring & 0.145 & 0.117 & 95.73 & 1.24 & 0.218 & -0.09 & 0.38\\*
  Reading Thoroughness & Expert Insights & -0.150 & 0.117 & 95.93 & -1.28 & 0.202 & -0.38 & 0.08\\*
  Reading Thoroughness & Order & -0.061 & 0.111 & 99.09 & -0.56 & 0.579 & -0.28 & 0.16\\*
  Reading Thoroughness & Interaction & 0.018 & 0.117 & 95.98 & 0.15 & 0.878 & -0.21 & 0.25\\
  \addlinespace
  \textit{Understanding} & \textit{Intercept} & \textit{4.462} & \textit{0.265} & \textit{19.11} & \textit{16.83} & \textit{< .001} & \textit{3.91} & \textit{5.02}\\*
  Understanding & Event Structuring & -0.125 & 0.100 & 95.80 & -1.24 & 0.217 & -0.32 & 0.08\\*
  Understanding & Expert Insights & 0.160 & 0.101 & 96.28 & 1.60 & 0.114 & -0.04 & 0.36\\*
  Understanding & Order & -0.069 & 0.093 & 97.74 & -0.74 & 0.463 & -0.25 & 0.12\\*
  Understanding & Interaction & -0.149 & 0.100 & 96.39 & -1.48 & 0.141 & -0.35 & 0.05\\
\end{longtable}
\normalsize

\begin{longtable}{l l l r r}
  \caption{\label{tab:random_effects}Random effect variance components (intercepts). No random slopes were included in the model.}\\
  \toprule
  Outcome & Grouping & Effect & variance & SD\\
  \midrule
  \endfirsthead
  \caption[]{Random effect variance component, intercepts (continued)}\\
  \toprule
  Outcome & Grouping & Effect & variance & SD\\
  \midrule
  \endhead
  \bottomrule
  \endlastfoot
  Competitive Nature & Participant & Intercept & 1.969 & 1.403\\*
  Competitive Nature & Match & Intercept & 0.024 & 0.156\\*
  Competitive Nature & Residual & Observation & 0.712 & 0.844\\
  \addlinespace
  Skill Improvement & Participant & Intercept & 1.759 & 1.326\\*
  Skill Improvement & Match & Intercept & 0.006 & 0.078\\*
  Skill Improvement & Residual & Observation & 0.466 & 0.682\\
  \addlinespace
  Game Knowledge & Participant & Intercept & 2.691 & 1.640\\*
  Game Knowledge & Match & Intercept & 0.039 & 0.198\\*
  Game Knowledge & Residual & Observation & 1.023 & 1.011\\
  \addlinespace
  Skill Appreciation & Participant & Intercept & 1.224 & 1.106\\*
  Skill Appreciation & Match & Intercept & 0.037 & 0.193\\*
  Skill Appreciation & Residual & Observation & 0.432 & 0.657\\
  \addlinespace
  Entertaining Nature & Participant & Intercept & 1.636 & 1.279\\*
  Entertaining Nature & Match & Intercept & 0.025 & 0.159\\*
  Entertaining Nature & Residual & Observation & 0.680 & 0.825\\
  \addlinespace
  Dramatic Nature & Participant & Intercept & 1.485 & 1.219\\*
  Dramatic Nature & Match & Intercept & 0.045 & 0.213\\*
  Dramatic Nature & Residual & Observation & 0.749 & 0.866\\
  \addlinespace
  Vicarious Sensation & Participant & Intercept & 1.483 & 1.218\\*
  Vicarious Sensation & Match & Intercept & 0.029 & 0.169\\*
  Vicarious Sensation & Residual & Observation & 0.770 & 0.878\\
  \addlinespace
  Overall Quality & Participant & Intercept & 0.977 & 0.988\\*
  Overall Quality & Match & Intercept & 0.000 & 0.000\\*
  Overall Quality & Residual & Observation & 1.415 & 1.189\\
  \addlinespace
  Word Choice & Participant & Intercept & 0.806 & 0.898\\*
  Word Choice & Match & Intercept & 0.000 & 0.000\\*
  Word Choice & Residual & Observation & 1.440 & 1.200\\
  \addlinespace
  Readability & Participant & Intercept & 1.038 & 1.019\\*
  Readability & Match & Intercept & 0.101 & 0.318\\*
  Readability & Residual & Observation & 1.534 & 1.239\\
  \addlinespace
  Logicality & Participant & Intercept & 0.400 & 0.633\\*
  Logicality & Match & Intercept & 0.008 & 0.088\\*
  Logicality & Residual & Observation & 0.966 & 0.983\\
  \addlinespace
  Detail Level & Participant & Intercept & 1.268 & 1.126\\*
  Detail Level & Match & Intercept & 0.142 & 0.377\\*
  Detail Level & Residual & Observation & 1.419 & 1.191\\
  \addlinespace
  Reading Thoroughness & Participant & Intercept & 0.813 & 0.901\\*
  Reading Thoroughness & Match & Intercept & 0.277 & 0.526\\*
  Reading Thoroughness & Residual & Observation & 1.788 & 1.337\\
  \addlinespace
  Understanding & Participant & Intercept & 1.563 & 1.250\\*
  Understanding & Match & Intercept & 0.052 & 0.227\\*
  Understanding & Residual & Observation & 1.324 & 1.151\\*
\end{longtable}

\section{Model Fit Metrics}

\begin{table}[H]
  \centering\small
  \begin{tabular}{lcc}
    \toprule
    \textbf{Outcome} & \textbf{Marginal $R^2$} & \textbf{Conditional $R^2$} \\
    \midrule
    Competitive Nature & 0.017 & 0.741 \\
    Skill Improvement & 0.019 & 0.795 \\
    Game Knowledge & 0.008 & 0.730 \\
    Skill Appreciation & 0.017 & 0.749 \\
    Entertaining Nature & 0.017 & 0.715 \\
    Dramatic Nature & 0.010 & 0.675 \\
    Vicarious Sensation & 0.034 & 0.674 \\
    Overall Quality & 0.025 & 0.423 \\
    Word Choice & 0.027 & 0.376 \\
    Readability & 0.008 & 0.431 \\
    Logicality & 0.062 & 0.341 \\
    Detail Level & 0.023 & 0.510 \\
    Reading Thoroughness & 0.017 & 0.389 \\
    Understanding & 0.023 & 0.560 \\
    \bottomrule
  \end{tabular}
  \caption{`mrginal and conditional $R^2$ for each linear mixed model (marginal: fixed effects only; conditional: fixed + random intercepts). Values are computed from the fitted LMMs reported in Table~\ref{app:fixed_effects}.}
  \label{tab:r2_nakagawa}
\end{table}

\begin{table}[H]
  \centering\small
  \begin{tabular}{lcc}
    \toprule
    \textbf{Outcome} & \textbf{ICC (Participant)} & \textbf{ICC (Match)} \\
    \midrule
    Competitive Nature & 0.728 & 0.009 \\
    Skill Improvement & 0.788 & 0.003 \\
    Game Knowledge & 0.717 & 0.010 \\
    Skill Appreciation & 0.723 & 0.022 \\
    Entertaining Nature & 0.699 & 0.011 \\
    Dramatic Nature & 0.651 & 0.020 \\
    Vicarious Sensation & 0.650 & 0.013 \\
    Overall Quality & 0.408 & 0.000 \\
    Word Choice & 0.359 & 0.000 \\
    Readability & 0.388 & 0.038 \\
    Logicality & 0.291 & 0.006 \\
    Detail Level & 0.448 & 0.050 \\
    Reading Thoroughness & 0.282 & 0.096 \\
    Understanding & 0.532 & 0.018 \\
    \bottomrule
  \end{tabular}
  \caption{Intraclass correlation coefficients (ICCs) derived from the linear mixed models. With crossed random intercepts, we report the proportion of total random-intercept-plus-residual variance attributable to participants versus matches (stimuli).}
  \label{tab:icc}
\end{table}

\endgroup

